\documentclass[12pt,a4paper]{report}
\usepackage[utf8]{inputenc}
\usepackage[T1]{fontenc}
\usepackage[french]{babel}
\usepackage{amsmath, amsfonts, amssymb}
\usepackage{geometry}
\geometry{a4paper, margin=1in}
\usepackage{tikz}
\usepackage{listings}
\usepackage{xcolor}

\lstset{
    basicstyle=\ttfamily\small,
    breaklines=true,
    keywordstyle=\color{blue},
    commentstyle=\color{gray},
    stringstyle=\color{red}
}

\title{Continuité des fonctions d'une variable réelle, théorème des valeurs intermédiaires et compacité locale}
\author{Charles EDOU NZE}
\date{}

\begin{document}
\maketitle
\tableofcontents
\newpage

\chapter{Intuition, Formalisation et Démonstrations}

\begin{center}
\begin{tikzpicture}[scale=1.5]
  % Axes
  \draw[->] (-0.5,0) -- (4,0) node[right] {$x$};
  \draw[->] (0,-1) -- (0,3) node[above] {$y$};

  % Points a, b and f(a), f(b)
  \draw[dashed] (1,0) node[below] {$a$} -- (1,-0.5) node[left] {$f(a)$};
  \draw[dashed] (3,0) node[below] {$b$} -- (3,2.5) node[left] {$f(b)$};
  \draw (0,-0.5) node[left] {$f(a)$} -- (0.1,-0.5);
  \draw (0,2.5) node[left] {$f(b)$} -- (0.1,2.5);

  % y_0
  \draw[dashed] (0,1.2) node[left] {$y_0$} -- (3.5,1.2);

  % Function curve (continuous)
  \draw[thick, blue] (1,-0.5) .. controls (1.5,0.5) and (2,0.8) .. (2.2,1.2) .. controls (2.4,1.6) and (2.8,2) .. (3,2.5);

  % Intersection c
  \filldraw[red] (2.2,1.2) circle (1.5pt);
  \draw[dashed, red] (2.2,1.2) -- (2.2,0) node[below] {$c$};

\end{tikzpicture}
\end{center}



\section{Jalon 18 : Continuité des fonctions d'une variable réelle, théorème des valeurs intermédiaires et compacité locale}

\subsection{1. Présentation du concept clé}
 Imaginez que vous dessiniez une courbe avec un crayon sur une feuille de papier. La \textbf{Continuité}, c'est simplement la règle d'or : vous n'avez pas le droit de lever le crayon. Si vous devez passer de la gauche à la droite de la feuille, votre trait doit être ininterrompu. Le \textbf{Théorème des Valeurs Intermédiaires (TVI)}, c'est comme dire que si vous commencez à dessiner en bas de la feuille et que vous finissez en haut, votre crayon a forcément dû traverser la ligne du milieu à un moment donné. Vous ne pouvez pas vous téléporter !
\begin{itemize}
\item \textbf{Le "Pourquoi on a inventé ça" :} Dans la nature, peu de choses changent instantanément. La température, la position d'une voiture, ou la croissance d'une plante sont des phénomènes continus. Les mathématiciens ont eu besoin de définir cette notion pour garantir que si on cherche une solution (un point où une fonction vaut zéro), elle existe vraiment. C'est fondamental pour modéliser des systèmes physiques et pour l'analyse numérique.
\end{itemize}
 Imaginez un élastique tendu. Si vous tirez sur un point, les points voisins suivent le mouvement. C'est la continuité : des causes proches produisent des effets proches. À l'inverse, une \textbf{discontinuité} serait une rupture nette de cet élastique, un "saut" ou un "trou" dans la courbe.

\subsection{2. Formalisation}
\textit{Le niveau bascule ici instantanément dans l'exigence pure des mathématiques supérieures.}

\subsubsection{A. Définitions Formelles}

Soit $I$ un intervalle de $\mathbb{R}$ et $f : I \to \mathbb{R}$ une fonction.

1.  \textbf{Continuité en un point $x_0 \in I$ :}
    La fonction $f$ est dite \textbf{continue en $x_0$} si et seulement si :
    $$\forall \epsilon > 0, \exists \delta > 0, \forall x \in I, (|x - x_0| < \delta \Rightarrow |f(x) - f(x_0)| < \epsilon)$$
    *   \textbf{Typage chirurgical :} Cette définition est la définition de Cauchy (ou $\epsilon-\delta$). Elle exprime que pour toute "tolérance" $\epsilon$ sur l'image, il existe une "tolérance" $\delta$ sur l'antécédent telle que tous les points $x$ suffisamment proches de $x_0$ ont leurs images $f(x)$ suffisamment proches de $f(x_0)$. Le $\delta$ dépend de $\epsilon$ et de $x_0$.
    *   \textbf{Définition séquentielle (Heine) :} La fonction $f$ est continue en $x_0$ si et seulement si pour toute suite $(x_n)_{n \in \mathbb{N}}$ d'éléments de $I$ qui converge vers $x_0$, la suite des images $(f(x_n))_{n \in \mathbb{N}}$ converge vers $f(x_0)$.
        $$\forall (x_n)_{n \in \mathbb{N}} \subset I, (\lim_{n \to \infty} x_n = x_0 \Rightarrow \lim_{n \to \infty} f(x_n) = f(x_0))$$
    *   \textbf{Exemple :} La fonction $f(x) = x\textasciicircum 2$ est continue en tout point $x_0 \in \mathbb{R}$. Pour un $\epsilon > 0$ donné, on cherche $\delta > 0$ tel que $|x - x_0| < \delta \Rightarrow |x^2 - x_0^2| < \epsilon$. On a $|x^2 - x_0^2| = |x - x_0||x + x_0|$. Si on suppose $\delta \le 1$, alors $|x - x_0| < 1 \Rightarrow x_0 - 1 < x < x_0 + 1 \Rightarrow |x + x_0| < |x_0 - 1 + x_0| + |x_0 + 1 + x_0| = |2x_0 - 1| + |2x_0 + 1|$. Plus simplement, si $|x-x_0|<1$, alors $|x|<|x_0|+1$, donc $|x+x_0| \le |x|+|x_0| < 2|x_0|+1$. Ainsi, $|x^2 - x_0^2| < \delta(2|x_0|+1)$. Pour que cela soit inférieur à $\epsilon$, on peut choisir $\delta = \min\left(1, \frac{\epsilon}{2|x_0|+1}\right)$.
    *   \textbf{Contre-exemple :} La fonction de Heaviside $H(x) = \begin{cases} 0 & \text{si } x < 0 \\ 1 & \text{si } x \ge 0 \end{cases}$ n'est pas continue en $x_0 = 0$.
        Pour $\epsilon = 1/2$, il n'existe aucun $\delta > 0$ tel que pour tout $x$ avec $|x - 0| < \delta$, on ait $|H(x) - H(0)| < 1/2$. En effet, si $x \in (-\delta, 0)$, alors $H(x) = 0$, donc $|H(x) - H(0)| = |0 - 1| = 1$, ce qui n'est pas inférieur à $1/2$.

2.  \textbf{Continuité à gauche et à droite en un point $x_0 \in I$ :}
    *   $f$ est \textbf{continue à droite en $x_0$} si :
        $$\forall \epsilon > 0, \exists \delta > 0, \forall x \in I, (x_0 \le x < x_0 + \delta \Rightarrow |f(x) - f(x_0)| < \epsilon)$$
    *   $f$ est \textbf{continue à gauche en $x_0$} si :
        $$\forall \epsilon > 0, \exists \delta > 0, \forall x \in I, (x_0 - \delta < x \le x_0 \Rightarrow |f(x) - f(x_0)| < \epsilon)$$
    *   \textbf{Proposition :} $f$ est continue en $x_0$ si et seulement si elle est continue à gauche et à droite en $x_0$.
    *   \textbf{Exemple :} La fonction $f(x) = \lfloor x \rfloor$ (partie entière) est continue à droite en tout $x_0 \in \mathbb{Z}$, mais pas continue à gauche.

3.  \textbf{Continuité sur un intervalle :}
    La fonction $f$ est dite \textbf{continue sur l'intervalle $I$} si elle est continue en tout point $x_0 \in I$. Si $I$ est un intervalle fermé $[a, b]$, la continuité en $a$ implique la continuité à droite en $a$, et la continuité en $b$ implique la continuité à gauche en $b$.

4.  \textbf{Types de discontinuités en un point $x_0$ :}
    *   \textbf{Discontinuité de première espèce (saut) :} Les limites à gauche et à droite existent et sont finies, mais différentes, ou l'une d'elles est différente de $f(x_0)$.
        $$\lim_{x \to x_0^-} f(x) \neq \lim_{x \to x_0^+} f(x)$$
        *   \textbf{Exemple :} La fonction $f(x) = \lfloor x \rfloor$ en $x_0 = 1$. $\lim_{x \to 1^-} f(x) = 0$ et $\lim_{x \to 1^+} f(x) = 1$.
    *   \textbf{Discontinuité de première espèce (éliminable ou "trou") :} Les limites à gauche et à droite existent et sont égales, mais différentes de $f(x_0)$, ou $f(x_0)$ n'est pas définie.
        *   \textbf{Exemple :} La fonction $f(x) = \frac{\sin(x)}{x}$ pour $x \neq 0$ et $f(0) = 0$. $\lim_{x \to 0} f(x) = 1 \neq f(0)$. On peut "recoller" la fonction en posant $f(0)=1$ pour la rendre continue.
    *   \textbf{Discontinuité de deuxième espèce (essentielle) :} Au moins une des limites à gauche ou à droite n'existe pas ou est infinie.
        *   \textbf{Exemple :} La fonction $f(x) = \sin(1/x)$ en $x_0 = 0$. La limite n'existe pas.

5.  \textbf{Continuité Uniforme :}
    La fonction $f$ est dite \textbf{uniformément continue sur $I$} si :
    $$\forall \epsilon > 0, \exists \delta > 0, \forall (x, y) \in I^2, (|x - y| < \delta \Rightarrow |f(x) - f(y)| < \epsilon)$$
    \textit{   \textbf{Typage chirurgical :} La différence cruciale avec la continuité simple est que le $\delta$ ne dépend }que* de $\epsilon$, et non du point $x_0$ (ou $x, y$) choisi dans l'intervalle $I$. Cela signifie que la "vitesse" à laquelle la fonction varie est bornée sur tout l'intervalle.
    *   \textbf{Exemple :} La fonction $f(x) = x$ est uniformément continue sur $\mathbb{R}$. Pour tout $\epsilon > 0$, on peut choisir $\delta = \epsilon$. Alors $|x-y| < \delta \Rightarrow |f(x)-f(y)| = |x-y| < \epsilon$.
    *   \textbf{Contre-exemple :} La fonction $f(x) = x\textasciicircum 2$ est continue sur $[0, +\infty[$ mais n'est pas uniformément continue sur cet intervalle.
        Pour le prouver, nous devons montrer qu'il existe un $\epsilon_0 > 0$ tel que pour tout $\delta > 0$, il existe $x, y \in [0, +\infty[$ avec $|x-y| < \delta$ et $|f(x)-f(y)| \ge \epsilon_0$.
        Prenons $\epsilon_0 = 1$. Soit $\delta > 0$ arbitraire.
        Choisissons $x = \frac{1}{\delta}$ et $y = \frac{1}{\delta} + \frac{\delta}{2}$.
        Alors $|x-y| = \left|\frac{1}{\delta} - \left(\frac{1}{\delta} + \frac{\delta}{2}\right)\right| = \frac{\delta}{2}$. Puisque $\delta > 0$, on a $\frac{\delta}{2} < \delta$.
        Calculons $|f(x)-f(y)| = |x^2 - y^2| = \left|\left(\frac{1}{\delta}\right)^2 - \left(\frac{1}{\delta} + \frac{\delta}{2}\right)^2\right|$.
        $|x^2 - y^2| = \left|\frac{1}{\delta^2} - \left(\frac{1}{\delta^2} + 1 + \frac{\delta^2}{4}\right)\right| = \left|-1 - \frac{\delta^2}{4}\right| = 1 + \frac{\delta^2}{4}$.
        Puisque $1 + \frac{\delta^2}{4} \ge 1 = \epsilon_0$, nous avons trouvé $x, y$ tels que $|x-y| < \delta$ mais $|f(x)-f(y)| \ge \epsilon_0$.
        Donc $f(x) = x\textasciicircum 2$ n'est pas uniformément continue sur $[0, +\infty[$.

6.  \textbf{Fonction Lipschitzienne :}
    La fonction $f$ est dite \textbf{Lipschitzienne sur $I$} s'il existe une constante $L \ge 0$ (appelée constante de Lipschitz) telle que :
    $$\forall (x, y) \in I^2, |f(x) - f(y)| \le L|x - y|$$
    *   \textbf{Typage chirurgical :} Une fonction Lipschitzienne est toujours uniformément continue. Pour un $\epsilon > 0$ donné, on choisit $\delta = \epsilon/L$ (si $L>0$). Alors $|x-y| < \delta \Rightarrow |f(x)-f(y)| \le L|x-y| < L(\epsilon/L) = \epsilon$. Si $L=0$, $f$ est constante, donc uniformément continue.
    *   \textbf{Exemple :} $f(x) = \sin(x)$ est Lipschitzienne sur $\mathbb{R}$ avec $L=1$, car $|\sin(x) - \sin(y)| \le |x-y|$ (d'après l'inégalité des accroissements finis, ou directement par la formule de soustraction des sinus).

\subsubsection{B. Théorèmes, Propositions \& Lemmes}

1.  \textbf{Propriétés algébriques des fonctions continues :}
    Soient $f, g : I \to \mathbb{R}$ deux fonctions continues en $x_0 \in I$.
    *   La somme $f+g$ est continue en $x_0$.
    *   Le produit $f \cdot g$ est continue en $x_0$.
    *   Pour tout réel $\lambda$, la fonction $\lambda f$ est continue en $x_0$.
    *   Si $g(x_0) \neq 0$, alors le quotient $f/g$ est continue en $x_0$.
    *   \textbf{Proposition :} Les fonctions polynomiales sont continues sur $\mathbb{R}$. Les fonctions rationnelles sont continues sur leur domaine de définition.

2.  \textbf{Composition de fonctions continues :}
    Soient $f : I \to J$ et $g : J \to \mathbb{R}$ deux fonctions. Si $f$ est continue en $x_0 \in I$ et $g$ est continue en $f(x_0) \in J$, alors la fonction composée $g \circ f : I \to \mathbb{R}$ est continue en $x_0$.

3.  \textbf{Théorème des Valeurs Intermédiaires (TVI) :}
    Soit $f$ une fonction continue sur un intervalle $[a, b]$. Alors pour tout réel $y$ compris entre $f(a)$ et $f(b)$ (c'est-à-dire $y \in [\min(f(a), f(b)), \max(f(a), f(b))]$), il existe au moins un réel $c \in [a, b]$ tel que $f(c) = y$.
    *   \textbf{Corollaire (Théorème de Bolzano) :} Si $f$ est continue sur $[a, b]$ et $f(a)f(b) < 0$ (i.e., $f(a)$ et $f(b)$ sont de signes opposés), alors il existe au moins un $c \in ]a, b[$ tel que $f(c) = 0$.

4.  \textbf{Théorème de Heine (ou Théorème de la continuité uniforme sur un segment) :}
    Toute fonction continue sur un segment $[a, b]$ (c'est-à-dire un intervalle fermé et borné) est uniformément continue sur ce segment.
    *   \textbf{Typage chirurgical :} Ce théorème est fondamental car il établit un lien fort entre la continuité simple et la continuité uniforme sur des ensembles compacts (les segments de $\mathbb{R}$). Il n'est pas vrai sur des intervalles ouverts ou non bornés.

5.  \textbf{Théorème des Bornes Atteintes (ou Théorème de Weierstrass) :}
    L'image d'un segment $[a, b]$ par une fonction continue $f$ est un segment $[m, M]$. Autrement dit, $f([a, b]) = [\min_{x \in [a,b]} f(x), \max_{x \in [a,b]} f(x)]$. La fonction $f$ est donc bornée sur $[a, b]$ et atteint ses bornes (il existe $x_m, x_M \in [a, b]$ tels que $f(x_m) = m$ et $f(x_M) = M$).
    *   \textbf{Typage chirurgical :} Ce théorème garantit l'existence d'un minimum et d'un maximum globaux pour une fonction continue sur un segment. C'est une propriété cruciale pour l'optimisation.

6.  \textbf{Théorème de la fonction inverse continue :}
    Soit $f : I \to \mathbb{R}$ une fonction continue et strictement monotone sur un intervalle $I$. Alors $f$ est une bijection de $I$ sur $J = f(I)$, et sa fonction réciproque $f^{-1} : J \to I$ est continue et strictement monotone sur $J$.
    *   \textbf{Typage chirurgical :} L'image $J$ est également un intervalle.

\subsection{3. Démonstrations}
\subsubsection{Démonstration du Théorème Pivot : Théorème des Valeurs Intermédiaires (par dichotomie)}
Nous allons démontrer le cas particulier du TVI où $f(a) < 0 < f(b)$, et montrer qu'il existe $c \in [a, b]$ tel que $f(c) = 0$. Le cas général du TVI s'en déduit facilement.

\textbf{Énoncé :} Soit $f : [a, b] \to \mathbb{R}$ une fonction continue sur le segment $[a, b]$ telle que $f(a) < 0$ et $f(b) > 0$. Alors il existe au moins un réel $c \in ]a, b[$ tel que $f(c) = 0$.

\textbf{Démonstration :}

1.  \textbf{Initialisation / Cadre :}
    *   Nous définissons deux suites $(a_n)_{n \in \mathbb{N}}$ et $(b_n)_{n \in \mathbb{N}}$ par récurrence, qui vont encadrer la racine $c$.
    *   Posons $a_0 = a$ et $b_0 = b$.
    *   Par hypothèse, nous avons $f(a_0) < 0$ et $f(b_0) > 0$.

2.  \textbf{Étape 1 : Construction des suites par dichotomie}
    Supposons que pour un certain $n \in \mathbb{N}$, les termes $a_n$ et $b_n$ sont définis et vérifient les propriétés suivantes :
    (i) $a_n \in [a, b]$ et $b_n \in [a, b]$
    (ii) $a_n \le b_n$
    (iii) $f(a_n) < 0$
    (iv) $f(b_n) > 0$

    Nous construisons $a_{n+1}$ et $b_{n+1}$ comme suit :
    *   Calculons le milieu du segment $[a_n, b_n]$ : $m_n = \frac{a_n + b_n}{2}$.
    *   \textbf{Cas 1 :} Si $f(m_n) = 0$.
        Alors $c = m_n$ est une racine de $f$. La démonstration est terminée.
    *   \textbf{Cas 2 :} Si $f(m_n) < 0$.
        Dans ce cas, la racine doit se trouver dans l'intervalle $[m_n, b_n]$.
        Nous posons $a_{n+1} = m_n$ et $b_{n+1} = b_n$.
        Vérifions les propriétés pour $n+1$:
        (i) $a_{n+1} = m_n \in [a_n, b_n] \subset [a, b]$ et $b_{n+1} = b_n \in [a, b]$.
        (ii) $a_{n+1} = m_n = \frac{a_n+b_n}{2} \le b_n = b_{n+1}$.
        (iii) $f(a_{n+1}) = f(m_n) < 0$ par hypothèse de ce cas.
        (iv) $f(b_{n+1}) = f(b_n) > 0$ par hypothèse de récurrence.
    *   \textbf{Cas 3 :} Si $f(m_n) > 0$.
        Dans ce cas, la racine doit se trouver dans l'intervalle $[a_n, m_n]$.
        Nous posons $a_{n+1} = a_n$ et $b_{n+1} = m_n$.
        Vérifions les propriétés pour $n+1$:
        (i) $a_{n+1} = a_n \in [a, b]$ et $b_{n+1} = m_n \in [a_n, b_n] \subset [a, b]$.
        (ii) $a_{n+1} = a_n \le \frac{a_n+b_n}{2} = m_n = b_{n+1}$.
        (iii) $f(a_{n+1}) = f(a_n) < 0$ par hypothèse de récurrence.
        (iv) $f(b_{n+1}) = f(m_n) > 0$ par hypothèse de ce cas.

    Dans tous les cas (sauf si $f(m_n)=0$ et la preuve est finie), nous avons construit des suites $(a_n)$ et $(b_n)$ telles que $a_n \le a_{n+1}$, $b_{n+1} \le b_n$, $f(a_n) < 0$, $f(b_n) > 0$ et $b_{n+1} - a_{n+1} = \frac{b_n - a_n}{2}$.

3.  \textbf{Étape 2 : Convergence des suites}
    *   La suite $(a_n)_{n \in \mathbb{N}}$ est croissante par construction ($a_{n+1} \ge a_n$).
    *   La suite $(b_n)_{n \in \mathbb{N}}$ est décroissante par construction ($b_{n+1} \le b_n$).
    *   De plus, pour tout $n$, $a_n \le b_n$.
    *   Les suites sont bornées : $a \le a_n \le b_n \le b$.
    *   Par le théorème de convergence monotone, $(a_n)$ converge vers une limite $c_1 \in [a, b]$ et $(b_n)$ converge vers une limite $c_2 \in [a, b]$.
    *   Calculons la différence entre les termes des suites :
        $b_n - a_n = \frac{b_{n-1} - a_{n-1}}{2} = \frac{b_{n-2} - a_{n-2}}{2^2} = \dots = \frac{b_0 - a_0}{2^n} = \frac{b - a}{2^n}$.
    *   Puisque $\lim_{n \to \infty} \frac{b - a}{2^n} = 0$, nous avons $\lim_{n \to \infty} (b_n - a_n) = 0$.
    *   Par conséquent, $c_1 = c_2$. Nous appelons cette limite commune $c$.
    *   Puisque $a_n \le c \le b_n$ pour tout $n$, et $a_0 = a$, $b_0 = b$, on a $c \in [a, b]$.
    *   De plus, comme $f(a) < 0$ et $f(b) > 0$, $c$ ne peut être ni $a$ ni $b$ (sauf si $f(a)=0$ ou $f(b)=0$, ce qui est exclu par l'hypothèse stricte $f(a)<0<f(b)$). Donc $c \in ]a, b[$.

4.  \textbf{Étape 3 : Utilisation de la continuité}
    *   La fonction $f$ est continue sur $[a, b]$, et donc en particulier en $c$.
    *   Par la définition séquentielle de la continuité (Heine), puisque $\lim_{n \to \infty} a_n = c$, nous avons $\lim_{n \to \infty} f(a_n) = f(c)$.
    *   De même, puisque $\lim_{n \to \infty} b_n = c$, nous avons $\lim_{n \to \infty} f(b_n) = f(c)$.
    *   Par construction, pour tout $n \in \mathbb{N}$, nous avons $f(a_n) < 0$.
    *   Par passage à la limite dans les inégalités larges (si une suite $(u_n)$ converge vers $L$ et $u_n \le K$ pour tout $n$, alors $L \le K$), nous obtenons $f(c) \le 0$.
    *   De même, pour tout $n \in \mathbb{N}$, nous avons $f(b_n) > 0$.
    *   Par passage à la limite dans les inégalités larges, nous obtenons $f(c) \ge 0$.

5.  \textbf{Conclusion :}
    *   Nous avons établi que $f(c) \le 0$ et $f(c) \ge 0$.
    *   La seule valeur possible pour $f(c)$ qui satisfait ces deux inégalités est $f(c) = 0$.
    *   Le point $c \in ]a, b[$ tel que $f(c) = 0$ existe bien. Le théorème est démontré pour le cas $f(a) < 0 < f(b)$.

\textbf{Démonstration du cas général du TVI :}
Soit $f$ continue sur $[a, b]$. Soit $y$ un réel compris entre $f(a)$ et $f(b)$.
Sans perte de généralité, supposons $f(a) \le y \le f(b)$.
Considérons la fonction auxiliaire $g(x) = f(x) - y$.
*   $g$ est continue sur $[a, b]$ car $f$ est continue et $x \mapsto y$ est continue.
*   $g(a) = f(a) - y \le 0$ (par hypothèse $f(a) \le y$).
*   $g(b) = f(b) - y \ge 0$ (par hypothèse $y \le f(b)$).
Si $g(a)=0$ ou $g(b)=0$, alors $a$ ou $b$ est le $c$ recherché.
Sinon, $g(a) < 0$ et $g(b) > 0$. D'après la démonstration précédente (cas particulier du TVI), il existe $c \in ]a, b[$ tel que $g(c) = 0$.
Ceci implique $f(c) - y = 0$, soit $f(c) = y$.
Le théorème des valeurs intermédiaires est ainsi démontré dans sa généralité.

\subsubsection{Démonstration des propriétés algébriques de la continuité (pour la somme)}
\textbf{Énoncé :} Soient $f, g : I \to \mathbb{R}$ deux fonctions continues en $x_0 \in I$. Alors la fonction $h = f+g$ est continue en $x_0$.

\textbf{Démonstration :}
Nous voulons montrer que pour tout $\epsilon > 0$, il existe $\delta > 0$ tel que pour tout $x \in I$, si $|x - x_0| < \delta$, alors $|h(x) - h(x_0)| < \epsilon$.

1.  Soit $\epsilon > 0$ arbitrairement choisi.
2.  Puisque $f$ est continue en $x_0$, par définition, il existe $\delta_1 > 0$ tel que pour tout $x \in I$, si $|x - x_0| < \delta_1$, alors $|f(x) - f(x_0)| < \frac{\epsilon}{2}$.
3.  Puisque $g$ est continue en $x_0$, par définition, il existe $\delta_2 > 0$ tel que pour tout $x \in I$, si $|x - x_0| < \delta_2$, alors $|g(x) - g(x_0)| < \frac{\epsilon}{2}$.
4.  Choisissons $\delta = \min(\delta_1, \delta_2)$. Puisque $\delta_1 > 0$ et $\delta_2 > 0$, $\delta$ est bien un nombre réel strictement positif.
5.  Maintenant, considérons un $x \in I$ tel que $|x - x_0| < \delta$.
    *   Puisque $\delta \le \delta_1$, nous avons $|x - x_0| < \delta_1$, ce qui implique $|f(x) - f(x_0)| < \frac{\epsilon}{2}$.
    *   Puisque $\delta \le \delta_2$, nous avons $|x - x_0| < \delta_2$, ce qui implique $|g(x) - g(x_0)| < \frac{\epsilon}{2}$.
6.  Nous voulons majorer $|h(x) - h(x_0)|$:
    $|h(x) - h(x_0)| = |(f(x) + g(x)) - (f(x_0) + g(x_0))|$
    $|h(x) - h(x_0)| = |(f(x) - f(x_0)) + (g(x) - g(x_0))|$
7.  Par l'inégalité triangulaire, nous avons :
    $|(f(x) - f(x_0)) + (g(x) - g(x_0))| \le |f(x) - f(x_0)| + |g(x) - g(x_0)|$
8.  En utilisant les majorations obtenues aux étapes 5 :
    $|f(x) - f(x_0)| + |g(x) - g(x_0)| < \frac{\epsilon}{2} + \frac{\epsilon}{2} = \epsilon$.
9.  Donc, pour tout $x \in I$ tel que $|x - x_0| < \delta$, nous avons $|h(x) - h(x_0)| < \epsilon$.
10. Puisque $\epsilon$ était arbitraire, $h = f+g$ est continue en $x_0$.

\subsubsection{Démonstration du Théorème de Weierstrass (Théorème des Bornes Atteintes)}
\textbf{Énoncé :} Soit $f : [a, b] \to \mathbb{R}$ une fonction continue sur le segment $[a, b]$. Alors $f$ est bornée sur $[a, b]$ et atteint ses bornes. C'est-à-dire, il existe $x_m, x_M \in [a, b]$ tels que $f(x_m) = \min_{x \in [a,b]} f(x)$ et $f(x_M) = \max_{x \in [a,b]} f(x)$.

\textbf{Démonstration :}

\textbf{Partie 1 : $f$ est bornée sur $[a, b]$.}
Nous allons démontrer que $f$ est bornée supérieurement. La démonstration pour la borne inférieure est similaire.
Supposons par l'absurde que $f$ n'est pas bornée supérieurement sur $[a, b]$.
1.  Alors, pour tout $n \in \mathbb{N}^*$, il existe un $x_n \in [a, b]$ tel que $f(x_n) > n$.
2.  La suite $(x_n)_{n \in \mathbb{N}^*}$ est une suite d'éléments du segment $[a, b]$.
3.  Puisque $[a, b]$ est un segment (fermé et borné), il est compact. D'après le théorème de Bolzano-Weierstrass, toute suite bornée admet une sous-suite convergente.
4.  Il existe donc une sous-suite $(x_{\phi(n)})_{n \in \mathbb{N}^*}$ qui converge vers un point $c \in [a, b]$ (car $[a, b]$ est fermé).
5.  Puisque $f$ est continue en $c$, par la définition séquentielle de la continuité, nous avons $\lim_{n \to \infty} f(x_{\phi(n)}) = f(c)$.
6.  Cependant, par construction de la suite $(x_n)$, nous avons $f(x_{\phi(n)}) > \phi(n)$.
7.  Comme $\phi(n)$ est une suite strictement croissante d'entiers, $\lim_{n \to \infty} \phi(n) = +\infty$.
8.  Donc, $\lim_{n \to \infty} f(x_{\phi(n)}) = +\infty$.
9.  Ceci contredit le fait que $\lim_{n \to \infty} f(x_{\phi(n)}) = f(c)$, car $f(c)$ est un nombre réel fini.
10. Par conséquent, l'hypothèse de départ est fausse : $f$ doit être bornée supérieurement sur $[a, b]$.
11. De manière analogue, en considérant $f(x_n) < -n$, on montre que $f$ est bornée inférieurement.
12. Donc, $f$ est bornée sur $[a, b]$.

\textbf{Partie 2 : $f$ atteint ses bornes.}
Puisque $f$ est bornée sur $[a, b]$, l'ensemble image $f([a, b])$ est borné.
1.  Soit $M = \sup_{x \in [a,b]} f(x)$. Puisque $f([a, b])$ est non vide et borné supérieurement, $M$ existe et est fini.
2.  Nous voulons montrer qu'il existe $x_M \in [a, b]$ tel que $f(x_M) = M$.
3.  Par la définition de la borne supérieure, pour tout $n \in \mathbb{N}^*$, il existe un $x_n \in [a, b]$ tel que $M - \frac{1}{n} < f(x_n) \le M$.
4.  La suite $(x_n)_{n \in \mathbb{N}^*}$ est une suite d'éléments du segment $[a, b]$.
5.  D'après le théorème de Bolzano-Weierstrass, il existe une sous-suite $(x_{\phi(n)})_{n \in \mathbb{N}^*}$ qui converge vers un point $x_M \in [a, b]$.
6.  Puisque $f$ est continue en $x_M$, par la définition séquentielle de la continuité, nous avons $\lim_{n \to \infty} f(x_{\phi(n)}) = f(x_M)$.
7.  D'autre part, pour la sous-suite, nous avons $M - \frac{1}{\phi(n)} < f(x_{\phi(n)}) \le M$.
8.  Par le théorème des gendarmes, puisque $\lim_{n \to \infty} (M - \frac{1}{\phi(n)}) = M$ et $\lim_{n \to \infty} M = M$, nous avons $\lim_{n \to \infty} f(x_{\phi(n)}) = M$.
9.  Par unicité de la limite, $f(x_M) = M$.
10. Donc, $f$ atteint sa borne supérieure sur $[a, b]$.
11. De manière analogue, en considérant la borne inférieure $m = \inf_{x \in [a,b]} f(x)$, on montre qu'il existe $x_m \in [a, b]$ tel que $f(x_m) = m$.
12. Le théorème de Weierstrass est démontré.

\subsubsection{Démonstration du Théorème de Heine (Théorème de la continuité uniforme sur un segment)}
\textbf{Énoncé :} Toute fonction continue sur un segment $[a, b]$ est uniformément continue sur ce segment.

\textbf{Démonstration :}
Supposons par l'absurde que $f$ est continue sur $[a, b]$ mais n'est pas uniformément continue sur $[a, b]$.
1.  Si $f$ n'est pas uniformément continue, alors la négation de la définition de la continuité uniforme est vraie :
    Il existe un $\epsilon_0 > 0$ tel que pour tout $\delta > 0$, il existe $x, y \in [a, b]$ tels que $|x - y| < \delta$ et $|f(x) - f(y)| \ge \epsilon_0$.
2.  En particulier, pour chaque $n \in \mathbb{N}^*$, en prenant $\delta = \frac{1}{n}$, il existe une paire de points $(x_n, y_n)$ dans $[a, b]^2$ telle que :
    (i) $|x_n - y_n| < \frac{1}{n}$
    (ii) $|f(x_n) - f(y_n)| \ge \epsilon_0$
3.  Les suites $(x_n)_{n \in \mathbb{N}^*}$ et $(y_n)_{n \in \mathbb{N}^*}$ sont des suites d'éléments du segment $[a, b]$.
4.  Puisque $[a, b]$ est un segment (fermé et borné), il est compact. D'après le théorème de Bolzano-Weierstrass, la suite $(x_n)$ admet une sous-suite convergente $(x_{\phi(n)})_{n \in \mathbb{N}^*}$.
5.  Soit $c = \lim_{n \to \infty} x_{\phi(n)}$. Puisque $[a, b]$ est fermé, $c \in [a, b]$.
6.  Considérons maintenant la sous-suite correspondante $(y_{\phi(n)})_{n \in \mathbb{N}^*}$.
    Nous avons $|x_{\phi(n)} - y_{\phi(n)}| < \frac{1}{\phi(n)}$.
    Puisque $\phi(n) \to \infty$ quand $n \to \infty$, nous avons $\frac{1}{\phi(n)} \to 0$ quand $n \to \infty$.
    Par le théorème des gendarmes, $\lim_{n \to \infty} |x_{\phi(n)} - y_{\phi(n)}| = 0$.
7.  Puisque $\lim_{n \to \infty} x_{\phi(n)} = c$ et $\lim_{n \to \infty} (x_{\phi(n)} - y_{\phi(n)}) = 0$, nous pouvons déduire la limite de $y_{\phi(n)}$:
    $\lim_{n \to \infty} y_{\phi(n)} = \lim_{n \to \infty} (x_{\phi(n)} - (x_{\phi(n)} - y_{\phi(n)})) = \lim_{n \to \infty} x_{\phi(n)} - \lim_{n \to \infty} (x_{\phi(n)} - y_{\phi(n)}) = c - 0 = c$.
    Donc, la sous-suite $(y_{\phi(n)})$ converge également vers $c$.
8.  Puisque $f$ est continue sur $[a, b]$, elle est continue en $c$.
    Par la définition séquentielle de la continuité :
    *   $\lim_{n \to \infty} f(x_{\phi(n)}) = f(c)$.
    *   $\lim_{n \to \infty} f(y_{\phi(n)}) = f(c)$.
9.  Par conséquent, $\lim_{n \to \infty} (f(x_{\phi(n)}) - f(y_{\phi(n)})) = f(c) - f(c) = 0$.
10. Cependant, par construction de la sous-suite, nous avons $|f(x_{\phi(n)}) - f(y_{\phi(n)})| \ge \epsilon_0$ pour tout $n$.
11. Par passage à la limite dans les inégalités larges, nous obtenons $|0| \ge \epsilon_0$, ce qui signifie $0 \ge \epsilon_0$.
12. Ceci contredit notre hypothèse initiale que $\epsilon_0 > 0$.
13. Par conséquent, l'hypothèse de départ est fausse : $f$ doit être uniformément continue sur $[a, b]$.
14. Le théorème de Heine est démontré.

\subsection{4. Exercices d'Application}
\textit{Proposer au moins 2 exercices progressifs corrigés de façon exhaustive, sans aucune ellipse.}

\subsubsection{Exercice 1 : Point Fixe (Application du TVI)}
\textbf{Énoncé :} Soit $f : [0, 1] \to [0, 1]$ une fonction continue. Démontrer qu'il existe au moins un réel $c \in [0, 1]$ tel que $f(c) = c$. Un tel $c$ est appelé un point fixe de $f$.

\textbf{Correction Détaillée :}

1.  \textbf{Définition de la fonction auxiliaire :}
    Considérons la fonction auxiliaire $g(x)$ définie sur l'intervalle $[0, 1]$ par $g(x) = f(x) - x$.
2.  \textbf{Continuité de la fonction auxiliaire :}
    *   La fonction $f$ est continue sur $[0, 1]$ par hypothèse.
    *   La fonction $h(x) = x$ (fonction identité) est continue sur $[0, 1]$ (c'est une fonction polynomiale).
    *   Puisque la différence de deux fonctions continues est continue, la fonction $g(x) = f(x) - x$ est continue sur l'intervalle $[0, 1]$.
3.  \textbf{Évaluation aux bornes de l'intervalle :}
    *   Calculons la valeur de $g$ en $x=0$:
        $g(0) = f(0) - 0 = f(0)$.
        Puisque $f : [0, 1] \to [0, 1]$, l'image $f(0)$ doit appartenir à l'intervalle $[0, 1]$.
        Par conséquent, $f(0) \ge 0$. Donc, $g(0) \ge 0$.
    *   Calculons la valeur de $g$ en $x=1$:
        $g(1) = f(1) - 1$.
        Puisque $f : [0, 1] \to [0, 1]$, l'image $f(1)$ doit appartenir à l'intervalle $[0, 1]$.
        Par conséquent, $f(1) \le 1$. Donc, $f(1) - 1 \le 0$. Donc, $g(1) \le 0$.
4.  \textbf{Application du Théorème des Valeurs Intermédiaires (TVI) :}
    Nous avons deux cas possibles :
    *   \textbf{Cas 1 :} Si $g(0) = 0$.
        Alors $f(0) - 0 = 0$, ce qui signifie $f(0) = 0$. Dans ce cas, $c=0$ est un point fixe.
    *   \textbf{Cas 2 :} Si $g(1) = 0$.
        Alors $f(1) - 1 = 0$, ce qui signifie $f(1) = 1$. Dans ce cas, $c=1$ est un point fixe.
    *   \textbf{Cas 3 :} Si $g(0) > 0$ et $g(1) < 0$.
        Dans ce cas, nous avons $g(1) < 0 < g(0)$.
        La fonction $g$ est continue sur le segment $[0, 1]$.
        D'après le Théorème des Valeurs Intermédiaires (TVI), pour toute valeur $y$ comprise entre $g(1)$ et $g(0)$, il existe un $c \in [0, 1]$ tel que $g(c) = y$.
        En particulier, puisque $0$ est compris entre $g(1)$ et $g(0)$, il existe au moins un réel $c \in ]0, 1[$ tel que $g(c) = 0$.
5.  \textbf{Conclusion :}
    Dans tous les cas, il existe un réel $c \in [0, 1]$ tel que $g(c) = 0$.
    Par définition de $g(x)$, cela signifie $f(c) - c = 0$, ce qui équivaut à $f(c) = c$.
    Ainsi, toute fonction continue d'un segment dans lui-même admet au moins un point fixe.

\subsubsection{Exercice 2 : Continuité Uniforme de $\sqrt{x}$}
\textbf{Énoncé :} Démontrer que la fonction $f(x) = \sqrt{x}$ est uniformément continue sur $[0, +\infty[$.

\textbf{Correction Détaillée :}
Pour démontrer l'uniforme continuité sur $[0, +\infty[$, nous allons diviser l'intervalle en deux parties : un segment $[0, M]$ et un intervalle non borné $[M, +\infty[$, puis recoller les résultats.

1.  \textbf{Uniforme continuité sur un segment $[0, 1]$ :}
    *   La fonction $f(x) = \sqrt{x}$ est continue sur l'intervalle $[0, 1]$.
    *   L'intervalle $[0, 1]$ est un segment (fermé et borné).
    *   D'après le Théorème de Heine, toute fonction continue sur un segment est uniformément continue sur ce segment.
    *   Donc, $f(x) = \sqrt{x}$ est uniformément continue sur $[0, 1]$.
    *   Cela signifie que pour tout $\epsilon > 0$, il existe $\delta_1 > 0$ tel que pour tout $(x, y) \in [0, 1]^2$, si $|x - y| < \delta_1$, alors $|f(x) - f(y)| < \epsilon$.

2.  \textbf{Uniforme continuité sur un intervalle non borné $[1, +\infty[$ :}
    *   Considérons $x, y \in [1, +\infty[$.
    *   Nous voulons majorer $|f(x) - f(y)| = |\sqrt{x} - \sqrt{y}|$.
    *   Pour éviter les racines au dénominateur, nous multiplions par l'expression conjuguée :
        $|\sqrt{x} - \sqrt{y}| = \left|\frac{(\sqrt{x} - \sqrt{y})(\sqrt{x} + \sqrt{y})}{\sqrt{x} + \sqrt{y}}\right| = \left|\frac{x - y}{\sqrt{x} + \sqrt{y}}\right| = \frac{|x - y|}{\sqrt{x} + \sqrt{y}}$.
    *   Puisque $x \ge 1$ et $y \ge 1$, nous avons $\sqrt{x} \ge 1$ et $\sqrt{y} \ge 1$.
    *   Par conséquent, $\sqrt{x} + \sqrt{y} \ge 1 + 1 = 2$.
    *   Donc, $\frac{1}{\sqrt{x} + \sqrt{y}} \le \frac{1}{2}$.
    *   En substituant cette inégalité, nous obtenons :
        $|\sqrt{x} - \sqrt{y}| \le \frac{1}{2} |x - y|$.
    *   Cette inégalité montre que $f(x) = \sqrt{x}$ est une fonction \textbf{Lipschitzienne} sur $[1, +\infty[$ avec une constante de Lipschitz $L = 1/2$.
    *   Toute fonction Lipschitzienne est uniformément continue.
    *   Pour un $\epsilon > 0$ donné, nous pouvons choisir $\delta_2 = 2\epsilon$.
    *   Alors, si $|x - y| < \delta_2 = 2\epsilon$, nous avons $|\sqrt{x} - \sqrt{y}| \le \frac{1}{2} |x - y| < \frac{1}{2} (2\epsilon) = \epsilon$.
    *   Donc, $f(x) = \sqrt{x}$ est uniformément continue sur $[1, +\infty[$.

3.  \textbf{Recollement des continuités uniformes sur $[0, +\infty[$ :}
    Nous avons montré que $f$ est uniformément continue sur $[0, 1]$ (avec $\delta_1$) et sur $[1, +\infty[$ (avec $\delta_2$). Nous voulons montrer qu'elle l'est sur l'union $[0, +\infty[$.
    Soit $\epsilon > 0$ donné.
    *   D'après le point 1, il existe $\delta_1 > 0$ tel que pour tout $(x, y) \in [0, 1]^2$, si $|x - y| < \delta_1$, alors $|f(x) - f(y)| < \epsilon$.
    *   D'après le point 2, il existe $\delta_2 > 0$ tel que pour tout $(x, y) \in [1, +\infty[^2$, si $|x - y| < \delta_2$, alors $|f(x) - f(y)| < \epsilon$.
    *   Choisissons $\delta = \min(\delta_1, \delta_2, 1)$. (Le choix de 1 est pour s'assurer que si $x$ et $y$ sont "proches" et de part et d'autre de 1, ils ne sont pas trop éloignés de 1).
    *   Soient $x, y \in [0, +\infty[$ tels que $|x - y| < \delta$.
    *   \textbf{Cas A :} $x, y \in [0, 1]$.
        Puisque $|x - y| < \delta \le \delta_1$, nous avons $|f(x) - f(y)| < \epsilon$.
    *   \textbf{Cas B :} $x, y \in [1, +\infty[$.
        Puisque $|x - y| < \delta \le \delta_2$, nous avons $|f(x) - f(y)| < \epsilon$.
    *   \textbf{Cas C :} $x \in [0, 1]$ et $y \in [1, +\infty[$ (ou inversement, par symétrie).
        Puisque $|x - y| < \delta \le 1$, nous avons $y - x < 1$.
        Comme $x \le 1$ et $y \ge 1$, on a $y-x \ge 0$.
        Considérons le point $1$. Nous pouvons écrire :
        $|f(x) - f(y)| = |f(x) - f(1) + f(1) - f(y)| \le |f(x) - f(1)| + |f(1) - f(y)|$ (par inégalité triangulaire).
        *   Pour le terme $|f(x) - f(1)|$: $x \in [0, 1]$ et $1 \in [0, 1]$. De plus, $|x - 1| \le |x - y| < \delta \le \delta_1$.
            Donc, $|f(x) - f(1)| < \epsilon$.
        *   Pour le terme $|f(1) - f(y)|$: $1 \in [1, +\infty[$ et $y \in [1, +\infty[$. De plus, $|1 - y| \le |x - y| < \delta \le \delta_2$.
            Donc, $|f(1) - f(y)| < \epsilon$.
        *   En combinant, $|f(x) - f(y)| < \epsilon + \epsilon = 2\epsilon$.
        *   Pour que cela soit inférieur à $\epsilon$, nous devons ajuster notre choix initial de $\delta_1$ et $\delta_2$. Si nous avions choisi $\delta_1'$ et $\delta_2'$ pour $\epsilon/2$, alors le recollement donnerait $\epsilon/2 + \epsilon/2 = \epsilon$.
        *   Reprenons : Pour $\epsilon > 0$, il existe $\delta_1'$ tel que pour $(x,y) \in [0,1]^2$, $|x-y|<\delta_1' \Rightarrow |f(x)-f(y)| < \epsilon/2$.
        *   Il existe $\delta_2'$ tel que pour $(x,y) \in [1,+\infty[^2$, $|x-y|<\delta_2' \Rightarrow |f(x)-f(y)| < \epsilon/2$.
        *   Choisissons $\delta = \min(\delta_1', \delta_2', 1)$.
        *   Si $x \in [0,1]$ et $y \in [1,+\infty[$ avec $|x-y|<\delta$:
            $|f(x)-f(y)| \le |f(x)-f(1)| + |f(1)-f(y)|$.
            Comme $|x-1| \le |x-y| < \delta \le \delta_1'$, alors $|f(x)-f(1)| < \epsilon/2$.
            Comme $|1-y| \le |x-y| < \delta \le \delta_2'$, alors $|f(1)-f(y)| < \epsilon/2$.
            Donc, $|f(x)-f(y)| < \epsilon/2 + \epsilon/2 = \epsilon$.

    Dans tous les cas, pour tout $\epsilon > 0$, nous avons trouvé un $\delta > 0$ tel que si $|x - y| < \delta$, alors $|f(x) - f(y)| < \epsilon$.
\textbf{Conclusion :} La fonction $f(x) = \sqrt{x}$ est bien uniformément continue sur $[0, +\infty[$.

\subsubsection{Exercice 3 : Fonction continue mais non uniformément continue}
\textbf{Énoncé :} Démontrer que la fonction $f(x) = \frac{1}{x}$ est continue sur $]0, 1]$ mais n'est pas uniformément continue sur cet intervalle.

\textbf{Correction Détaillée :}

\textbf{Partie 1 : Continuité sur $]0, 1]$}
1.  Soit $x_0 \in ]0, 1]$. Nous voulons montrer que $f$ est continue en $x_0$.
2.  Soit $\epsilon > 0$ arbitrairement choisi.
3.  Nous cherchons $\delta > 0$ tel que pour tout $x \in ]0, 1]$, si $|x - x_0| < \delta$, alors $|f(x) - f(x_0)| < \epsilon$.
4.  $|f(x) - f(x_0)| = \left|\frac{1}{x} - \frac{1}{x_0}\right| = \left|\frac{x_0 - x}{x x_0}\right| = \frac{|x - x_0|}{|x x_0|}$.
5.  Pour majorer cette expression, nous devons minorer $|x x_0|$.
    Puisque $x_0 \in ]0, 1]$, $x_0 > 0$.
    Choisissons $\delta$ tel que $\delta \le \frac{x_0}{2}$.
    Alors, si $|x - x_0| < \delta$, nous avons $x_0 - \delta < x < x_0 + \delta$.
    Puisque $\delta \le \frac{x_0}{2}$, on a $x_0 - \delta \ge x_0 - \frac{x_0}{2} = \frac{x_0}{2}$.
    Donc, $x > \frac{x_0}{2}$.
    Par conséquent, $x x_0 > \frac{x_0}{2} \cdot x_0 = \frac{x_0^2}{2}$.
    Donc, $\frac{1}{|x x_0|} < \frac{2}{x_0^2}$.
6.  En substituant dans l'expression de $|f(x) - f(x_0)|$:
    $|f(x) - f(x_0)| < |x - x_0| \cdot \frac{2}{x_0^2}$.
7.  Nous voulons que cette expression soit inférieure à $\epsilon$. Donc, nous voulons $|x - x_0| \cdot \frac{2}{x_0^2} < \epsilon$.
    Cela implique $|x - x_0| < \frac{\epsilon x_0^2}{2}$.
8.  Nous choisissons $\delta = \min\left(\frac{x_0}{2}, \frac{\epsilon x_0^2}{2}\right)$.
9.  Avec ce choix de $\delta$, pour tout $x \in ]0, 1]$ tel que $|x - x_0| < \delta$, nous avons :
    $|f(x) - f(x_0)| < \frac{\epsilon x_0^2}{2} \cdot \frac{2}{x_0^2} = \epsilon$.
10. Donc, $f(x) = \frac{1}{x}$ est continue en tout point $x_0 \in ]0, 1]$. Par conséquent, elle est continue sur $]0, 1]$.

\textbf{Partie 2 : Non-uniforme continuité sur $]0, 1]$}
Nous devons montrer qu'il existe un $\epsilon_0 > 0$ tel que pour tout $\delta > 0$, il existe $x, y \in ]0, 1]$ avec $|x - y| < \delta$ et $|f(x) - f(y)| \ge \epsilon_0$.

1.  Choisissons $\epsilon_0 = 1$.
2.  Soit $\delta > 0$ arbitrairement choisi.
3.  Nous devons trouver $x, y \in ]0, 1]$ tels que $|x - y| < \delta$ et $\left|\frac{1}{x} - \frac{1}{y}\right| \ge 1$.
4.  Considérons des points $x$ et $y$ très proches de $0$.
    Soit $y = \min\left(1, \frac{\delta}{2}\right)$. (Pour s'assurer que $y \in ]0, 1]$ et que $y$ est petit).
    Soit $x = \frac{y}{2}$.
    Alors $x \in ]0, 1]$ et $y \in ]0, 1]$.
5.  Calculons la distance entre $x$ et $y$:
    $|x - y| = \left|\frac{y}{2} - y\right| = \left|-\frac{y}{2}\right| = \frac{y}{2}$.
    Puisque $y \le \frac{\delta}{2}$, on a $\frac{y}{2} \le \frac{\delta}{4}$.
    Donc, $|x - y| < \delta$.
6.  Calculons la différence des images :
    $|f(x) - f(y)| = \left|\frac{1}{x} - \frac{1}{y}\right| = \left|\frac{1}{y/2} - \frac{1}{y}\right| = \left|\frac{2}{y} - \frac{1}{y}\right| = \left|\frac{1}{y}\right| = \frac{1}{y}$.
7.  Puisque $y = \min\left(1, \frac{\delta}{2}\right)$, nous avons $y \le 1$ et $y \le \frac{\delta}{2}$.
    Donc, $\frac{1}{y} \ge \frac{1}{\delta/2} = \frac{2}{\delta}$.
    Et $\frac{1}{y} \ge \frac{1}{1} = 1$.
    En particulier, $\frac{1}{y} \ge 1 = \epsilon_0$.
8.  Nous avons trouvé $x, y \in ]0, 1]$ tels que $|x - y| < \delta$ et $|f(x) - f(y)| \ge \epsilon_0$.
9.  Puisque ce raisonnement est valable pour tout $\delta > 0$, la fonction $f(x) = \frac{1}{x}$ n'est pas uniformément continue sur $]0, 1]$.

\textbf{Conclusion :} La fonction $f(x) = \frac{1}{x}$ est continue sur $]0, 1]$ mais n'est pas uniformément continue sur cet intervalle. L'intuition est que la pente de la fonction devient arbitrairement grande à l'approche de 0, ce qui empêche de trouver un $\delta$ unique pour toutes les paires de points.

\subsection{5. Application en Intelligence Artificielle}
\begin{itemize}
\item \textbf{Le Pont Théorique :} En Intelligence Artificielle, et plus particulièrement dans les réseaux de neurones profonds, les fonctions de transformation (activations, couches linéaires, convolutions) sont souvent choisies pour être continues, voire différentiables. La continuité garantit que de petites modifications des données d'entrée (par exemple, un léger bruit sur une image, une petite variation dans un texte) n'entraînent pas de changements brutaux et imprévisibles dans la prédiction ou la sortie du modèle. C'est la base de la \textbf{Robustesse} des modèles d'IA. Sans continuité, un modèle pourrait donner des résultats complètement différents pour des entrées presque identiques, le rendant inutilisable dans des applications critiques.
\item \textbf{Exemple Concret 1 : Robustesse et Adversarial Examples :}
\end{itemize}
    Dans la classification d'images, un réseau de neurones apprend une fonction $f: \mathbb{R}^d \to \mathbb{R}^k$ (où $d$ est la dimension de l'image et $k$ le nombre de classes). Si $f$ est continue, on s'attend à ce que $f(x)$ soit proche de $f(x')$ si $x$ est proche de $x'$. Cependant, les fonctions apprises par les réseaux de neurones sont souvent très complexes et peuvent être "localement non-Lipschitziennes" ou avoir des régions où la pente est extrêmement raide. Cela conduit au phénomène des \textbf{exemples adversariaux} : une perturbation infime et imperceptible à l'œil humain sur une image $x$ (créant $x'$) peut faire en sorte que le modèle classifie $x'$ dans une classe complètement différente de $x$. Bien que la fonction soit globalement continue, l'absence de continuité uniforme ou de bornes de Lipschitz globales peut rendre le modèle vulnérable. La recherche vise à rendre ces fonctions plus "lisses" ou plus robustes, souvent en imposant des contraintes de régularisation qui favorisent la continuité uniforme ou la Lipschitzianité.
\begin{itemize}
\item \textbf{Exemple Concret 2 : Génération d'images et Espace Latent :}
\end{itemize}
    Dans la génération d'images par des modèles comme les \textbf{GAN (Generative Adversarial Networks)} ou les \textbf{Auto-encodeurs variationnels (VAE)}, on manipule un "Espace Latent" (un espace de vecteurs de faible dimension). Le générateur (ou décodeur) est une fonction $G: \mathbb{R}^m \to \mathbb{R}^d$ qui transforme un vecteur latent $z$ en une image $G(z)$. On veut que cet espace latent soit continu : si on se déplace doucement entre le vecteur latent $z_1$ d'un "Chien" et le vecteur latent $z_2$ d'un "Chat", le décodeur doit générer une suite d'images qui se transforment graduellement de l'un à l'autre sans sauter d'une image à une autre sans transition logique. Si la fonction $G$ apprise par le réseau n'était pas continue, le modèle "hallucinerait" des images incohérentes ou des transitions abruptes au moindre petit changement de paramètre dans l'espace latent. La continuité du générateur est essentielle pour l'exploration fluide et l'interpolation dans l'espace des images.
\begin{itemize}
\item \textbf{Exemple Concret 3 : Optimisation et Convergence :}
\end{itemize}
    Les algorithmes d'apprentissage automatique reposent souvent sur l'optimisation de fonctions de coût. La continuité de ces fonctions de coût (et de leurs dérivées, si elles existent) est une propriété fondamentale qui garantit que les algorithmes d'optimisation (comme la descente de gradient) peuvent converger vers un minimum local. Si la fonction de coût présentait des discontinuités, l'algorithme pourrait "sauter" des minima ou ne jamais trouver de direction de descente stable.

\subsection{6. Liens Sémantiques}
\begin{itemize}
\item \textbf{Concepts Précédents requis :}
\end{itemize}
    *   [[Jalon 13 (Structure de R)]] : Compréhension des propriétés de $\mathbb{R}$ (ordre, complétude, intervalles, segments), qui sont fondamentales pour les définitions et théorèmes de continuité.
    *   [[Jalon 14 (Suites réelles et complexes)]] : Maîtrise des notions de convergence de suites, de suites adjacentes, de suites de Cauchy, et du théorème de Bolzano-Weierstrass, qui sont utilisées dans les démonstrations de la continuité et des théorèmes associés (TVI, Weierstrass, Heine).
\begin{itemize}
\item \textbf{Concepts Futurs dépendants :}
\end{itemize}
    *   [[Jalon 19 (Dérivabilité)]] : La dérivabilité est une condition plus forte que la continuité. Toute fonction dérivable est continue, mais la réciproque est fausse. La continuité est un prérequis essentiel à la notion de dérivabilité.
    *   [[Jalon 44 (Fonctions de plusieurs variables)]] : Les concepts de continuité s'étendent aux fonctions de plusieurs variables, nécessitant l'introduction de topologies plus générales (normes, métriques).
    *   [[Jalon 52 (Applications continues entre espaces topologiques et définition fine des homéomorphismes.)]] : La continuité est un concept central en topologie générale. La définition $\epsilon-\delta$ est une instanciation de la continuité dans les espaces métriques. Ce jalon généralisera la notion à des espaces topologiques arbitraires via les ouverts et les voisinages, et introduira les homéomorphismes comme bijections continues dont la réciproque est continue.
    *   \textbf{Intégration (Théorème fondamental de l'analyse) :} La continuité est une condition suffisante pour l'intégrabilité au sens de Riemann.
    *   \textbf{Équations Différentielles :} L'existence et l'unicité des solutions d'équations différentielles dépendent souvent de la continuité (et de la Lipschitzianité) des fonctions impliquées (Théorème de Cauchy-Lipschitz).

\chapter{Exercices d'application et de concours}
\section{Exo-01}
\section{Exercice 1 - Exploration de la Continuité des Fonctions d'une Variable Réelle}

Cet exercice est conçu pour évaluer votre compréhension approfondie de la continuité des fonctions d'une variable réelle, depuis la définition fondamentale jusqu'à l'application de théorèmes clés. La rigueur de votre rédaction est primordiale.

---

\subsection{Énoncé de l'Exercice}

\subsubsection{Partie A : Continuité par la Définition Epsilon-Delta}

Soit la fonction $f: \mathbb{R} \to \mathbb{R}$ définie par $f(x) = x\textasciicircum 2$.
En utilisant la définition formelle de la continuité (la définition $\epsilon-\delta$), démontrez que la fonction $f$ est continue au point $x_0 = 2$.

\subsubsection{Partie B : Continuité d'une Fonction Définie par Morceaux}

Considérons la fonction $g: \mathbb{R} \to \mathbb{R}$ définie par :
$$
g(x) = \begin{cases}
    \frac{x^2 - 4}{x - 2} & \text{si } x < 2 \\
    ax + b & \text{si } 2 \le x \le 3 \\
    \sqrt{x + 1} & \text{si } x > 3
\end{cases}
$$
Déterminez les valeurs des constantes réelles $a$ et $b$ pour que la fonction $g$ soit continue sur l'ensemble de tous les nombres réels $\mathbb{R}$.

\subsubsection{Partie C : Application du Théorème des Valeurs Intermédiaires}

Soit la fonction $h: \mathbb{R} \to \mathbb{R}$ définie par $h(x) = x^3 - 3x + 1$.
Démontrez que l'équation $h(x) = 0$ possède au moins trois racines réelles distinctes.

\subsubsection{Partie D : Continuité Epsilon-Delta pour une Fonction Rationnelle}

Soit la fonction $k: \mathbb{R}^* \to \mathbb{R}$ définie par $k(x) = \frac{1}{x}$.
En utilisant la définition formelle de la continuité ($\epsilon-\delta$), démontrez que la fonction $k$ est continue au point $x_0 = 1$.

---

\subsection{Correction Détaillée}

\subsubsection{Partie A : Continuité par la Définition Epsilon-Delta}

Nous voulons démontrer que la fonction $f(x) = x\textasciicircum 2$ est continue au point $x_0 = 2$.
La définition de la continuité d'une fonction $f$ au point $x_0$ est la suivante :
Pour tout $\epsilon > 0$, il existe un $\delta > 0$ tel que si $|x - x_0| < \delta$, alors $|f(x) - f(x_0)| < \epsilon$.

1.  \textbf{Identification de $f(x_0)$ :}
    Nous avons $x_0 = 2$, donc $f(x_0) = f(2) = 2^2 = 4$.

2.  \textbf{Expression de $|f(x) - f(x_0)|$ :}
    Nous devons analyser l'expression $|f(x) - f(2)| = |x^2 - 4|$.
    Par factorisation de la différence de carrés, nous obtenons :
    $$ |x^2 - 4| = |(x - 2)(x + 2)| = |x - 2| |x + 2| $$

3.  \textbf{Majoration du terme $|x+2|$ :}
    Nous voulons que $|x - 2|$ soit petit. Choisissons une première contrainte sur $\delta$. Supposons que $\delta \le 1$.
    Si $|x - 2| < \delta$ et $\delta \le 1$, alors $|x - 2| < 1$.
    Ceci implique que $-1 < x - 2 < 1$.
    En ajoutant $2$ à toutes les parties de l'inégalité, nous obtenons :
    $$ 2 - 1 < x < 2 + 1 $$
    $$ 1 < x < 3 $$
    Maintenant, nous voulons majorer $|x + 2|$. Puisque $1 < x < 3$, nous avons :
    $$ 1 + 2 < x + 2 < 3 + 2 $$
    $$ 3 < x + 2 < 5 $$
    Par conséquent, $|x + 2| < 5$.

4.  \textbf{Combinaison et choix de $\delta$ :}
    En utilisant la majoration précédente, nous avons :
    $$ |f(x) - f(2)| = |x - 2| |x + 2| < |x - 2| \cdot 5 $$
    Nous voulons que cette expression soit inférieure à $\epsilon$. Donc, nous voulons que $5|x - 2| < \epsilon$, ce qui implique $|x - 2| < \frac{\epsilon}{5}$.
    Nous avons deux contraintes sur $\delta$ : $\delta \le 1$ (pour majorer $|x+2|$) et $\delta \le \frac{\epsilon}{5}$ (pour obtenir la condition finale).
    Nous choisissons $\delta = \min\left(1, \frac{\epsilon}{5}\right)$.

5.  \textbf{Démonstration formelle :}
    Soit $\epsilon > 0$ un nombre réel arbitraire.
    Choisissons $\delta = \min\left(1, \frac{\epsilon}{5}\right)$. Par définition, $\delta > 0$.
    Supposons que $x$ est un nombre réel tel que $|x - 2| < \delta$.
    Puisque $\delta \le 1$, nous avons $|x - 2| < 1$.
    Ceci implique $-1 < x - 2 < 1$, ce qui donne $1 < x < 3$.
    En ajoutant $2$ à l'inégalité $1 < x < 3$, nous obtenons $3 < x + 2 < 5$.
    Par conséquent, $|x + 2| < 5$.
    Maintenant, considérons $|f(x) - f(2)|$:
    $$ |f(x) - f(2)| = |x^2 - 4| $$
    $$ |f(x) - f(2)| = |(x - 2)(x + 2)| \quad \text{(factorisation de la différence de carrés)} $$
    $$ |f(x) - f(2)| = |x - 2| |x + 2| \quad \text{(propriété du module : } |ab| = |a||b|) $$
    Puisque $|x - 2| < \delta$ et $|x + 2| < 5$, nous pouvons écrire :
    $$ |f(x) - f(2)| < \delta \cdot 5 $$
    De plus, par notre choix de $\delta$, nous savons que $\delta \le \frac{\epsilon}{5}$.
    Donc, $\delta \cdot 5 \le \left(\frac{\epsilon}{5}\right) \cdot 5 = \epsilon$.
    En combinant ces inégalités, nous obtenons :
    $$ |f(x) - f(2)| < \epsilon $$
    Ainsi, pour tout $\epsilon > 0$, il existe un $\delta = \min\left(1, \frac{\epsilon}{5}\right) > 0$ tel que si $|x - 2| < \delta$, alors $|f(x) - f(2)| < \epsilon$.
    Par conséquent, la fonction $f(x) = x\textasciicircum 2$ est continue au point $x_0 = 2$.

\subsubsection{Partie B : Continuité d'une Fonction Définie par Morceaux}

La fonction $g(x)$ est définie par morceaux. Pour qu'elle soit continue sur $\mathbb{R}$, elle doit être continue sur chaque intervalle ouvert de sa définition et aux points de raccordement ($x=2$ et $x=3$).

1.  \textbf{Continuité sur les intervalles ouverts :}
    *   \textbf{Pour $x < 2$ :} $g(x) = \frac{x^2 - 4}{x - 2}$. Pour $x \ne 2$, nous pouvons simplifier cette expression.
        $$ g(x) = \frac{(x - 2)(x + 2)}{x - 2} = x + 2 $$
        La fonction $x \mapsto x+2$ est une fonction polynomiale, et les fonctions polynomiales sont continues sur $\mathbb{R}$. Donc, $g$ est continue sur $(-\infty, 2)$.
    *   \textbf{Pour $2 < x < 3$ :} $g(x) = ax + b$. Cette fonction est une fonction polynomiale (linéaire), donc elle est continue sur $(2, 3)$.
    *   \textbf{Pour $x > 3$ :} $g(x) = \sqrt{x + 1}$. La fonction $u \mapsto \sqrt{u}$ est continue sur son domaine de définition $[0, +\infty)$. La fonction $x \mapsto x+1$ est continue sur $\mathbb{R}$. Pour $x > 3$, nous avons $x+1 > 4 > 0$. Par composition de fonctions continues, $g(x) = \sqrt{x+1}$ est continue sur $(3, +\infty)$.

2.  \textbf{Continuité aux points de raccordement :}
    Pour que $g$ soit continue en un point $x_0$, il faut que $\lim_{x \to x_0^-} g(x) = \lim_{x \to x_0^+} g(x) = g(x_0)$.

    *   \textbf{Au point $x_0 = 2$ :}
        *   \textbf{Valeur de la fonction en $x=2$ :} $g(2) = a(2) + b = 2a + b$.
        *   \textbf{Limite à gauche en $x=2$ :}
            $$ \lim_{x \to 2^-} g(x) = \lim_{x \to 2^-} \frac{x^2 - 4}{x - 2} $$
            Puisque $x \ne 2$ dans le calcul de la limite, nous pouvons factoriser le numérateur :
            $$ \lim_{x \to 2^-} \frac{(x - 2)(x + 2)}{x - 2} $$
            En simplifiant par $(x-2)$ (car $x \ne 2$), nous obtenons :
            $$ \lim_{x \to 2^-} (x + 2) $$
            La fonction $x \mapsto x+2$ est polynomiale, donc sa limite est obtenue par substitution directe :
            $$ \lim_{x \to 2^-} (x + 2) = 2 + 2 = 4 $$
        *   \textbf{Limite à droite en $x=2$ :}
            $$ \lim_{x \to 2^+} g(x) = \lim_{x \to 2^+} (ax + b) $$
            La fonction $x \mapsto ax+b$ est polynomiale, donc sa limite est obtenue par substitution directe :
            $$ \lim_{x \to 2^+} (ax + b) = a(2) + b = 2a + b $$
        *   \textbf{Condition de continuité en $x=2$ :} Pour que $g$ soit continue en $x=2$, il faut que $\lim_{x \to 2^-} g(x) = \lim_{x \to 2^+} g(x) = g(2)$.
            Donc, nous devons avoir $4 = 2a + b$. (Équation 1)

    *   \textbf{Au point $x_0 = 3$ :}
        *   \textbf{Valeur de la fonction en $x=3$ :} $g(3) = a(3) + b = 3a + b$.
        *   \textbf{Limite à gauche en $x=3$ :}
            $$ \lim_{x \to 3^-} g(x) = \lim_{x \to 3^-} (ax + b) $$
            La fonction $x \mapsto ax+b$ est polynomiale, donc sa limite est obtenue par substitution directe :
            $$ \lim_{x \to 3^-} (ax + b) = a(3) + b = 3a + b $$
        *   \textbf{Limite à droite en $x=3$ :}
            $$ \lim_{x \to 3^+} g(x) = \lim_{x \to 3^+} \sqrt{x + 1} $$
            La fonction $u \mapsto \sqrt{u}$ est continue pour $u \ge 0$. La fonction $x \mapsto x+1$ est continue. Pour $x \to 3^+$, $x+1 \to 4$. Puisque $4 > 0$, nous pouvons utiliser la continuité de la racine carrée :
            $$ \lim_{x \to 3^+} \sqrt{x + 1} = \sqrt{\lim_{x \to 3^+} (x + 1)} = \sqrt{3 + 1} = \sqrt{4} = 2 $$
        *   \textbf{Condition de continuité en $x=3$ :} Pour que $g$ soit continue en $x=3$, il faut que $\lim_{x \to 3^-} g(x) = \lim_{x \to 3^+} g(x) = g(3)$.
            Donc, nous devons avoir $3a + b = 2$. (Équation 2)

3.  \textbf{Résolution du système d'équations :}
    Nous avons un système de deux équations linéaires à deux inconnues $a$ et $b$ :
    1.  $2a + b = 4$
    2.  $3a + b = 2$

    Soustrayons l'Équation 1 de l'Équation 2 :
    $$ (3a + b) - (2a + b) = 2 - 4 $$
    $$ 3a - 2a + b - b = -2 $$
    $$ a = -2 $$
    Substituons la valeur de $a$ dans l'Équation 1 :
    $$ 2(-2) + b = 4 $$
    $$ -4 + b = 4 $$
    $$ b = 4 + 4 $$
    $$ b = 8 $$

    Par conséquent, pour que la fonction $g$ soit continue sur $\mathbb{R}$, les constantes doivent être $a = -2$ et $b = 8$.

\subsubsection{Partie C : Application du Théorème des Valeurs Intermédiaires}

Nous voulons démontrer que l'équation $h(x) = x^3 - 3x + 1 = 0$ possède au moins trois racines réelles distinctes.

1.  \textbf{Vérification de la continuité :}
    La fonction $h(x) = x^3 - 3x + 1$ est une fonction polynomiale. Les fonctions polynomiales sont continues sur l'ensemble de tous les nombres réels $\mathbb{R}$. Par conséquent, $h$ est continue sur tout intervalle fermé $[c, d] \subset \mathbb{R}$.

2.  \textbf{Recherche d'intervalles où la fonction change de signe :}
    Pour appliquer le Théorème des Valeurs Intermédiaires (TVI), nous devons trouver des intervalles $[c, d]$ sur lesquels $h$ est continue et où $h(c)$ et $h(d)$ ont des signes opposés. Si c'est le cas, le TVI garantit l'existence d'au moins une racine dans l'intervalle ouvert $(c, d)$.

    Calculons les valeurs de $h(x)$ pour quelques points :
    *   $h(-2) = (-2)^3 - 3(-2) + 1 = -8 + 6 + 1 = -1$.
    *   $h(0) = (0)^3 - 3(0) + 1 = 0 - 0 + 1 = 1$.
    *   $h(1) = (1)^3 - 3(1) + 1 = 1 - 3 + 1 = -1$.
    *   $h(2) = (2)^3 - 3(2) + 1 = 8 - 6 + 1 = 3$.

3.  \textbf{Application du Théorème des Valeurs Intermédiaires :}

    *   \textbf{Premier intervalle : $[-2, 0]$}
        1.  La fonction $h$ est continue sur l'intervalle fermé $[-2, 0]$ car elle est polynomiale.
        2.  Nous avons $h(-2) = -1$ et $h(0) = 1$.
        3.  Puisque $h(-2) < 0$ et $h(0) > 0$, et que $0$ est une valeur intermédiaire entre $h(-2)$ et $h(0)$, le Théorème des Valeurs Intermédiaires garantit qu'il existe au moins une racine $c_1 \in (-2, 0)$ telle que $h(c_1) = 0$.

    *   \textbf{Deuxième intervalle : $[0, 1]$}
        1.  La fonction $h$ est continue sur l'intervalle fermé $[0, 1]$ car elle est polynomiale.
        2.  Nous avons $h(0) = 1$ et $h(1) = -1$.
        3.  Puisque $h(0) > 0$ et $h(1) < 0$, et que $0$ est une valeur intermédiaire entre $h(0)$ et $h(1)$, le Théorème des Valeurs Intermédiaires garantit qu'il existe au moins une racine $c_2 \in (0, 1)$ telle que $h(c_2) = 0$.

    *   \textbf{Troisième intervalle : $[1, 2]$}
        1.  La fonction $h$ est continue sur l'intervalle fermé $[1, 2]$ car elle est polynomiale.
        2.  Nous avons $h(1) = -1$ et $h(2) = 3$.
        3.  Puisque $h(1) < 0$ et $h(2) > 0$, et que $0$ est une valeur intermédiaire entre $h(1)$ et $h(2)$, le Théorème des Valeurs Intermédiaires garantit qu'il existe au moins une racine $c_3 \in (1, 2)$ telle que $h(c_3) = 0$.

4.  \textbf{Conclusion :}
    Nous avons trouvé trois racines $c_1, c_2, c_3$ qui appartiennent respectivement aux intervalles $(-2, 0)$, $(0, 1)$ et $(1, 2)$. Ces intervalles sont disjoints.
    *   $c_1 \in (-2, 0)$
    *   $c_2 \in (0, 1)$
    *   $c_3 \in (1, 2)$
    Puisque ces racines se trouvent dans des intervalles disjoints, elles sont nécessairement distinctes.
    Par conséquent, l'équation $x^3 - 3x + 1 = 0$ possède au moins trois racines réelles distinctes.

\subsubsection{Partie D : Continuité Epsilon-Delta pour une Fonction Rationnelle}

Nous voulons démontrer que la fonction $k(x) = \frac{1}{x}$ est continue au point $x_0 = 1$.
La définition de la continuité d'une fonction $k$ au point $x_0$ est la suivante :
Pour tout $\epsilon > 0$, il existe un $\delta > 0$ tel que si $|x - x_0| < \delta$, alors $|k(x) - k(x_0)| < \epsilon$.

1.  \textbf{Identification de $k(x_0)$ :}
    Nous avons $x_0 = 1$, donc $k(x_0) = k(1) = \frac{1}{1} = 1$.

2.  \textbf{Expression de $|k(x) - k(x_0)|$ :}
    Nous devons analyser l'expression $|k(x) - k(1)| = \left|\frac{1}{x} - 1\right|$.
    Pour combiner les termes, nous mettons sur un dénominateur commun :
    $$ \left|\frac{1}{x} - 1\right| = \left|\frac{1 - x}{x}\right| $$
    En utilisant la propriété du module $|a/b| = |a|/|b|$ et $|1-x| = |x-1|$ :
    $$ \left|\frac{1 - x}{x}\right| = \frac{|1 - x|}{|x|} = \frac{|x - 1|}{|x|} $$

3.  \textbf{Majoration du terme $\frac{1}{|x|}$ :}
    Nous voulons que $|x - 1|$ soit petit. Choisissons une première contrainte sur $\delta$. Supposons que $\delta \le \frac{1}{2}$.
    Si $|x - 1| < \delta$ et $\delta \le \frac{1}{2}$, alors $|x - 1| < \frac{1}{2}$.
    Ceci implique $-\frac{1}{2} < x - 1 < \frac{1}{2}$.
    En ajoutant $1$ à toutes les parties de l'inégalité, nous obtenons :
    $$ 1 - \frac{1}{2} < x < 1 + \frac{1}{2} $$
    $$ \frac{1}{2} < x < \frac{3}{2} $$
    Puisque $x > \frac{1}{2}$, nous avons $0 < \frac{1}{x} < 2$.
    Par conséquent, $\frac{1}{|x|} < 2$.

4.  \textbf{Combinaison et choix de $\delta$ :}
    En utilisant la majoration précédente, nous avons :
    $$ |k(x) - k(1)| = \frac{|x - 1|}{|x|} < |x - 1| \cdot 2 $$
    Nous voulons que cette expression soit inférieure à $\epsilon$. Donc, nous voulons que $2|x - 1| < \epsilon$, ce qui implique $|x - 1| < \frac{\epsilon}{2}$.
    Nous avons deux contraintes sur $\delta$ : $\delta \le \frac{1}{2}$ (pour majorer $\frac{1}{|x|}$) et $\delta \le \frac{\epsilon}{2}$ (pour obtenir la condition finale).
    Nous choisissons $\delta = \min\left(\frac{1}{2}, \frac{\epsilon}{2}\right)$.

5.  \textbf{Démonstration formelle :}
    Soit $\epsilon > 0$ un nombre réel arbitraire.
    Choisissons $\delta = \min\left(\frac{1}{2}, \frac{\epsilon}{2}\right)$. Par définition, $\delta > 0$.
    Supposons que $x$ est un nombre réel tel que $|x - 1| < \delta$.
    Puisque $\delta \le \frac{1}{2}$, nous avons $|x - 1| < \frac{1}{2}$.
    Ceci implique $-\frac{1}{2} < x - 1 < \frac{1}{2}$, ce qui donne $\frac{1}{2} < x < \frac{3}{2}$.
    Puisque $x > \frac{1}{2}$, $x$ est non nul, et nous pouvons prendre l'inverse :
    $$ \frac{1}{3/2} < \frac{1}{x} < \frac{1}{1/2} $$
    $$ \frac{2}{3} < \frac{1}{x} < 2 $$
    Par conséquent, $\frac{1}{|x|} < 2$.
    Maintenant, considérons $|k(x) - k(1)|$:
    $$ |k(x) - k(1)| = \left|\frac{1}{x} - 1\right| $$
    $$ |k(x) - k(1)| = \left|\frac{1 - x}{x}\right| \quad \text{(mise au même dénominateur)} $$
    $$ |k(x) - k(1)| = \frac{|1 - x|}{|x|} \quad \text{(propriété du module : } |a/b| = |a|/|b|) $$
    $$ |k(x) - k(1)| = \frac{|x - 1|}{|x|} \quad \text{(propriété du module : } |1-x| = |x-1|) $$
    Puisque $|x - 1| < \delta$ et $\frac{1}{|x|} < 2$, nous pouvons écrire :
    $$ |k(x) - k(1)| < \delta \cdot 2 $$
    De plus, par notre choix de $\delta$, nous savons que $\delta \le \frac{\epsilon}{2}$.
    Donc, $\delta \cdot 2 \le \left(\frac{\epsilon}{2}\right) \cdot 2 = \epsilon$.
    En combinant ces inégalités, nous obtenons :
    $$ |k(x) - k(1)| < \epsilon $$
    Ainsi, pour tout $\epsilon > 0$, il existe un $\delta = \min\left(\frac{1}{2}, \frac{\epsilon}{2}\right) > 0$ tel que si $|x - 1| < \delta$, alors $|k(x) - k(1)| < \epsilon$.
    Par conséquent, la fonction $k(x) = \frac{1}{x}$ est continue au point $x_0 = 1$.

\section{Exo-02}
\section{Exercice 2 - Exploration de la Continuité des Fonctions Réelles}

Ce jalon explore la notion fondamentale de continuité pour les fonctions d'une variable réelle. Nous allons aborder la définition formelle, l'analyse de fonctions définies par morceaux et l'application de théorèmes clés.

---

\subsection{Énoncé de l'Exercice}

\textbf{Partie A : Continuité par la définition $\epsilon-\delta$}

Soit la fonction $f: \mathbb{R} \to \mathbb{R}$ définie par $f(x) = x\textasciicircum 2 + 3x$.
En utilisant la définition formelle de la continuité ($\epsilon-\delta$), démontrez que $f$ est continue au point $x_0 = 1$.

\textbf{Partie B : Continuité d'une fonction définie par morceaux}

Considérons la fonction $g: \mathbb{R} \to \mathbb{R}$ définie par :
$$ g(x) = \begin{cases} \frac{x^2 - 4}{x - 2} & \text{si } x \neq 2 \\ k & \text{si } x = 2 \end{cases} $$
Déterminez la valeur de la constante $k$ pour laquelle la fonction $g$ est continue au point $x = 2$.

\textbf{Partie C : Continuité sur un intervalle et Théorème des Valeurs Intermédiaires}

Soit la fonction $h: D \to \mathbb{R}$ définie par $h(x) = \sqrt{x+1}$.

1.  Déterminez le domaine de définition $D$ de la fonction $h$.
2.  Démontrez que la fonction $h$ est continue sur son domaine $D$.
3.  En utilisant le Théorème des Valeurs Intermédiaires (TVI), montrez qu'il existe au moins une valeur $c \in [0, 3]$ telle que $h(c) = 2$.

---

\subsection{Correction Détaillée}

\subsubsection{Partie A : Continuité par la définition $\epsilon-\delta$}

Nous devons démontrer que la fonction $f(x) = x\textasciicircum 2 + 3x$ est continue au point $x_0 = 1$ en utilisant la définition $\epsilon-\delta$.
La fonction $f$ est continue en $x_0 = 1$ si et seulement si pour tout $\epsilon > 0$, il existe un $\delta > 0$ tel que pour tout $x \in \mathbb{R}$, si $|x - 1| < \delta$, alors $|f(x) - f(1)| < \epsilon$.

\textbf{Étape 1 : Calcul de $f(x_0)$}
Nous calculons la valeur de la fonction au point $x_0 = 1$:
$$ f(1) = (1)^2 + 3(1) = 1 + 3 = 4 $$

\textbf{Étape 2 : Analyse de l'expression $|f(x) - f(1)|$}
Nous considérons la différence $|f(x) - f(1)|$:
$$ |f(x) - f(1)| = |(x^2 + 3x) - 4| $$
Nous factorisons le polynôme $x^2 + 3x - 4$. Nous cherchons deux nombres dont la somme est $3$ et le produit est $-4$. Ces nombres sont $4$ et $-1$.
Ainsi, $x^2 + 3x - 4 = (x-1)(x+4)$.
L'expression devient :
$$ |f(x) - f(1)| = |(x-1)(x+4)| = |x-1||x+4| $$

\textbf{Étape 3 : Borner le terme $|x+4|$}
Nous voulons que $|f(x) - f(1)| < \epsilon$. Nous avons le terme $|x-1|$ que nous pouvons contrôler directement avec $\delta$. Le terme $|x+4|$ doit être borné.
Choisissons une première valeur pour $\delta$, par exemple $\delta_1 = 1$.
Si $|x-1| < \delta_1 = 1$, cela signifie que $-1 < x-1 < 1$.
En ajoutant $1$ à toutes les parties de l'inégalité, nous obtenons :
$$ 0 < x < 2 $$
Maintenant, nous voulons borner $|x+4|$. Si $0 < x < 2$, alors en ajoutant $4$ à toutes les parties de l'inégalité :
$$ 0+4 < x+4 < 2+4 $$
$$ 4 < x+4 < 6 $$
Puisque $x+4$ est positif dans cet intervalle, $|x+4| = x+4$.
Donc, nous avons $|x+4| < 6$.

\textbf{Étape 4 : Détermination de $\delta$ en fonction de $\epsilon$}
En utilisant la borne trouvée pour $|x+4|$, nous avons :
$$ |f(x) - f(1)| = |x-1||x+4| < |x-1| \cdot 6 $$
Nous voulons que cette expression soit inférieure à $\epsilon$:
$$ 6|x-1| < \epsilon $$
Ceci implique :
$$ |x-1| < \frac{\epsilon}{6} $$
Pour que les deux conditions $|x-1| < 1$ (pour borner $|x+4|$) et $|x-1| < \frac{\epsilon}{6}$ soient satisfaites simultanément, nous devons choisir $\delta$ comme le minimum de ces deux valeurs.
Soit $\delta = \min\left(1, \frac{\epsilon}{6}\right)$.

\textbf{Étape 5 : Conclusion de la démonstration}
Pour tout $\epsilon > 0$, nous avons trouvé un $\delta = \min\left(1, \frac{\epsilon}{6}\right) > 0$.
Si $|x-1| < \delta$, alors :
1.  $|x-1| < 1$, ce qui implique $0 < x < 2$, et donc $4 < x+4 < 6$, d'où $|x+4| < 6$.
2.  $|x-1| < \frac{\epsilon}{6}$.
En combinant ces deux inégalités, nous obtenons :
$$ |f(x) - f(1)| = |x-1||x+4| < \left(\frac{\epsilon}{6}\right) \cdot 6 = \epsilon $$
Par conséquent, la fonction $f(x) = x\textasciicircum 2 + 3x$ est continue au point $x_0 = 1$ selon la définition $\epsilon-\delta$.

\subsubsection{Partie B : Continuité d'une fonction définie par morceaux}

Pour que la fonction $g(x)$ soit continue au point $x = 2$, il faut que la limite de $g(x)$ lorsque $x$ tend vers $2$ soit égale à la valeur de la fonction en $x = 2$. C'est-à-dire :
$$ \lim_{x \to 2} g(x) = g(2) $$

\textbf{Étape 1 : Calcul de $g(2)$}
Par la définition de la fonction $g(x)$, lorsque $x = 2$, la valeur de la fonction est $k$.
$$ g(2) = k $$

\textbf{Étape 2 : Calcul de $\lim_{x \to 2} g(x)$}
Pour $x \neq 2$, la fonction $g(x)$ est définie par $g(x) = \frac{x^2 - 4}{x - 2}$.
Nous calculons la limite de cette expression lorsque $x$ tend vers $2$:
$$ \lim_{x \to 2} g(x) = \lim_{x \to 2} \frac{x^2 - 4}{x - 2} $$
Le numérateur $x^2 - 4$ est une différence de carrés, qui peut être factorisée comme $(x-2)(x+2)$.
$$ \lim_{x \to 2} \frac{(x-2)(x+2)}{x - 2} $$
Puisque $x \to 2$ signifie que $x$ s'approche de $2$ mais n'est pas égal à $2$, le terme $(x-2)$ est non nul. Nous pouvons donc simplifier l'expression en divisant le numérateur et le dénominateur par $(x-2)$:
$$ \lim_{x \to 2} (x+2) $$
Cette limite est celle d'un polynôme. Par les propriétés des limites (la limite d'une somme est la somme des limites, la limite d'une constante est la constante, la limite de $x$ est $x_0$), nous pouvons substituer $x=2$:
$$ \lim_{x \to 2} (x+2) = 2 + 2 = 4 $$

\textbf{Étape 3 : Détermination de $k$ pour la continuité}
Pour que $g(x)$ soit continue en $x=2$, nous devons avoir $\lim_{x \to 2} g(x) = g(2)$.
En substituant les valeurs calculées :
$$ 4 = k $$
Par conséquent, la fonction $g(x)$ est continue au point $x=2$ si et seulement si $k=4$.

\subsubsection{Partie C : Continuité sur un intervalle et Théorème des Valeurs Intermédiaires}

Soit la fonction $h(x) = \sqrt{x+1}$.

\textbf{1. Détermination du domaine de définition $D$ de la fonction $h$}

Pour que la fonction $h(x) = \sqrt{x+1}$ soit définie dans l'ensemble des nombres réels, l'expression sous le radical (le radicande) doit être supérieure ou égale à zéro.
$$ x+1 \ge 0 $$
En soustrayant $1$ des deux côtés de l'inégalité, nous obtenons :
$$ x \ge -1 $$
Le domaine de définition $D$ de la fonction $h$ est l'ensemble de tous les nombres réels $x$ tels que $x \ge -1$.
$$ D = [-1, \infty) $$

\textbf{2. Démonstration de la continuité de $h$ sur son domaine $D$}

Nous allons utiliser les propriétés de continuité des fonctions élémentaires et des compositions de fonctions.
Considérons deux fonctions :
*   $f_1(x) = x+1$. Cette fonction est un polynôme de degré 1. Les fonctions polynomiales sont continues sur tout $\mathbb{R}$. Par conséquent, $f_1(x)$ est continue sur $[-1, \infty)$.
*   $f_2(y) = \sqrt{y}$. Cette fonction est la fonction racine carrée. La fonction racine carrée est continue sur son domaine de définition, qui est $[0, \infty)$.

La fonction $h(x)$ est la composition de $f_2$ et $f_1$, c'est-à-dire $h(x) = f_2(f_1(x))$.
Pour qu'une composition de fonctions $f_2 \circ f_1$ soit continue, il faut que :
a.  $f_1$ soit continue au point $x$.
b.  $f_2$ soit continue au point $f_1(x)$.

Pour tout $x \in D = [-1, \infty)$, nous avons $x+1 \ge 0$.
Donc, pour tout $x \in D$, la valeur $f_1(x) = x+1$ appartient au domaine de définition de $f_2$, qui est $[0, \infty)$.
Puisque $f_1(x)$ est continue sur $[-1, \infty)$ et $f_2(y)$ est continue sur $[0, \infty)$, la composition $h(x) = f_2(f_1(x)) = \sqrt{x+1}$ est continue sur son domaine $D = [-1, \infty)$.

\textbf{3. Application du Théorème des Valeurs Intermédiaires (TVI)}

Nous voulons montrer qu'il existe au moins une valeur $c \in [0, 3]$ telle que $h(c) = 2$.
Pour appliquer le Théorème des Valeurs Intermédiaires, nous devons vérifier trois conditions :

\textbf{Condition 1 : La fonction doit être continue sur l'intervalle fermé $[a, b]$.}
L'intervalle donné est $[0, 3]$.
D'après la partie C.2, la fonction $h(x) = \sqrt{x+1}$ est continue sur son domaine $D = [-1, \infty)$.
Puisque l'intervalle $[0, 3]$ est un sous-ensemble de $[-1, \infty)$, la fonction $h(x)$ est continue sur l'intervalle fermé $[0, 3]$.

\textbf{Condition 2 : Calculer les valeurs de la fonction aux bornes de l'intervalle.}
Nous calculons $h(0)$ et $h(3)$:
$$ h(0) = \sqrt{0+1} = \sqrt{1} = 1 $$
$$ h(3) = \sqrt{3+1} = \sqrt{4} = 2 $$

\textbf{Condition 3 : Vérifier que la valeur cible $y_0$ est comprise entre $h(a)$ et $h(b)$.}
La valeur cible est $y_0 = 2$.
Les valeurs aux bornes sont $h(0) = 1$ et $h(3) = 2$.
Nous observons que $1 \le 2 \le 2$.
La valeur $y_0 = 2$ est bien comprise dans l'intervalle $[h(0), h(3)] = [1, 2]$.

\textbf{Conclusion par le TVI :}
Puisque toutes les conditions du Théorème des Valeurs Intermédiaires sont satisfaites (la fonction $h$ est continue sur $[0, 3]$ et la valeur $2$ est entre $h(0)$ et $h(3)$), le TVI garantit qu'il existe au moins une valeur $c \in [0, 3]$ telle que $h(c) = 2$.
Dans ce cas précis, nous avons même trouvé que $h(3) = 2$, donc $c=3$ est une telle valeur.

\section{Exo-03}
\section{Exercice 3 - Exploration Approfondie de la Continuité}

Cet exercice vise à consolider votre compréhension de la continuité des fonctions d'une variable réelle, depuis sa définition formelle jusqu'à ses applications fondamentales.

\textbf{Partie A : Maîtrise de la définition formelle de la continuité.}

Soit la fonction $f$ définie sur $\mathbb{R}$ par $f(x) = x\textasciicircum 2 - 2x + 3$.
En utilisant la définition $\epsilon-\delta$ de la continuité, démontrez que $f$ est continue au point $x_0 = 1$.

\textbf{Partie B : Analyse de la continuité sur un intervalle.}

Considérons la fonction $g$ définie par morceaux sur $\mathbb{R}$ comme suit :
$$ g(x) = \begin{cases} \frac{x^2 - 4}{x - 2} & \text{si } x < 2 \\ ax + b & \text{si } 2 \le x \le 3 \\ \sqrt{x^2 + 7} & \text{si } x > 3 \end{cases} $$
Déterminez les valeurs des constantes réelles $a$ et $b$ pour que la fonction $g$ soit continue sur l'ensemble de son domaine de définition, $\mathbb{R}$.

\textbf{Partie C : Application du Théorème des Valeurs Intermédiaires.}

Soit l'équation $(E) : x^3 - 3x + 1 = 0$.
1.  Montrez que l'équation $(E)$ possède au moins une solution réelle dans l'intervalle $(0, 1)$.
2.  Montrez que l'équation $(E)$ possède au moins une solution réelle dans l'intervalle $(-2, -1)$.
3.  En déduire que l'équation $(E)$ possède au moins deux solutions réelles distinctes.

---

\section{Correction de l'Exercice 3 - Exploration Approfondie de la Continuité}

\subsection{Partie A : Maîtrise de la définition formelle de la continuité.}

Nous devons démontrer que la fonction $f(x) = x\textasciicircum 2 - 2x + 3$ est continue au point $x_0 = 1$ en utilisant la définition $\epsilon-\delta$.
La fonction $f$ est continue en $x_0 = 1$ si et seulement si pour tout $\epsilon > 0$, il existe un $\delta > 0$ tel que pour tout $x \in \mathbb{R}$, si $|x - 1| < \delta$, alors $|f(x) - f(1)| < \epsilon$.

Calculons d'abord la valeur de la fonction $f$ au point $x_0 = 1$:
$$ f(1) = (1)^2 - 2(1) + 3 = 1 - 2 + 3 = 2 $$

Nous devons donc montrer que pour tout $\epsilon > 0$, il existe un $\delta > 0$ tel que si $|x - 1| < \delta$, alors $|(x^2 - 2x + 3) - 2| < \epsilon$.
Simplifions l'expression $|f(x) - f(1)|$:
$$ |f(x) - f(1)| = |x^2 - 2x + 3 - 2| $$
$$ |f(x) - f(1)| = |x^2 - 2x + 1| $$
Nous reconnaissons l'identité remarquable $A^2 - 2AB + B^2 = (A - B)^2$. Ici, $A=x$ et $B=1$.
Donc, nous pouvons factoriser l'expression $x^2 - 2x + 1$ comme suit :
$$ x^2 - 2x + 1 = (x - 1)^2 $$
En substituant cette factorisation dans l'expression de la valeur absolue, nous obtenons :
$$ |f(x) - f(1)| = |(x - 1)^2| $$
Puisque le carré d'un nombre réel est toujours positif ou nul, c'est-à-dire $(x-1)^2 \ge 0$ pour tout $x \in \mathbb{R}$, la valeur absolue de $(x-1)^2$ est $(x-1)^2$ elle-même.
$$ |f(x) - f(1)| = (x - 1)^2 $$
Nous voulons que cette quantité soit inférieure à $\epsilon$, c'est-à-dire $(x - 1)^2 < \epsilon$.
Pour que $(x - 1)^2 < \epsilon$, il est nécessaire et suffisant que $|x - 1| < \sqrt{\epsilon}$.

Soit $\epsilon > 0$ un nombre réel arbitrairement choisi.
Nous cherchons un $\delta > 0$ tel que si $|x - 1| < \delta$, alors $|f(x) - f(1)| < \epsilon$.
Choisissons $\delta = \sqrt{\epsilon}$. Puisque $\epsilon > 0$, $\sqrt{\epsilon}$ est un nombre réel positif, donc $\delta > 0$.
Maintenant, supposons que $|x - 1| < \delta$.
En substituant la valeur de $\delta$, nous avons :
$$ |x - 1| < \sqrt{\epsilon} $$
Puisque les deux membres de cette inégalité sont positifs (la valeur absolue est non négative et $\sqrt{\epsilon}$ est positif), nous pouvons élever les deux membres au carré sans changer le sens de l'inégalité :
$$ (|x - 1|)^2 < (\sqrt{\epsilon})^2 $$
$$ (x - 1)^2 < \epsilon $$
Comme nous l'avons établi précédemment, $|f(x) - f(1)| = (x - 1)^2$.
Donc, nous avons :
$$ |f(x) - f(1)| < \epsilon $$
Nous avons ainsi démontré que pour tout $\epsilon > 0$, il existe un $\delta = \sqrt{\epsilon} > 0$ tel que si $|x - 1| < \delta$, alors $|f(x) - f(1)| < \epsilon$.
Par conséquent, la fonction $f$ est continue au point $x_0 = 1$ selon la définition $\epsilon-\delta$.

\subsection{Partie B : Analyse de la continuité sur un intervalle.}

La fonction $g$ est définie par morceaux. Pour qu'elle soit continue sur $\mathbb{R}$, elle doit être continue sur chaque intervalle ouvert où elle est définie par une seule expression, et elle doit être continue aux points de raccordement.

1.  \textbf{Continuité sur les intervalles ouverts :}
    *   Pour $x < 2$, $g(x) = \frac{x^2 - 4}{x - 2}$. Cette fonction est un quotient de deux polynômes, $P(x) = x^2 - 4$ et $Q(x) = x - 2$. Les polynômes sont des fonctions continues sur $\mathbb{R}$. Le dénominateur $Q(x) = x - 2$ est non nul pour $x < 2$. Par conséquent, $g$ est continue sur l'intervalle $(-\infty, 2)$ en tant que quotient de fonctions continues dont le dénominateur ne s'annule pas.
    *   Pour $2 < x < 3$, $g(x) = ax + b$. Cette fonction est un polynôme de degré 1. Les polynômes sont des fonctions continues sur $\mathbb{R}$. Par conséquent, $g$ est continue sur l'intervalle $(2, 3)$.
    *   Pour $x > 3$, $g(x) = \sqrt{x^2 + 7}$. Cette fonction est une composition de deux fonctions : $h_1(x) = x^2 + 7$ et $h_2(y) = \sqrt{y}$. La fonction $h_1(x)$ est un polynôme, donc continue sur $\mathbb{R}$. Pour $x > 3$, $x^2 > 9$, donc $x^2 + 7 > 16 > 0$. La fonction $h_2(y) = \sqrt{y}$ est continue pour tout $y > 0$. Puisque $h_1(x) > 0$ pour $x > 3$, la composition $g(x) = h_2(h_1(x))$ est continue sur l'intervalle $(3, +\infty)$.

2.  \textbf{Continuité aux points de raccordement :}
    Nous devons assurer la continuité aux points $x = 2$ et $x = 3$. Pour qu'une fonction $g$ soit continue en un point $x_0$, il faut que la limite de $g(x)$ lorsque $x$ tend vers $x_0$ existe et soit égale à $g(x_0)$. Cela implique que la limite à gauche et la limite à droite doivent exister et être égales à $g(x_0)$. C'est-à-dire $\lim_{x \to x_0^-} g(x) = \lim_{x \to x_0^+} g(x) = g(x_0)$.

    \textbf{Au point $x_0 = 2$ :}
    *   Calculons la limite à gauche de $g(x)$ lorsque $x$ tend vers $2$ :
        $$ \lim_{x \to 2^-} g(x) = \lim_{x \to 2^-} \frac{x^2 - 4}{x - 2} $$
        Nous observons une forme indéterminée de type $\frac{0}{0}$. Nous factorisons le numérateur en utilisant l'identité remarquable $A^2 - B^2 = (A - B)(A + B)$, où $A=x$ et $B=2$.
        $$ x^2 - 4 = (x - 2)(x + 2) $$
        Substituons cette factorisation dans la limite :
        $$ \lim_{x \to 2^-} \frac{(x - 2)(x + 2)}{x - 2} $$
        Pour $x \ne 2$ (ce qui est le cas lorsque $x \to 2^-$), nous pouvons simplifier le terme $(x - 2)$ au numérateur et au dénominateur :
        $$ \lim_{x \to 2^-} (x + 2) $$
        La fonction $x+2$ est un polynôme, donc elle est continue. Nous pouvons évaluer la limite par substitution directe :
        $$ \lim_{x \to 2^-} (x + 2) = 2 + 2 = 4 $$
    *   Calculons la limite à droite de $g(x)$ lorsque $x$ tend vers $2$ :
        $$ \lim_{x \to 2^+} g(x) = \lim_{x \to 2^+} (ax + b) $$
        La fonction $ax+b$ est un polynôme, donc elle est continue. Nous pouvons évaluer la limite par substitution directe :
        $$ \lim_{x \to 2^+} (ax + b) = a(2) + b = 2a + b $$
    *   Calculons la valeur de la fonction $g$ au point $x = 2$ :
        Selon la définition de $g(x)$, pour $x = 2$, nous utilisons la deuxième expression :
        $$ g(2) = a(2) + b = 2a + b $$
    Pour que $g$ soit continue en $x = 2$, les trois valeurs doivent être égales :
    $$ \lim_{x \to 2^-} g(x) = \lim_{x \to 2^+} g(x) = g(2) $$
    Cela nous donne la première équation :
    $$ 4 = 2a + b \quad (1) $$

    \textbf{Au point $x_0 = 3$ :}
    *   Calculons la limite à gauche de $g(x)$ lorsque $x$ tend vers $3$ :
        $$ \lim_{x \to 3^-} g(x) = \lim_{x \to 3^-} (ax + b) $$
        La fonction $ax+b$ est un polynôme, donc elle est continue. Nous pouvons évaluer la limite par substitution directe :
        $$ \lim_{x \to 3^-} (ax + b) = a(3) + b = 3a + b $$
    *   Calculons la limite à droite de $g(x)$ lorsque $x$ tend vers $3$ :
        $$ \lim_{x \to 3^+} g(x) = \lim_{x \to 3^+} \sqrt{x^2 + 7} $$
        La fonction $\sqrt{x^2+7}$ est continue pour $x > 3$. Nous pouvons évaluer la limite par substitution directe :
        $$ \lim_{x \to 3^+} \sqrt{x^2 + 7} = \sqrt{(3)^2 + 7} = \sqrt{9 + 7} = \sqrt{16} = 4 $$
    *   Calculons la valeur de la fonction $g$ au point $x = 3$ :
        Selon la définition de $g(x)$, pour $x = 3$, nous utilisons la deuxième expression :
        $$ g(3) = a(3) + b = 3a + b $$
    Pour que $g$ soit continue en $x = 3$, les trois valeurs doivent être égales :
    $$ \lim_{x \to 3^-} g(x) = \lim_{x \to 3^+} g(x) = g(3) $$
    Cela nous donne la deuxième équation :
    $$ 3a + b = 4 \quad (2) $$

    Nous avons maintenant un système de deux équations linéaires à deux inconnues $a$ et $b$ :
    (1) $2a + b = 4$
    (2) $3a + b = 4$

    Pour résoudre ce système, nous pouvons soustraire l'équation (1) de l'équation (2) :
    $$ (3a + b) - (2a + b) = 4 - 4 $$
    $$ 3a - 2a + b - b = 0 $$
    $$ a = 0 $$
    Maintenant que nous avons la valeur de $a$, nous la substituons dans l'équation (1) pour trouver $b$ :
    $$ 2(0) + b = 4 $$
    $$ 0 + b = 4 $$
    $$ b = 4 $$

    Par conséquent, pour que la fonction $g$ soit continue sur l'ensemble de son domaine de définition $\mathbb{R}$, les constantes $a$ et $b$ doivent prendre les valeurs $a = 0$ et $b = 4$.

\subsection{Partie C : Application du Théorème des Valeurs Intermédiaires.}

Soit l'équation $(E) : x^3 - 3x + 1 = 0$.
Considérons la fonction $h(x) = x^3 - 3x + 1$.
La fonction $h(x)$ est un polynôme. Les polynômes sont des fonctions continues sur l'ensemble de tous les nombres réels, $\mathbb{R}$. Par conséquent, $h(x)$ est continue sur tout intervalle fermé $[c, d]$ inclus dans $\mathbb{R}$.

1.  \textbf{Montrons que l'équation $(E)$ possède au moins une solution réelle dans l'intervalle $(0, 1)$.}
    Nous allons appliquer le Théorème des Valeurs Intermédiaires (TVI).
    *   \textbf{Condition 1 (Continuité) :} La fonction $h(x) = x^3 - 3x + 1$ est continue sur l'intervalle fermé $[0, 1]$ car c'est un polynôme.
    *   \textbf{Condition 2 (Valeurs aux bornes) :} Calculons les valeurs de $h$ aux bornes de l'intervalle $[0, 1]$ :
        $$ h(0) = (0)^3 - 3(0) + 1 = 0 - 0 + 1 = 1 $$
        $$ h(1) = (1)^3 - 3(1) + 1 = 1 - 3 + 1 = -1 $$
    *   \textbf{Condition 3 (Existence de $k$) :} Nous observons que $h(0) = 1$ et $h(1) = -1$. Puisque $h(1) < 0 < h(0)$, la valeur $k=0$ est comprise entre $h(1)$ et $h(0)$.
    *   \textbf{Conclusion par le TVI :} Le Théorème des Valeurs Intermédiaires stipule que si une fonction $h$ est continue sur un intervalle fermé $[c, d]$ et si $k$ est un nombre réel compris entre $h(c)$ et $h(d)$, alors il existe au moins un $x_1 \in (c, d)$ tel que $h(x_1) = k$.
        Dans notre cas, $c=0$, $d=1$, et $k=0$. Toutes les conditions du TVI sont remplies.
        Par conséquent, il existe au moins une solution $x_1 \in (0, 1)$ telle que $h(x_1) = 0$.
    Ainsi, l'équation $(E)$ possède au moins une solution réelle dans l'intervalle $(0, 1)$.

2.  \textbf{Montrons que l'équation $(E)$ possède au moins une solution réelle dans l'intervalle $(-2, -1)$.}
    Nous allons appliquer le Théorème des Valeurs Intermédiaires (TVI) de la même manière.
    *   \textbf{Condition 1 (Continuité) :} La fonction $h(x) = x^3 - 3x + 1$ est continue sur l'intervalle fermé $[-2, -1]$ car c'est un polynôme.
    *   \textbf{Condition 2 (Valeurs aux bornes) :} Calculons les valeurs de $h$ aux bornes de l'intervalle $[-2, -1]$ :
        $$ h(-2) = (-2)^3 - 3(-2) + 1 = -8 + 6 + 1 = -1 $$
        $$ h(-1) = (-1)^3 - 3(-1) + 1 = -1 + 3 + 1 = 3 $$
    *   \textbf{Condition 3 (Existence de $k$) :} Nous observons que $h(-2) = -1$ et $h(-1) = 3$. Puisque $h(-2) < 0 < h(-1)$, la valeur $k=0$ est comprise entre $h(-2)$ et $h(-1)$.
    *   \textbf{Conclusion par le TVI :} Toutes les conditions du TVI sont remplies pour l'intervalle $[-2, -1]$ et $k=0$.
        Par conséquent, il existe au moins une solution $x_2 \in (-2, -1)$ telle que $h(x_2) = 0$.
    Ainsi, l'équation $(E)$ possède au moins une solution réelle dans l'intervalle $(-2, -1)$.

3.  \textbf{En déduire que l'équation $(E)$ possède au moins deux solutions réelles distinctes.}
    D'après la question 1, nous avons montré l'existence d'une solution $x_1$ telle que $h(x_1) = 0$ et $x_1 \in (0, 1)$.
    D'après la question 2, nous avons montré l'existence d'une solution $x_2$ telle que $h(x_2) = 0$ et $x_2 \in (-2, -1)$.
    Les intervalles $(0, 1)$ et $(-2, -1)$ sont des intervalles disjoints.
    En effet, l'intervalle $(0, 1)$ contient des nombres strictement positifs, tandis que l'intervalle $(-2, -1)$ contient des nombres strictement négatifs. Il n'y a aucun nombre commun à ces deux intervalles, donc leur intersection est vide : $(0, 1) \cap (-2, -1) = \emptyset$.
    Puisque $x_1$ appartient à l'intervalle $(0, 1)$ et $x_2$ appartient à l'intervalle $(-2, -1)$, il est impossible que $x_1$ soit égal à $x_2$.
    Par conséquent, $x_1$ et $x_2$ sont deux solutions réelles distinctes de l'équation $(E)$.
    L'équation $(E)$ possède donc au moins deux solutions réelles distinctes.

\section{Exo-04}
\section{Exercice 4 - Continuité et Propriétés Fondamentales des Fonctions Réelles}

\subsection{Partie 1 : Analyse de la continuité ponctuelle et globale}

Soit $f: \mathbb{R} \to \mathbb{R}$ la fonction définie par:
$$ f(x) = \begin{cases} \frac{\sin(x^2)}{x} & \text{si } x \neq 0 \\ 0 & \text{si } x = 0 \end{cases} $$

1.  Démontrer, en utilisant la définition $\epsilon-\delta$ de la continuité, que $f$ est continue en $x=0$.
2.  Démontrer que $f$ est continue sur $\mathbb{R}^* = \mathbb{R} \setminus \{0\}$.
3.  En déduire que $f$ est continue sur $\mathbb{R}$.

\subsection{Partie 2 : Applications des Théorèmes Fondamentaux}

1.  Soit $h: [0, 1] \to [0, 1]$ une fonction continue. Démontrer que $h$ admet au moins un point fixe, c'est-à-dire qu'il existe $c \in [0, 1]$ tel que $h(c) = c$.
2.  Soit $P(x) = a_n x^n + a_{n-1} x^{n-1} + \dots + a_1 x + a_0$ un polynôme à coefficients réels, où $n$ est un entier impair et $a_n \neq 0$. Démontrer que $P(x)$ admet au moins une racine réelle.

\subsection{Partie 3 : Continuité Uniforme}

1.  Rappeler la définition de la continuité uniforme d'une fonction $f: I \to \mathbb{R}$ sur un intervalle $I$.
2.  Démontrer que la fonction $f(x) = x\textasciicircum 2$ n'est pas uniformément continue sur $\mathbb{R}$.
3.  Démontrer que la fonction $g(x) = \sqrt{x}$ est uniformément continue sur $[0, +\infty)$.

---

\section{Correction de l'Exercice 4}

\subsection{Partie 1 : Analyse de la continuité ponctuelle et globale}

1.  \textbf{Démonstration de la continuité de $f$ en $x=0$ par la définition $\epsilon-\delta$ :}
    La fonction $f$ est continue en $x=0$ si et seulement si pour tout $\epsilon > 0$, il existe un $\delta > 0$ tel que pour tout $x \in \mathbb{R}$, si $|x - 0| < \delta$, alors $|f(x) - f(0)| < \epsilon$.
    Nous avons $f(0) = 0$. Pour $x \neq 0$, $f(x) = \frac{\sin(x^2)}{x}$.
    Nous devons donc montrer que pour tout $\epsilon > 0$, il existe $\delta > 0$ tel que si $0 < |x| < \delta$, alors $\left| \frac{\sin(x^2)}{x} - 0 \right| < \epsilon$.
    Considérons l'expression $\left| \frac{\sin(x^2)}{x} \right|$ pour $x \neq 0$.
    Nous pouvons la réécrire comme $\left| \frac{\sin(x^2)}{x^2} \cdot x \right| = \left| \frac{\sin(x^2)}{x^2} \right| \cdot |x|$.
    Nous utilisons l'inégalité fondamentale $|\sin(u)| \le |u|$ pour tout $u \in \mathbb{R}$.
    Si $u \neq 0$, alors $\left| \frac{\sin(u)}{u} \right| \le 1$.
    En particulier, pour $x \neq 0$, $x^2 \neq 0$, donc nous avons $\left| \frac{\sin(x^2)}{x^2} \right| \le 1$.
    Par conséquent, pour tout $x \neq 0$, nous avons $|f(x) - f(0)| = \left| \frac{\sin(x^2)}{x^2} \right| \cdot |x| \le 1 \cdot |x| = |x|$.
    Soit $\epsilon > 0$ donné. Nous voulons que $|f(x) - f(0)| < \epsilon$.
    Puisque nous avons montré que $|f(x) - f(0)| \le |x|$, il suffit de choisir $\delta$ tel que $|x| < \delta$ implique $|x| < \epsilon$.
    Nous pouvons donc choisir $\delta = \epsilon$.
    Vérification : Soit $\epsilon > 0$. Choisissons $\delta = \epsilon$.
    Si $x \in \mathbb{R}$ et $0 < |x| < \delta$, alors $0 < |x| < \epsilon$.
    Pour ces valeurs de $x$, nous avons $|f(x) - f(0)| = \left| \frac{\sin(x^2)}{x} \right| = \left| \frac{\sin(x^2)}{x^2} \right| \cdot |x|$.
    Puisque $x \neq 0$, $x^2 \neq 0$, et nous savons que $\left| \frac{\sin(u)}{u} \right| \le 1$ pour $u \neq 0$.
    Donc, $\left| \frac{\sin(x^2)}{x^2} \right| \le 1$.
    Par conséquent, $|f(x) - f(0)| \le 1 \cdot |x| = |x|$.
    Comme $|x| < \delta = \epsilon$, nous avons $|f(x) - f(0)| < \epsilon$.
    La fonction $f$ est donc continue en $x=0$.

2.  \textbf{Démonstration de la continuité de $f$ sur $\mathbb{R}^*$ :}
    Pour tout $x \in \mathbb{R}^*$, la fonction $f(x)$ est définie par $f(x) = \frac{\sin(x^2)}{x}$.
    Considérons les fonctions suivantes :
    *   La fonction $g_1(x) = x^2$. C'est une fonction polynomiale, donc elle est continue sur $\mathbb{R}$. En particulier, elle est continue sur $\mathbb{R}^*$.
    *   La fonction $g_2(u) = \sin(u)$. C'est une fonction trigonométrique élémentaire, donc elle est continue sur $\mathbb{R}$.
    La fonction $g_3(x) = \sin(x^2)$ est la composition de $g_2$ et $g_1$, c'est-à-dire $g_3(x) = g_2(g_1(x))$. Puisque $g_1$ est continue sur $\mathbb{R}$ et $g_2$ est continue sur $\mathbb{R}$, leur composition $g_3$ est continue sur $\mathbb{R}$. En particulier, $g_3$ est continue sur $\mathbb{R}^*$.
    *   La fonction $g_4(x) = x$. C'est une fonction polynomiale, donc elle est continue sur $\mathbb{R}$. En particulier, elle est continue sur $\mathbb{R}^*$.
    De plus, pour tout $x \in \mathbb{R}^*$, $g_4(x) = x \neq 0$.
    La fonction $f(x)$ est le quotient de $g_3(x)$ et $g_4(x)$, c'est-à-dire $f(x) = \frac{g_3(x)}{g_4(x)}$.
    Le théorème sur la continuité des quotients de fonctions stipule que si $A(x)$ et $B(x)$ sont continues sur un intervalle $I$ et $B(x) \neq 0$ pour tout $x \in I$, alors $\frac{A(x)}{B(x)}$ est continue sur $I$.
    Ici, $A(x) = \sin(x^2)$ est continue sur $\mathbb{R}^*$ et $B(x) = x$ est continue sur $\mathbb{R}^*$ et non nulle sur $\mathbb{R}^*$.
    Par conséquent, $f(x) = \frac{\sin(x^2)}{x}$ est continue sur $\mathbb{R}^*$.

3.  \textbf{Conclusion sur la continuité de $f$ sur $\mathbb{R}$ :}
    D'après la question 1, nous avons démontré que $f$ est continue au point $x=0$.
    D'après la question 2, nous avons démontré que $f$ est continue sur l'ensemble $\mathbb{R}^* = \mathbb{R} \setminus \{0\}$.
    L'ensemble $\mathbb{R}$ est l'union de $\mathbb{R}^*$ et du point $\{0\}$, c'est-à-dire $\mathbb{R} = \mathbb{R}^* \cup \{0\}$.
    Puisque $f$ est continue en chaque point de $\mathbb{R}^*$ et également au point $0$, nous pouvons conclure que $f$ est continue sur tout $\mathbb{R}$.

\subsection{Partie 2 : Applications des Théorèmes Fondamentaux}

1.  \textbf{Démonstration de l'existence d'un point fixe pour $h: [0, 1] \to [0, 1]$ continue :}
    Soit $h: [0, 1] \to [0, 1]$ une fonction continue. Nous cherchons à montrer qu'il existe $c \in [0, 1]$ tel que $h(c) = c$.
    Considérons la fonction auxiliaire $g: [0, 1] \to \mathbb{R}$ définie par $g(x) = h(x) - x$.
    *   \textbf{Continuité de $g$ :} La fonction $h$ est continue sur $[0, 1]$ par hypothèse. La fonction $k(x) = x$ est une fonction polynomiale, donc elle est continue sur $\mathbb{R}$, et par conséquent continue sur $[0, 1]$. La différence de deux fonctions continues est continue, donc $g(x) = h(x) - x$ est continue sur l'intervalle fermé et borné $[0, 1]$.
    *   \textbf{Évaluation aux bornes de l'intervalle :}
        *   En $x=0$: $g(0) = h(0) - 0 = h(0)$. Puisque le codomaine de $h$ est $[0, 1]$, nous savons que $h(0) \in [0, 1]$. Par conséquent, $h(0) \ge 0$, ce qui implique $g(0) \ge 0$.
        *   En $x=1$: $g(1) = h(1) - 1$. Puisque le codomaine de $h$ est $[0, 1]$, nous savons que $h(1) \in [0, 1]$. Par conséquent, $h(1) \le 1$, ce qui implique $h(1) - 1 \le 0$, donc $g(1) \le 0$.
    *   \textbf{Application du Théorème des Valeurs Intermédiaires (TVI) :}
        Nous avons trois cas possibles :
        1.  Si $g(0) = 0$, alors $h(0) - 0 = 0$, ce qui signifie $h(0) = 0$. Dans ce cas, $c=0$ est un point fixe.
        2.  Si $g(1) = 0$, alors $h(1) - 1 = 0$, ce qui signifie $h(1) = 1$. Dans ce cas, $c=1$ est un point fixe.
        3.  Si $g(0) > 0$ et $g(1) < 0$. Dans ce cas, $g(0)$ et $g(1)$ ont des signes opposés. Puisque $g$ est continue sur l'intervalle fermé $[0, 1]$, et que $0$ est une valeur comprise entre $g(1)$ et $g(0)$, le Théorème des Valeurs Intermédiaires garantit l'existence d'au moins un $c \in (0, 1)$ tel que $g(c) = 0$.
            Le TVI stipule que si une fonction $f$ est continue sur un intervalle fermé $[a, b]$, alors pour toute valeur $y_0$ entre $f(a)$ et $f(b)$ (inclusivement), il existe au moins un $c \in [a, b]$ tel que $f(c) = y_0$.
            Ici, $a=0$, $b=1$, et $y_0=0$ est entre $g(1)$ et $g(0)$.
            Donc, il existe $c \in (0, 1)$ tel que $g(c) = 0$.
            Cela signifie $h(c) - c = 0$, d'où $h(c) = c$.
    Dans tous les cas, il existe au moins un $c \in [0, 1]$ tel que $h(c) = c$. Ce $c$ est un point fixe de $h$.

2.  \textbf{Démonstration de l'existence d'une racine réelle pour un polynôme de degré impair :}
    Soit $P(x) = a_n x^n + a_{n-1} x^{n-1} + \dots + a_1 x + a_0$ un polynôme à coefficients réels, où $n$ est un entier impair et $a_n \neq 0$. Nous voulons démontrer que $P(x)$ admet au moins une racine réelle.
    *   \textbf{Continuité de $P$ :} Toute fonction polynomiale est continue sur $\mathbb{R}$. Par conséquent, $P(x)$ est continue sur $\mathbb{R}$.
    *   \textbf{Comportement aux limites :} Nous allons examiner les limites de $P(x)$ lorsque $x \to +\infty$ et $x \to -\infty$. Le comportement d'un polynôme pour de grandes valeurs absolues de $x$ est dominé par son terme de plus haut degré, $a_n x^n$.
        Nous pouvons factoriser $a_n x^n$:
        $$ P(x) = a_n x^n \left( 1 + \frac{a_{n-1}}{a_n x} + \frac{a_{n-2}}{a_n x^2} + \dots + \frac{a_0}{a_n x^n} \right) $$
        Lorsque $x \to +\infty$, tous les termes $\frac{a_{k}}{a_n x^{n-k}}$ pour $k < n$ tendent vers $0$.
        Donc, $\lim_{x \to +\infty} \left( 1 + \frac{a_{n-1}}{a_n x} + \dots + \frac{a_0}{a_n x^n} \right) = 1$.
        Par conséquent, $\lim_{x \to +\infty} P(x) = \lim_{x \to +\infty} (a_n x^n \cdot 1) = \lim_{x \to +\infty} a_n x^n$.
        Puisque $n$ est un entier impair :
        *   Si $a_n > 0$, alors $\lim_{x \to +\infty} x^n = +\infty$, donc $\lim_{x \to +\infty} P(x) = +\infty$.
        *   Si $a_n < 0$, alors $\lim_{x \to +\infty} x^n = +\infty$, donc $\lim_{x \to +\infty} P(x) = -\infty$.

        De même, lorsque $x \to -\infty$:
        $\lim_{x \to -\infty} P(x) = \lim_{x \to -\infty} a_n x^n$.
        Puisque $n$ est un entier impair, $x^n$ a le même signe que $x$. Donc, $\lim_{x \to -\infty} x^n = -\infty$.
        *   Si $a_n > 0$, alors $\lim_{x \to -\infty} P(x) = -\infty$.
        *   Si $a_n < 0$, alors $\lim_{x \to -\infty} P(x) = +\infty$.

    *   \textbf{Synthèse des limites et application du TVI :}
        Nous observons que dans tous les cas, les limites de $P(x)$ en $+\infty$ et en $-\infty$ sont de signes opposés :
        *   Si $a_n > 0$: $\lim_{x \to -\infty} P(x) = -\infty$ et $\lim_{x \to +\infty} P(x) = +\infty$.
        *   Si $a_n < 0$: $\lim_{x \to -\infty} P(x) = +\infty$ et $\lim_{x \to +\infty} P(x) = -\infty$.

        Cela signifie qu'il existe des valeurs $x_1$ et $x_2$ telles que $P(x_1)$ et $P(x_2)$ ont des signes opposés.
        Plus précisément :
        *   Si $a_n > 0$: Puisque $\lim_{x \to -\infty} P(x) = -\infty$, il existe un $x_1 \in \mathbb{R}$ (suffisamment petit) tel que $P(x_1) < 0$. Puisque $\lim_{x \to +\infty} P(x) = +\infty$, il existe un $x_2 \in \mathbb{R}$ (suffisamment grand) tel que $P(x_2) > 0$. Nous pouvons choisir $x_1 < x_2$.
        *   Si $a_n < 0$: Puisque $\lim_{x \to -\infty} P(x) = +\infty$, il existe un $x_1 \in \mathbb{R}$ (suffisamment petit) tel que $P(x_1) > 0$. Puisque $\lim_{x \to +\infty} P(x) = -\infty$, il existe un $x_2 \in \mathbb{R}$ (suffisamment grand) tel que $P(x_2) < 0$. Nous pouvons choisir $x_1 < x_2$.

        Dans les deux cas, nous avons trouvé un intervalle fermé et borné $[x_1, x_2]$ sur lequel $P(x)$ est continue (car $P$ est continue sur $\mathbb{R}$), et les valeurs $P(x_1)$ et $P(x_2)$ ont des signes opposés.
        Par le Théorème des Valeurs Intermédiaires (TVI), puisque $0$ est une valeur comprise entre $P(x_1)$ et $P(x_2)$, il existe au moins un $c \in (x_1, x_2)$ tel que $P(c) = 0$.
        Ce $c$ est une racine réelle du polynôme $P(x)$.
        Par conséquent, tout polynôme de degré impair à coefficients réels admet au moins une racine réelle.

\subsection{Partie 3 : Continuité Uniforme}

1.  \textbf{Définition de la continuité uniforme :}
    Une fonction $f: I \to \mathbb{R}$ est dite uniformément continue sur un intervalle $I$ si pour tout $\epsilon > 0$, il existe un $\delta > 0$ (qui ne dépend que de $\epsilon$ et de $f$, mais pas des points $x$ et $y$ spécifiques dans $I$) tel que pour tous $x, y \in I$, si $|x - y| < \delta$, alors $|f(x) - f(y)| < \epsilon$.

2.  \textbf{Démonstration que $f(x) = x\textasciicircum 2$ n'est pas uniformément continue sur $\mathbb{R}$ :}
    Pour montrer que $f(x) = x\textasciicircum 2$ n'est pas uniformément continue sur $\mathbb{R}$, nous devons nier la définition de la continuité uniforme. Cela signifie qu'il existe un $\epsilon_0 > 0$ tel que pour tout $\delta > 0$, il existe des $x, y \in \mathbb{R}$ tels que $|x - y| < \delta$ mais $|f(x) - f(y)| \ge \epsilon_0$.
    Choisissons $\epsilon_0 = 1$.
    Soit $\delta > 0$ un nombre réel arbitraire. Nous devons trouver $x, y \in \mathbb{R}$ tels que $|x - y| < \delta$ et $|x^2 - y^2| \ge 1$.
    Considérons l'expression $|x^2 - y^2| = |(x - y)(x + y)| = |x - y| |x + y|$.
    Pour satisfaire la condition $|x - y| < \delta$, nous pouvons choisir $y = x + \frac{\delta}{2}$.
    Alors $|x - y| = \left| x - \left(x + \frac{\delta}{2}\right) \right| = \left| -\frac{\delta}{2} \right| = \frac{\delta}{2}$.
    Puisque $\delta > 0$, nous avons $\frac{\delta}{2} < \delta$, donc la première condition est satisfaite.
    Maintenant, substituons ces valeurs dans l'expression de $|f(x) - f(y)|$:
    $|f(x) - f(y)| = \left| \frac{\delta}{2} \right| \left| x + \left(x + \frac{\delta}{2}\right) \right| = \frac{\delta}{2} \left| 2x + \frac{\delta}{2} \right|$.
    Nous voulons que cette expression soit supérieure ou égale à $\epsilon_0 = 1$.
    C'est-à-dire, nous voulons $\frac{\delta}{2} \left| 2x + \frac{\delta}{2} \right| \ge 1$.
    Cela est équivalent à $\left| 2x + \frac{\delta}{2} \right| \ge \frac{2}{\delta}$.
    Nous pouvons choisir $x$ de manière à satisfaire cette inégalité. Par exemple, choisissons $2x + \frac{\delta}{2} = \frac{2}{\delta}$.
    Alors $2x = \frac{2}{\delta} - \frac{\delta}{2}$, ce qui donne $x = \frac{1}{\delta} - \frac{\delta}{4}$.
    Avec ce choix de $x$, $y = x + \frac{\delta}{2} = \left( \frac{1}{\delta} - \frac{\delta}{4} \right) + \frac{\delta}{2} = \frac{1}{\delta} + \frac{\delta}{4}$.
    Pour tout $\delta > 0$ donné, nous avons trouvé $x = \frac{1}{\delta} - \frac{\delta}{4}$ et $y = \frac{1}{\delta} + \frac{\delta}{4}$.
    Vérifions les conditions :
    1.  $|x - y| = \left| \left(\frac{1}{\delta} - \frac{\delta}{4}\right) - \left(\frac{1}{\delta} + \frac{\delta}{4}\right) \right| = \left| -\frac{\delta}{2} \right| = \frac{\delta}{2}$. Puisque $\frac{\delta}{2} < \delta$, la première condition est satisfaite.
    2.  $|f(x) - f(y)| = |x^2 - y^2| = |(x-y)(x+y)| = \left| -\frac{\delta}{2} \right| \left| \left(\frac{1}{\delta} - \frac{\delta}{4}\right) + \left(\frac{1}{\delta} + \frac{\delta}{4}\right) \right|$
        $= \frac{\delta}{2} \left| \frac{2}{\delta} \right| = \frac{\delta}{2} \cdot \frac{2}{\delta} = 1$.
    Nous avons donc trouvé, pour tout $\delta > 0$, des points $x, y \in \mathbb{R}$ tels que $|x - y| < \delta$ mais $|f(x) - f(y)| = 1$.
    Puisque $1 \ge \epsilon_0 = 1$, la condition de non-uniforme continuité est satisfaite.
    Par conséquent, la fonction $f(x) = x\textasciicircum 2$ n'est pas uniformément continue sur $\mathbb{R}$.

3.  \textbf{Démonstration que $g(x) = \sqrt{x}$ est uniformément continue sur $[0, +\infty)$ :}
    Nous voulons montrer que pour tout $\epsilon > 0$, il existe un $\delta > 0$ tel que pour tous $x, y \in [0, +\infty)$, si $|x - y| < \delta$, alors $|\sqrt{x} - \sqrt{y}| < \epsilon$.
    Soit $\epsilon > 0$ donné.
    Considérons l'expression $|\sqrt{x} - \sqrt{y}|$.
    Nous allons d'abord établir une inégalité utile : $|\sqrt{x} - \sqrt{y}| \le \sqrt{|x - y|}$ pour tous $x, y \ge 0$.
    \textbf{Preuve de l'inégalité $|\sqrt{x} - \sqrt{y}| \le \sqrt{|x - y|}$ :}
    Sans perte de généralité, supposons $x \ge y$. Alors $|x - y| = x - y$ et $|\sqrt{x} - \sqrt{y}| = \sqrt{x} - \sqrt{y}$ (puisque $\sqrt{x} \ge \sqrt{y}$ car la fonction racine carrée est croissante).
    Nous voulons montrer $\sqrt{x} - \sqrt{y} \le \sqrt{x - y}$.
    Puisque les deux membres de l'inégalité sont non-négatifs (car $x \ge y \ge 0$), nous pouvons élever les deux côtés au carré sans changer le sens de l'inégalité :
    $(\sqrt{x} - \sqrt{y})^2 \le (\sqrt{x - y})^2$
    $x - 2\sqrt{xy} + y \le x - y$
    Soustraire $x$ des deux côtés :
    $-2\sqrt{xy} + y \le -y$
    Ajouter $y$ aux deux côtés :
    $-2\sqrt{xy} \le -2y$
    Diviser par $-2$ et inverser le sens de l'inégalité :
    $\sqrt{xy} \ge y$
    Si $y=0$, l'inégalité devient $\sqrt{x \cdot 0} \ge 0$, soit $0 \ge 0$, ce qui est vrai.
    Si $y > 0$, nous pouvons diviser par $\sqrt{y}$ (qui est positif) :
    $\sqrt{x} \ge \sqrt{y}$
    Cette dernière inégalité est vraie car nous avons supposé $x \ge y$ et la fonction racine carrée est croissante.
    Puisque toutes les étapes sont réversibles (en respectant les conditions de non-négativité pour le carré et de non-nullité pour la division), l'inégalité originale $|\sqrt{x} - \sqrt{y}| \le \sqrt{|x - y|}$ est vraie pour $x \ge y$.
    Par symétrie, si $y \ge x$, alors $|\sqrt{y} - \sqrt{x}| \le \sqrt{|y - x|}$, ce qui est équivalent à $|\sqrt{x} - \sqrt{y}| \le \sqrt{|x - y|}$.
    L'inégalité $|\sqrt{x} - \sqrt{y}| \le \sqrt{|x - y|}$ est donc établie pour tous $x, y \in [0, +\infty)$.

    \textbf{Application de l'inégalité pour la continuité uniforme :}
    Nous voulons que $|\sqrt{x} - \sqrt{y}| < \epsilon$.
    En utilisant l'inégalité que nous venons de prouver, si nous assurons que $\sqrt{|x - y|} < \epsilon$, alors nous aurons $|\sqrt{x} - \sqrt{y}| < \epsilon$.
    Pour que $\sqrt{|x - y|} < \epsilon$, il suffit que $|x - y| < \epsilon^2$.
    Nous pouvons donc choisir $\delta = \epsilon^2$.

    \textbf{Vérification finale :}
    Soit $\epsilon > 0$ donné.
    Choisissons $\delta = \epsilon^2$.
    Soient $x, y \in [0, +\infty)$ tels que $|x - y| < \delta$.
    Nous voulons montrer que $|\sqrt{x} - \sqrt{y}| < \epsilon$.
    D'après l'inégalité établie, nous avons $|\sqrt{x} - \sqrt{y}| \le \sqrt{|x - y|}$.
    Puisque $|x - y| < \delta$, nous avons $\sqrt{|x - y|} < \sqrt{\delta}$.
    En substituant $\delta = \epsilon^2$, nous obtenons $\sqrt{\delta} = \sqrt{\epsilon^2} = \epsilon$ (car $\epsilon > 0$).
    Par conséquent, $|\sqrt{x} - \sqrt{y}| \le \sqrt{|x - y|} < \sqrt{\delta} = \epsilon$.
    Ainsi, pour tout $\epsilon > 0$, nous avons trouvé un $\delta = \epsilon^2 > 0$ tel que pour tous $x, y \in [0, +\infty)$, si $|x - y| < \delta$, alors $|\sqrt{x} - \sqrt{y}| < \epsilon$.
    Ceci démontre que la fonction $g(x) = \sqrt{x}$ est uniformément continue sur $[0, +\infty)$.

\section{Exo-05}
\section{Exercice 5 - Exploration Approfondie de la Continuité et de ses Applications}

\subsection{Énoncé de l'exercice}

\subsubsection{Partie A: Maîtrise de la Définition Epsilon-Delta}

Soit la fonction $f: \mathbb{R}^* \to \mathbb{R}$ définie par $f(x) = \frac{\sqrt{1+x^2} - 1}{x^2}$.

1.  Montrer que $\lim_{x \to 0} f(x)$ existe et déterminer sa valeur $L$.
2.  En utilisant la définition $\epsilon-\delta$ de la continuité, montrer que la fonction $g: \mathbb{R} \to \mathbb{R}$ définie par $g(x) = f(x)$ pour $x \neq 0$ et $g(0) = L$ est continue en $x=0$.

\subsubsection{Partie B: Continuité des Fonctions Définies par Morceaux et Paramètres}

Soit la fonction $h: \mathbb{R} \to \mathbb{R}$ définie par:
$$ h(x) = \begin{cases} \frac{\sin(ax)}{x} & \text{si } x < 0 \\ b & \text{si } x = 0 \\ \frac{e^{2x} - 1}{x} & \text{si } x > 0 \end{cases} $$
où $a$ et $b$ sont des constantes réelles.

1.  Déterminer les valeurs de $a$ et $b$ pour lesquelles $h$ est continue en $x=0$.
2.  Pour les valeurs de $a$ et $b$ trouvées, montrer que $h$ est continue sur $\mathbb{R}$.

\subsubsection{Partie C: Application du Théorème des Valeurs Intermédiaires (TVI) et Propriétés Globales}

1.  Soit $P(x)$ un polynôme de degré impair. Montrer que $P(x)$ admet au moins une racine réelle.
2.  Soit $f: [0, 1] \to [0, 1]$ une fonction continue. Montrer qu'il existe au moins un point fixe $c \in [0, 1]$ tel que $f(c) = c$.
3.  Soit $f: \mathbb{R} \to \mathbb{R}$ une fonction continue telle que $\lim_{x \to -\infty} f(x) = -\infty$ et $\lim_{x \to +\infty} f(x) = +\infty$. Montrer que pour tout $y \in \mathbb{R}$, il existe $x \in \mathbb{R}$ tel que $f(x) = y$.

\subsubsection{Partie D: Un Défi Théorique (Équation Fonctionnelle de Cauchy)}

Soit $f: \mathbb{R} \to \mathbb{R}$ une fonction continue telle que $f(x+y) = f(x) + f(y)$ pour tout $x, y \in \mathbb{R}$.

1.  Montrer que $f(nx) = nf(x)$ pour tout $n \in \mathbb{Z}$ et $x \in \mathbb{R}$.
2.  Montrer que $f(qx) = qf(x)$ pour tout $q \in \mathbb{Q}$ et $x \in \mathbb{R}$.
3.  En utilisant la continuité de $f$, montrer que $f(x) = cx$ pour une certaine constante $c \in \mathbb{R}$ et pour tout $x \in \mathbb{R}$.

---

\subsection{Correction de l'exercice}

\subsubsection{Partie A: Maîtrise de la Définition Epsilon-Delta}

1.  \textbf{Montrer que $\lim_{x \to 0} f(x)$ existe et déterminer sa valeur $L$.}

    Nous devons évaluer la limite $\lim_{x \to 0} \frac{\sqrt{1+x^2} - 1}{x^2}$.
    Lorsque $x \to 0$, le numérateur tend vers $\sqrt{1+0^2} - 1 = \sqrt{1} - 1 = 0$, et le dénominateur tend vers $0^2 = 0$. Il s'agit d'une forme indéterminée de type $\frac{0}{0}$.
    Pour lever l'indétermination, nous allons multiplier le numérateur et le dénominateur par l'expression conjuguée du numérateur, qui est $\sqrt{1+x^2} + 1$.

    $$ f(x) = \frac{\sqrt{1+x^2} - 1}{x^2} = \frac{(\sqrt{1+x^2} - 1)(\sqrt{1+x^2} + 1)}{x^2(\sqrt{1+x^2} + 1)} $$
    En utilisant l'identité $(A-B)(A+B) = A^2 - B^2$ pour le numérateur, avec $A = \sqrt{1+x^2}$ et $B = 1$:
    $$ f(x) = \frac{(\sqrt{1+x^2})^2 - 1^2}{x^2(\sqrt{1+x^2} + 1)} = \frac{(1+x^2) - 1}{x^2(\sqrt{1+x^2} + 1)} $$
    $$ f(x) = \frac{x^2}{x^2(\sqrt{1+x^2} + 1)} $$
    Puisque nous calculons la limite lorsque $x \to 0$, nous considérons des valeurs de $x$ qui sont proches de 0 mais non nulles. Par conséquent, $x^2 \neq 0$, et nous pouvons simplifier l'expression en divisant le numérateur et le dénominateur par $x^2$:
    $$ f(x) = \frac{1}{\sqrt{1+x^2} + 1} \quad \text{pour } x \neq 0 $$
    Maintenant, nous pouvons évaluer la limite par substitution directe, car le dénominateur ne tend pas vers zéro:
    $$ \lim_{x \to 0} f(x) = \lim_{x \to 0} \frac{1}{\sqrt{1+x^2} + 1} $$
    La fonction $x \mapsto x^2$ est continue en $x=0$, donc $\lim_{x \to 0} x^2 = 0$.
    La fonction $u \mapsto 1+u$ est continue, donc $\lim_{x \to 0} (1+x^2) = 1 + \lim_{x \to 0} x^2 = 1+0 = 1$.
    La fonction $v \mapsto \sqrt{v}$ est continue pour $v > 0$. Puisque $1+x^2 \to 1 > 0$, $\lim_{x \to 0} \sqrt{1+x^2} = \sqrt{\lim_{x \to 0} (1+x^2)} = \sqrt{1} = 1$.
    Par conséquent, $\lim_{x \to 0} (\sqrt{1+x^2} + 1) = \lim_{x \to 0} \sqrt{1+x^2} + \lim_{x \to 0} 1 = 1+1 = 2$.
    Enfin, la fonction $w \mapsto \frac{1}{w}$ est continue pour $w \neq 0$. Puisque le dénominateur tend vers $2 \neq 0$:
    $$ \lim_{x \to 0} f(x) = \frac{1}{2} $$
    La limite existe et sa valeur est $L = \frac{1}{2}$.

2.  \textbf{En utilisant la définition $\epsilon-\delta$ de la continuité, montrer que la fonction $g: \mathbb{R} \to \mathbb{R}$ définie par $g(x) = f(x)$ pour $x \neq 0$ et $g(0) = L$ est continue en $x=0$.}

    La fonction $g(x)$ est définie par:
    $$ g(x) = \begin{cases} \frac{1}{\sqrt{1+x^2} + 1} & \text{si } x \neq 0 \\ \frac{1}{2} & \text{si } x = 0 \end{cases} $$
    Pour montrer que $g$ est continue en $x=0$ en utilisant la définition $\epsilon-\delta$, nous devons montrer que pour tout $\epsilon > 0$, il existe un $\delta > 0$ tel que si $|x - 0| < \delta$, alors $|g(x) - g(0)| < \epsilon$.

    Soit $\epsilon > 0$ donné.
    Nous voulons trouver $\delta > 0$ tel que pour tout $x \in \mathbb{R}$ avec $|x| < \delta$, nous ayons $|g(x) - g(0)| < \epsilon$.

    Considérons deux cas:
    *   \textbf{Cas 1: $x = 0$}.
        Alors $|g(0) - g(0)| = |L - L| = 0$. Puisque $0 < \epsilon$ pour tout $\epsilon > 0$, la condition est satisfaite pour $x=0$.

    *   \textbf{Cas 2: $x \neq 0$}.
        Alors $g(x) = \frac{1}{\sqrt{1+x^2} + 1}$ et $g(0) = \frac{1}{2}$.
        Nous devons majorer l'expression $|g(x) - g(0)|$:
        $$ |g(x) - g(0)| = \left| \frac{1}{\sqrt{1+x^2} + 1} - \frac{1}{2} \right| $$
        Mettons les fractions sur un dénominateur commun:
        $$ = \left| \frac{2 - (\sqrt{1+x^2} + 1)}{2(\sqrt{1+x^2} + 1)} \right| = \left| \frac{1 - \sqrt{1+x^2}}{2(\sqrt{1+x^2} + 1)} \right| $$
        Pour simplifier le numérateur, multiplions-le par son conjugué $1 + \sqrt{1+x^2}$. Pour maintenir l'égalité, nous devons aussi multiplier le dénominateur par cette même expression:
        $$ = \left| \frac{(1 - \sqrt{1+x^2})(1 + \sqrt{1+x^2})}{2(\sqrt{1+x^2} + 1)(1 + \sqrt{1+x^2})} \right| $$
        En utilisant l'identité $(A-B)(A+B) = A^2 - B^2$ pour le numérateur, avec $A=1$ et $B=\sqrt{1+x^2}$:
        $$ = \left| \frac{1^2 - (\sqrt{1+x^2})^2}{2(\sqrt{1+x^2} + 1)^2} \right| = \left| \frac{1 - (1+x^2)}{2(\sqrt{1+x^2} + 1)^2} \right| $$
        $$ = \left| \frac{-x^2}{2(\sqrt{1+x^2} + 1)^2} \right| = \frac{|-x^2|}{2(\sqrt{1+x^2} + 1)^2} = \frac{x^2}{2(\sqrt{1+x^2} + 1)^2} $$
        Maintenant, nous devons trouver une borne inférieure pour le dénominateur afin de majorer l'expression.
        Puisque $x^2 \ge 0$, nous avons $1+x^2 \ge 1$.
        En prenant la racine carrée (la fonction racine carrée est croissante), $\sqrt{1+x^2} \ge \sqrt{1} = 1$.
        En ajoutant 1, $\sqrt{1+x^2} + 1 \ge 1+1 = 2$.
        En élevant au carré (les deux membres sont positifs), $(\sqrt{1+x^2} + 1)^2 \ge 2^2 = 4$.
        En multipliant par 2, $2(\sqrt{1+x^2} + 1)^2 \ge 2 \times 4 = 8$.
        Par conséquent, $\frac{1}{2(\sqrt{1+x^2} + 1)^2} \le \frac{1}{8}$.
        Ainsi, nous pouvons majorer $|g(x) - g(0)|$:
        $$ |g(x) - g(0)| = \frac{x^2}{2(\sqrt{1+x^2} + 1)^2} \le \frac{x^2}{8} $$
        Nous voulons que cette expression soit inférieure à $\epsilon$:
        $$ \frac{x^2}{8} < \epsilon $$
        $$ x^2 < 8\epsilon $$
        $$ |x| < \sqrt{8\epsilon} $$
        Nous pouvons donc choisir $\delta = \sqrt{8\epsilon}$.

    \textbf{Conclusion pour la définition $\epsilon-\delta$}:
    Pour tout $\epsilon > 0$, choisissons $\delta = \sqrt{8\epsilon}$.
    Si $|x - 0| < \delta$, c'est-à-dire $|x| < \sqrt{8\epsilon}$:
    *   Si $x=0$, alors $|g(0) - g(0)| = 0 < \epsilon$.
    *   Si $x \neq 0$, alors $|g(x) - g(0)| = \frac{x^2}{2(\sqrt{1+x^2} + 1)^2}$.
        Comme nous l'avons montré, $2(\sqrt{1+x^2} + 1)^2 \ge 8$, donc $\frac{1}{2(\sqrt{1+x^2} + 1)^2} \le \frac{1}{8}$.
        Ainsi, $|g(x) - g(0)| \le \frac{x^2}{8}$.
        Puisque $|x| < \sqrt{8\epsilon}$, nous avons $x^2 < 8\epsilon$.
        Par conséquent, $|g(x) - g(0)| \le \frac{x^2}{8} < \frac{8\epsilon}{8} = \epsilon$.
    Dans tous les cas, si $|x| < \delta$, alors $|g(x) - g(0)| < \epsilon$.
    La fonction $g$ est donc continue en $x=0$.

\subsubsection{Partie B: Continuité des Fonctions Définies par Morceaux et Paramètres}

1.  \textbf{Déterminer les valeurs de $a$ et $b$ pour lesquelles $h$ est continue en $x=0$.}

    Pour que la fonction $h$ soit continue en $x=0$, il faut que la limite de $h(x)$ lorsque $x$ tend vers $0$ existe et soit égale à $h(0)$. Cela implique que les limites à gauche et à droite en $x=0$ doivent exister et être égales à $h(0)$.
    $$ \lim_{x \to 0^-} h(x) = \lim_{x \to 0^+} h(x) = h(0) $$

    *   \textbf{Valeur de $h(0)$}:
        Par définition, $h(0) = b$.

    *   \textbf{Limite à gauche en $x=0$}:
        Pour $x < 0$, $h(x) = \frac{\sin(ax)}{x}$.
        $$ \lim_{x \to 0^-} h(x) = \lim_{x \to 0^-} \frac{\sin(ax)}{x} $$
        Si $a=0$, alors $\lim_{x \to 0^-} \frac{\sin(0)}{x} = \lim_{x \to 0^-} \frac{0}{x} = 0$.
        Si $a \neq 0$, nous utilisons la limite fondamentale $\lim_{u \to 0} \frac{\sin u}{u} = 1$.
        Posons $u = ax$. Lorsque $x \to 0^-$, $u \to 0$.
        $$ \lim_{x \to 0^-} \frac{\sin(ax)}{x} = \lim_{x \to 0^-} a \cdot \frac{\sin(ax)}{ax} $$
        Par la propriété des limites, $\lim_{x \to 0^-} a \cdot \frac{\sin(ax)}{ax} = a \cdot \lim_{u \to 0} \frac{\sin u}{u} = a \cdot 1 = a$.
        Donc, $\lim_{x \to 0^-} h(x) = a$.

    *   \textbf{Limite à droite en $x=0$}:
        Pour $x > 0$, $h(x) = \frac{e^{2x} - 1}{x}$.
        $$ \lim_{x \to 0^+} h(x) = \lim_{x \to 0^+} \frac{e^{2x} - 1}{x} $$
        Nous utilisons la limite fondamentale $\lim_{u \to 0} \frac{e^u - 1}{u} = 1$.
        Posons $u = 2x$. Lorsque $x \to 0^+$, $u \to 0$.
        $$ \lim_{x \to 0^+} \frac{e^{2x} - 1}{x} = \lim_{x \to 0^+} 2 \cdot \frac{e^{2x} - 1}{2x} $$
        Par la propriété des limites, $\lim_{x \to 0^+} 2 \cdot \frac{e^{2x} - 1}{2x} = 2 \cdot \lim_{u \to 0} \frac{e^u - 1}{u} = 2 \cdot 1 = 2$.
        Donc, $\lim_{x \to 0^+} h(x) = 2$.

    Pour que $h$ soit continue en $x=0$, nous devons avoir $a = b = 2$.

2.  \textbf{Pour les valeurs de $a$ et $b$ trouvées, montrer que $h$ est continue sur $\mathbb{R}$.}

    Avec $a=2$ et $b=2$, la fonction $h(x)$ s'écrit:
    $$ h(x) = \begin{cases} \frac{\sin(2x)}{x} & \text{si } x < 0 \\ 2 & \text{si } x = 0 \\ \frac{e^{2x} - 1}{x} & \text{si } x > 0 \end{cases} $$
    Nous allons étudier la continuité sur les intervalles ouverts $(-\infty, 0)$, $(0, +\infty)$ et au point $x=0$.

    *   \textbf{Continuité sur $(-\infty, 0)$}:
        Pour $x < 0$, $h(x) = \frac{\sin(2x)}{x}$.
        La fonction $x \mapsto 2x$ est un polynôme, donc elle est continue sur $\mathbb{R}$.
        La fonction $u \mapsto \sin(u)$ est continue sur $\mathbb{R}$.
        Par composition, la fonction $x \mapsto \sin(2x)$ est continue sur $\mathbb{R}$.
        La fonction $x \mapsto x$ est un polynôme, donc elle est continue sur $\mathbb{R}$.
        Pour $x < 0$, le dénominateur $x$ est non nul.
        Puisque $h(x)$ est le quotient de deux fonctions continues dont le dénominateur est non nul sur $(-\infty, 0)$, $h$ est continue sur $(-\infty, 0)$.

    *   \textbf{Continuité sur $(0, +\infty)$}:
        Pour $x > 0$, $h(x) = \frac{e^{2x} - 1}{x}$.
        La fonction $x \mapsto 2x$ est continue sur $\mathbb{R}$.
        La fonction $u \mapsto e^u$ est continue sur $\mathbb{R}$.
        Par composition, la fonction $x \mapsto e^{2x}$ est continue sur $\mathbb{R}$.
        La fonction $x \mapsto e^{2x} - 1$ est la différence de fonctions continues, donc elle est continue sur $\mathbb{R}$.
        La fonction $x \mapsto x$ est continue sur $\mathbb{R}$.
        Pour $x > 0$, le dénominateur $x$ est non nul.
        Puisque $h(x)$ est le quotient de deux fonctions continues dont le dénominateur est non nul sur $(0, +\infty)$, $h$ est continue sur $(0, +\infty)$.

    *   \textbf{Continuité en $x=0$}:
        D'après la question B.1, nous avons choisi $a=2$ et $b=2$ précisément pour assurer la continuité en $x=0$.
        Nous avons montré que $\lim_{x \to 0^-} h(x) = a = 2$.
        Nous avons montré que $\lim_{x \to 0^+} h(x) = 2$.
        Et $h(0) = b = 2$.
        Puisque $\lim_{x \to 0^-} h(x) = \lim_{x \to 0^+} h(x) = h(0) = 2$, la fonction $h$ est continue en $x=0$.

    En combinant ces résultats, la fonction $h$ est continue sur $(-\infty, 0)$, sur $(0, +\infty)$ et au point $x=0$. Par conséquent, $h$ est continue sur $\mathbb{R}$.

\subsubsection{Partie C: Application du Théorème des Valeurs Intermédiaires (TVI) et Propriétés Globales}

1.  \textbf{Soit $P(x)$ un polynôme de degré impair. Montrer que $P(x)$ admet au moins une racine réelle.}

    Soit $P(x) = a_n x^n + a_{n-1} x^{n-1} + \dots + a_1 x + a_0$, où $a_n \neq 0$ et $n$ est un entier impair.
    Les polynômes sont des fonctions continues sur $\mathbb{R}$.
    Pour montrer que $P(x)$ admet au moins une racine réelle, nous allons utiliser le Théorème des Valeurs Intermédiaires (TVI). Le TVI stipule que si une fonction $f$ est continue sur un intervalle $[u, v]$, alors pour toute valeur $k$ entre $f(u)$ et $f(v)$, il existe au moins un $c \in [u, v]$ tel que $f(c) = k$. Dans notre cas, nous voulons montrer l'existence d'un $c$ tel que $P(c) = 0$.

    Considérons le comportement de $P(x)$ lorsque $x \to \pm \infty$. Le comportement d'un polynôme est dominé par son terme de plus haut degré.
    $$ \lim_{x \to +\infty} P(x) = \lim_{x \to +\infty} a_n x^n $$
    $$ \lim_{x \to -\infty} P(x) = \lim_{x \to -\infty} a_n x^n $$
    Puisque $n$ est un entier impair:
    *   Si $a_n > 0$:
        $\lim_{x \to +\infty} x^n = +\infty$, donc $\lim_{x \to +\infty} P(x) = +\infty$.
        $\lim_{x \to -\infty} x^n = -\infty$, donc $\lim_{x \to -\infty} P(x) = -\infty$.
    *   Si $a_n < 0$:
        $\lim_{x \to +\infty} x^n = +\infty$, donc $\lim_{x \to +\infty} P(x) = -\infty$.
        $\lim_{x \to -\infty} x^n = -\infty$, donc $\lim_{x \to -\infty} P(x) = +\infty$.

    Dans les deux cas, les limites de $P(x)$ en $+\infty$ et en $-\infty$ sont de signes opposés.
    *   Si $a_n > 0$:
        Puisque $\lim_{x \to +\infty} P(x) = +\infty$, il existe un nombre réel $x_1$ suffisamment grand tel que $P(x_1) > 0$.
        Puisque $\lim_{x \to -\infty} P(x) = -\infty$, il existe un nombre réel $x_2$ suffisamment petit (négatif) tel que $P(x_2) < 0$.
        Nous avons donc $P(x_2) < 0 < P(x_1)$.
    *   Si $a_n < 0$:
        Puisque $\lim_{x \to +\infty} P(x) = -\infty$, il existe un nombre réel $x_1$ suffisamment grand tel que $P(x_1) < 0$.
        Puisque $\lim_{x \to -\infty} P(x) = +\infty$, il existe un nombre réel $x_2$ suffisamment petit (négatif) tel que $P(x_2) > 0$.
        Nous avons donc $P(x_1) < 0 < P(x_2)$.

    Dans les deux scénarios, nous pouvons trouver un intervalle $[x_a, x_b]$ (en prenant $x_a = \min(x_1, x_2)$ et $x_b = \max(x_1, x_2)$) tel que $P(x_a)$ et $P(x_b)$ sont de signes opposés.
    Puisque $P(x)$ est une fonction continue sur $\mathbb{R}$, elle est en particulier continue sur l'intervalle fermé $[x_a, x_b]$.
    Comme $P(x_a)$ et $P(x_b)$ sont de signes opposés, $0$ est une valeur intermédiaire entre $P(x_a)$ et $P(x_b)$.
    Par le Théorème des Valeurs Intermédiaires, il existe au moins un $c \in (x_a, x_b)$ tel que $P(c) = 0$.
    Ce $c$ est une racine réelle du polynôme $P(x)$.

2.  \textbf{Soit $f: [0, 1] \to [0, 1]$ une fonction continue. Montrer qu'il existe au moins un point fixe $c \in [0, 1]$ tel que $f(c) = c$.}

    Un point fixe $c$ est une valeur telle que $f(c) = c$. Cela est équivalent à $f(c) - c = 0$.
    Considérons la fonction auxiliaire $g: [0, 1] \to \mathbb{R}$ définie par $g(x) = f(x) - x$.
    *   \textbf{Continuité de $g$}:
        La fonction $f$ est continue sur $[0, 1]$ par hypothèse.
        La fonction $x \mapsto x$ est un polynôme, donc elle est continue sur $\mathbb{R}$, et en particulier sur $[0, 1]$.
        La fonction $g(x)$ est la différence de deux fonctions continues sur $[0, 1]$, elle est donc continue sur $[0, 1]$.

    *   \textbf{Évaluation de $g$ aux bornes de l'intervalle}:
        Calculons les valeurs de $g(x)$ aux extrémités de l'intervalle $[0, 1]$:
        *   $g(0) = f(0) - 0 = f(0)$.
            Puisque $f: [0, 1] \to [0, 1]$, cela signifie que l'image de $f$ est contenue dans $[0, 1]$. Donc $0 \le f(0) \le 1$.
            Par conséquent, $g(0) = f(0) \ge 0$.
        *   $g(1) = f(1) - 1$.
            Puisque $f: [0, 1] \to [0, 1]$, nous avons $0 \le f(1) \le 1$.
            Par conséquent, $f(1) - 1 \le 1 - 1 = 0$.
            Donc, $g(1) = f(1) - 1 \le 0$.

    *   \textbf{Application du Théorème des Valeurs Intermédiaires}:
        Nous avons $g(0) \ge 0$ et $g(1) \le 0$.
        Trois cas sont possibles:
        1.  Si $g(0) = 0$, alors $f(0) = 0$. Dans ce cas, $c=0$ est un point fixe.
        2.  Si $g(1) = 0$, alors $f(1) = 1$. Dans ce cas, $c=1$ est un point fixe.
        3.  Si $g(0) > 0$ et $g(1) < 0$. Dans ce cas, nous avons $g(1) < 0 < g(0)$.
            Puisque $g$ est continue sur l'intervalle fermé $[0, 1]$ et que $0$ est une valeur intermédiaire entre $g(1)$ et $g(0)$, le Théorème des Valeurs Intermédiaires garantit l'existence d'au moins un $c \in (0, 1)$ tel que $g(c) = 0$.
            Si $g(c) = 0$, alors $f(c) - c = 0$, ce qui signifie $f(c) = c$.

    Dans tous les cas, il existe au moins un point $c \in [0, 1]$ tel que $f(c) = c$.

3.  \textbf{Soit $f: \mathbb{R} \to \mathbb{R}$ une fonction continue telle que $\lim_{x \to -\infty} f(x) = -\infty$ et $\lim_{x \to +\infty} f(x) = +\infty$. Montrer que pour tout $y \in \mathbb{R}$, il existe $x \in \mathbb{R}$ tel que $f(x) = y$.}

    Nous voulons montrer que la fonction $f$ est surjective, c'est-à-dire que pour tout $y \in \mathbb{R}$ (valeur cible), il existe au moins un $x \in \mathbb{R}$ (antécédent) tel que $f(x) = y$.

    Soit $y \in \mathbb{R}$ une valeur arbitraire.
    *   \textbf{Utilisation de la limite en $-\infty$}:
        Puisque $\lim_{x \to -\infty} f(x) = -\infty$, par définition de la limite, pour tout nombre réel $M_1$, il existe un $A \in \mathbb{R}$ tel que pour tout $x < A$, $f(x) < M_1$.
        En particulier, choisissons $M_1 = y$. Il existe donc un $x_1 \in \mathbb{R}$ (en prenant $x_1 < A$) tel que $f(x_1) < y$.

    *   \textbf{Utilisation de la limite en $+\infty$}:
        Puisque $\lim_{x \to +\infty} f(x) = +\infty$, par définition de la limite, pour tout nombre réel $M_2$, il existe un $B \in \mathbb{R}$ tel que pour tout $x > B$, $f(x) > M_2$.
        En particulier, choisissons $M_2 = y$. Il existe donc un $x_2 \in \mathbb{R}$ (en prenant $x_2 > B$) tel que $f(x_2) > y$.

    *   \textbf{Application du Théorème des Valeurs Intermédiaires}:
        Nous avons trouvé deux points $x_1$ et $x_2$ tels que $f(x_1) < y$ et $f(x_2) > y$. Nous pouvons toujours choisir $x_1 < x_2$ en prenant $A$ suffisamment petit et $B$ suffisamment grand.
        La fonction $f$ est continue sur $\mathbb{R}$ par hypothèse, donc elle est continue sur l'intervalle fermé $[x_1, x_2]$.
        Puisque $f(x_1) < y < f(x_2)$, la valeur $y$ est une valeur intermédiaire entre $f(x_1)$ et $f(x_2)$.
        Par le Théorème des Valeurs Intermédiaires, il existe au moins un $c \in (x_1, x_2)$ tel que $f(c) = y$.

    Puisque $y$ a été choisi arbitrairement dans $\mathbb{R}$, nous avons montré que pour toute valeur réelle $y$, il existe un $x \in \mathbb{R}$ tel que $f(x) = y$. Cela signifie que $f$ est surjective.

\subsubsection{Partie D: Un Défi Théorique (Équation Fonctionnelle de Cauchy)}

Soit $f: \mathbb{R} \to \mathbb{R}$ une fonction continue telle que $f(x+y) = f(x) + f(y)$ pour tout $x, y \in \mathbb{R}$.

1.  \textbf{Montrer que $f(nx) = nf(x)$ pour tout $n \in \mathbb{Z}$ et $x \in \mathbb{R}$.}

    Nous allons démontrer cette propriété par étapes.

    *   \textbf{Cas $n \in \mathbb{N}$ (entiers naturels)}:
        Nous utilisons une preuve par récurrence sur $n$.
        *   \textbf{Initialisation}: Pour $n=1$, $f(1x) = f(x)$, et $1f(x) = f(x)$. Donc $f(1x) = 1f(x)$ est vraie.
        *   \textbf{Hérédité}: Supposons que $f(kx) = kf(x)$ est vraie pour un certain entier naturel $k \ge 1$.
            Nous voulons montrer que $f((k+1)x) = (k+1)f(x)$.
            $f((k+1)x) = f(kx + x)$.
            En utilisant la propriété donnée $f(u+v) = f(u) + f(v)$ avec $u=kx$ et $v=x$:
            $f(kx + x) = f(kx) + f(x)$.
            Par l'hypothèse de récurrence, $f(kx) = kf(x)$.
            Donc, $f(kx) + f(x) = kf(x) + f(x) = (k+1)f(x)$.
            Ainsi, $f((k+1)x) = (k+1)f(x)$.
        *   \textbf{Conclusion}: Par le principe de récurrence, $f(nx) = nf(x)$ pour tout $n \in \mathbb{N}^*$.

    *   \textbf{Cas $n=0$}:
        Posons $x=0$ dans l'équation fonctionnelle: $f(0+0) = f(0) + f(0)$, ce qui donne $f(0) = 2f(0)$.
        En soustrayant $f(0)$ des deux côtés, nous obtenons $f(0) = 0$.
        Alors $f(0x) = f(0) = 0$. Et $0f(x) = 0$.
        Donc $f(0x) = 0f(x)$ est vraie.

    *   \textbf{Cas $n \in \mathbb{Z}^-$ (entiers négatifs)}:
        Soit $n$ un entier négatif. Alors $n = -m$ pour un certain entier naturel $m \in \mathbb{N}^*$.
        Nous savons que $f(0) = 0$.
        Nous pouvons écrire $f(0) = f(mx + (-m)x)$.
        En utilisant la propriété $f(u+v) = f(u) + f(v)$ avec $u=mx$ et $v=(-m)x$:
        $f(0) = f(mx) + f((-m)x)$.
        Puisque $f(0)=0$ et que $m \in \mathbb{N}^*$, nous avons $f(mx) = mf(x)$ d'après le cas des entiers naturels.
        Donc, $0 = mf(x) + f((-m)x)$.
        Cela implique $f((-m)x) = -mf(x)$.
        En substituant $n = -m$, nous obtenons $f(nx) = nf(x)$.

    En combinant tous les cas, nous avons montré que $f(nx) = nf(x)$ pour tout $n \in \mathbb{Z}$ et $x \in \mathbb{R}$.

2.  \textbf{Montrer que $f(qx) = qf(x)$ pour tout $q \in \mathbb{Q}$ et $x \in \mathbb{R}$.}

    Soit $q \in \mathbb{Q}$. Par définition, $q$ peut s'écrire sous la forme $\frac{m}{n}$, où $m \in \mathbb{Z}$ et $n \in \mathbb{N}^*$ (donc $n \neq 0$).
    Nous voulons montrer que $f\left(\frac{m}{n}x\right) = \frac{m}{n}f(x)$.

    De la question D.1, nous savons que $f(kx) = kf(x)$ pour tout $k \in \mathbb{Z}$.
    Considérons l'expression $f(nx)$. Nous pouvons écrire $nx = n \left(\frac{m}{n}x\right)$.
    En utilisant la propriété pour l'entier $n$:
    $$ f\left(n \left(\frac{m}{n}x\right)\right) = n f\left(\frac{m}{n}x\right) $$
    D'autre part, nous pouvons écrire $f(nx)$ comme $f(mx)$ en considérant $y = \frac{m}{n}x$, alors $ny = mx$.
    Donc, $f(ny) = nf(y)$.
    Et $f(mx) = mf(x)$.
    Ainsi, $nf\left(\frac{m}{n}x\right) = mf(x)$.
    Puisque $n \neq 0$, nous pouvons diviser par $n$:
    $$ f\left(\frac{m}{n}x\right) = \frac{m}{n}f(x) $$
    Donc, $f(qx) = qf(x)$ pour tout $q \in \mathbb{Q}$ et $x \in \mathbb{R}$.

    Un cas particulier important est lorsque $x=1$. Posons $c = f(1)$.
    Alors pour tout $q \in \mathbb{Q}$, $f(q) = f(q \cdot 1) = qf(1) = cq$.

3.  \textbf{En utilisant la continuité de $f$, montrer que $f(x) = cx$ pour une certaine constante $c \in \mathbb{R}$ et pour tout $x \in \mathbb{R}$.}

    De la question D.2, nous avons établi que $f(x) = cx$ pour tout $x \in \mathbb{Q}$, où $c = f(1)$.
    Nous voulons maintenant étendre cette propriété à tous les nombres réels $x \in \mathbb{R}$. C'est ici que la continuité de $f$ est cruciale.

    Soit $x \in \mathbb{R}$ un nombre réel arbitraire.
    Nous savons que l'ensemble des nombres rationnels $\mathbb{Q}$ est dense dans $\mathbb{R}$. Cela signifie que pour tout nombre réel $x$, il existe une suite de nombres rationnels $(q_k)_{k \in \mathbb{N}}$ telle que $\lim_{k \to \infty} q_k = x$.

    Puisque $f$ est continue sur $\mathbb{R}$ (par hypothèse), la continuité de $f$ implique que si une suite $(q_k)$ converge vers $x$, alors la suite des images $(f(q_k))$ converge vers $f(x)$.
    C'est-à-dire:
    $$ \lim_{k \to \infty} f(q_k) = f\left(\lim_{k \to \infty} q_k\right) = f(x) $$
    D'autre part, puisque chaque $q_k$ est un nombre rationnel, nous savons d'après la question D.2 que $f(q_k) = cq_k$ (où $c = f(1)$).
    Donc, nous pouvons écrire:
    $$ \lim_{k \to \infty} f(q_k) = \lim_{k \to \infty} (cq_k) $$
    Puisque $c$ est une constante, nous pouvons la sortir de la limite:
    $$ \lim_{k \to \infty} (cq_k) = c \lim_{k \to \infty} q_k $$
    Comme $\lim_{k \to \infty} q_k = x$, nous avons:
    $$ c \lim_{k \to \infty} q_k = cx $$
    En égalant les deux expressions pour la limite de $f(q_k)$:
    $$ f(x) = cx $$
    Cette relation est vraie pour tout $x \in \mathbb{R}$.
    La constante $c$ est déterminée par la valeur de $f(1)$. Si $f(1)=0$, alors $f(x)=0$ pour tout $x$. Si $f(1)=1$, alors $f(x)=x$ pour tout $x$.

    Ainsi, toute fonction continue $f: \mathbb{R} \to \mathbb{R}$ satisfaisant l'équation fonctionnelle de Cauchy $f(x+y) = f(x) + f(y)$ est nécessairement de la forme $f(x) = cx$ pour une constante $c \in \mathbb{R}$.

\section{Exo-06}
\section{Exercice 6 - Analyse Approfondie de la Continuité d'une Fonction Définie par Morceaux}

Considérons la fonction $f: \mathbb{R} \to \mathbb{R}$ définie par :
$$ f(x) = \begin{cases} \frac{\sin(3x)}{x} & \text{si } x < 0 \\ A & \text{si } x = 0 \\ \frac{e^{3x} - 1}{x} & \text{si } x > 0 \end{cases} $$
où $A$ est un paramètre réel.

---

\subsubsection{Question 1 : Détermination du paramètre $A$}

Déterminer la valeur du paramètre réel $A$ pour laquelle la fonction $f$ est continue au point $x_0 = 0$. Justifier rigoureusement chaque étape de calcul de limite.

\paragraph{Correction Question 1}

Pour que la fonction $f$ soit continue au point $x_0 = 0$, il est nécessaire et suffisant que les trois conditions suivantes soient satisfaites :
1.  $f(0)$ est défini. (C'est le cas, $f(0) = A$).
2.  La limite de $f(x)$ lorsque $x$ tend vers $0$ existe. Cela signifie que les limites à gauche et à droite doivent exister et être égales.
    $$ \lim_{x \to 0^-} f(x) = \lim_{x \to 0^+} f(x) $$
3.  Cette limite doit être égale à $f(0)$.
    $$ \lim_{x \to 0} f(x) = f(0) $$

Calculons la limite à gauche :
$$ \lim_{x \to 0^-} f(x) = \lim_{x \to 0^-} \frac{\sin(3x)}{x} $$
Pour évaluer cette limite, nous allons utiliser un changement de variable.
Soit $u = 3x$. Lorsque $x \to 0^-$, $u$ tend également vers $0^-$.
L'expression devient :
$$ \lim_{u \to 0^-} \frac{\sin(u)}{u/3} = \lim_{u \to 0^-} 3 \cdot \frac{\sin(u)}{u} $$
Nous utilisons le résultat fondamental de limite bien connu : $\lim_{u \to 0} \frac{\sin(u)}{u} = 1$.
Par conséquent,
$$ \lim_{x \to 0^-} f(x) = 3 \cdot 1 = 3 $$

Calculons la limite à droite :
$$ \lim_{x \to 0^+} f(x) = \lim_{x \to 0^+} \frac{e^{3x} - 1}{x} $$
Pour évaluer cette limite, nous allons utiliser un changement de variable.
Soit $v = 3x$. Lorsque $x \to 0^+$, $v$ tend également vers $0^+$.
L'expression devient :
$$ \lim_{v \to 0^+} \frac{e^{v} - 1}{v/3} = \lim_{v \to 0^+} 3 \cdot \frac{e^{v} - 1}{v} $$
Nous utilisons le résultat fondamental de limite bien connu : $\lim_{v \to 0} \frac{e^{v} - 1}{v} = 1$.
Par conséquent,
$$ \lim_{x \to 0^+} f(x) = 3 \cdot 1 = 3 $$

Pour que la limite de $f(x)$ lorsque $x$ tend vers $0$ existe, il faut que les limites à gauche et à droite soient égales. Dans notre cas, $\lim_{x \to 0^-} f(x) = 3$ et $\lim_{x \to 0^+} f(x) = 3$. Elles sont égales.
Donc, $\lim_{x \to 0} f(x) = 3$.

Pour que $f$ soit continue en $x_0 = 0$, il faut que $\lim_{x \to 0} f(x) = f(0)$.
Nous avons $f(0) = A$.
Ainsi, nous devons avoir $A = 3$.

La valeur du paramètre $A$ pour laquelle la fonction $f$ est continue au point $x_0 = 0$ est $A=3$.

---

\subsubsection{Question 2 : Preuve $\epsilon-\delta$ de la continuité}

Pour la valeur de $A$ trouvée à la question précédente, c'est-à-dire $A=3$, démontrer la continuité de $f$ au point $x_0 = 0$ en utilisant la définition formelle $(\epsilon, \delta)$ de la continuité. Chaque étape doit être explicitement justifiée.

\paragraph{Correction Question 2}

Nous devons démontrer que $f$ est continue en $x_0 = 0$ avec $f(0) = 3$.
La définition formelle de la continuité en un point $x_0$ est la suivante :
Pour tout $\epsilon > 0$, il existe un $\delta > 0$ tel que pour tout $x \in \mathbb{R}$, si $|x - x_0| < \delta$, alors $|f(x) - f(x_0)| < \epsilon$.
Dans notre cas, $x_0 = 0$ et $f(x_0) = f(0) = 3$. Nous devons donc montrer que pour tout $\epsilon > 0$, il existe un $\delta > 0$ tel que pour tout $x \in \mathbb{R}$, si $|x| < \delta$, alors $|f(x) - 3| < \epsilon$.

Soit $\epsilon > 0$ un nombre réel arbitrairement petit. Nous devons trouver un $\delta > 0$.

Nous allons analyser trois cas pour $x$: $x=0$, $x<0$ et $x>0$.

\textbf{Cas 1 : $x = 0$}
Si $x=0$, alors $|f(0) - 3| = |3 - 3| = 0$. Puisque $0 < \epsilon$ pour tout $\epsilon > 0$, la condition est satisfaite pour $x=0$.

\textbf{Cas 2 : $x < 0$}
Pour $x < 0$, $f(x) = \frac{\sin(3x)}{x}$. Nous devons montrer que $|\frac{\sin(3x)}{x} - 3| < \epsilon$.
Nous savons que $\frac{\sin(3x)}{x} - 3 = 3 \left( \frac{\sin(3x)}{3x} - 1 \right)$.
Donc, nous devons montrer que $\left| 3 \left( \frac{\sin(3x)}{3x} - 1 \right) \right| < \epsilon$, ce qui est équivalent à $\left| \frac{\sin(3x)}{3x} - 1 \right| < \frac{\epsilon}{3}$.

Pour $u \in ]-\pi/2, \pi/2[ \setminus \{0\}$, nous utilisons l'inégalité fondamentale suivante :
$1 - \frac{u^2}{2} < \frac{\sin u}{u} < 1$ pour $u \in ]0, \pi/2[$ et $1 < \frac{\sin u}{u} < 1 - \frac{u^2}{2}$ pour $u \in ]-\pi/2, 0[$.
Plus précisément, pour $u \in ]-\pi/2, \pi/2[ \setminus \{0\}$, nous avons l'inégalité $|\frac{\sin u}{u} - 1| < \frac{u^2}{2}$.
(Cette inégalité peut être démontrée en utilisant le développement en série de Taylor de $\sin u = u - \frac{u^3}{3!} + \dots$ ou par des arguments géométriques et le théorème des accroissements finis si ces outils sont disponibles. Dans le cadre de la continuité, elle est souvent admise ou démontrée par des moyens élémentaires.)

Soit $u = 3x$. Puisque $x < 0$, $u < 0$.
Nous voulons que $|u| < \pi/2$, ce qui signifie $|3x| < \pi/2$, donc $|x| < \pi/6$.
Si $|x| < \pi/6$, alors $|3x| < \pi/2$, et nous pouvons appliquer l'inégalité :
$$ \left| \frac{\sin(3x)}{3x} - 1 \right| < \frac{(3x)^2}{2} = \frac{9x^2}{2} $$
Nous voulons que cette quantité soit inférieure à $\frac{\epsilon}{3}$.
$$ \frac{9x^2}{2} < \frac{\epsilon}{3} $$
$$ 9x^2 < \frac{2\epsilon}{3} $$
$$ x^2 < \frac{2\epsilon}{27} $$
$$ |x| < \sqrt{\frac{2\epsilon}{27}} $$
Soit $\delta_1 = \min\left(\frac{\pi}{6}, \sqrt{\frac{2\epsilon}{27}}\right)$. Si $0 < |x| < \delta_1$ et $x<0$, alors $|f(x) - 3| < \epsilon$.

\textbf{Cas 3 : $x > 0$}
Pour $x > 0$, $f(x) = \frac{e^{3x} - 1}{x}$. Nous devons montrer que $|\frac{e^{3x} - 1}{x} - 3| < \epsilon$.
Nous savons que $\frac{e^{3x} - 1}{x} - 3 = 3 \left( \frac{e^{3x} - 1}{3x} - 1 \right)$.
Donc, nous devons montrer que $\left| 3 \left( \frac{e^{3x} - 1}{3x} - 1 \right) \right| < \epsilon$, ce qui est équivalent à $\left| \frac{e^{3x} - 1}{3x} - 1 \right| < \frac{\epsilon}{3}$.

Pour $u > 0$ et $u$ suffisamment petit (par exemple $u \in ]0, 1[$), nous utilisons l'inégalité fondamentale suivante :
$1 < \frac{e^u - 1}{u} < e^u$.
(Cette inégalité peut être démontrée en utilisant le développement en série de Taylor de $e^u = 1 + u + \frac{u^2}{2!} + \dots$ ou par des arguments de convexité et le théorème des accroissements finis.)
De cette inégalité, nous déduisons que $0 < \frac{e^u - 1}{u} - 1 < e^u - 1$.
Nous voulons que $e^u - 1 < \frac{\epsilon}{3}$.
Soit $u = 3x$. Puisque $x > 0$, $u > 0$.
Nous voulons que $u$ soit suffisamment petit. Par exemple, si $u < 1$, alors $e^u - 1 < e - 1 \approx 1.718$.
Nous voulons que $e^{3x} - 1 < \frac{\epsilon}{3}$.
Puisque la fonction $y \mapsto e^y - 1$ est strictement croissante, nous pouvons prendre le logarithme naturel :
$3x < \ln\left(1 + \frac{\epsilon}{3}\right)$.
$$ x < \frac{1}{3} \ln\left(1 + \frac{\epsilon}{3}\right) $$
Soit $\delta_2 = \min\left(\frac{1}{3}, \frac{1}{3} \ln\left(1 + \frac{\epsilon}{3}\right)\right)$. Si $0 < x < \delta_2$, alors $|f(x) - 3| < \epsilon$.
(Note : $\ln(1+y) \approx y$ pour $y$ petit, donc $\frac{1}{3} \ln(1+\frac{\epsilon}{3}) \approx \frac{1}{3} \frac{\epsilon}{3} = \frac{\epsilon}{9}$. Nous pouvons donc choisir $\delta_2$ de l'ordre de $\epsilon/9$ pour $\epsilon$ petit.)

\textbf{Conclusion pour $\delta$}
Nous avons trouvé $\delta_1$ pour $x<0$ et $\delta_2$ pour $x>0$.
Pour que la condition $|f(x) - 3| < \epsilon$ soit satisfaite pour tout $x$ tel que $|x| < \delta$, nous devons choisir $\delta$ comme le minimum de ces valeurs.
Soit $\delta = \min\left(\delta_1, \delta_2\right) = \min\left(\frac{\pi}{6}, \sqrt{\frac{2\epsilon}{27}}, \frac{1}{3}, \frac{1}{3} \ln\left(1 + \frac{\epsilon}{3}\right)\right)$.
Avec ce choix de $\delta$, pour tout $x \in \mathbb{R}$ tel que $|x| < \delta$:
*   Si $x=0$, $|f(0)-3|=0 < \epsilon$.
*   Si $x<0$, alors $|x| < \delta_1$, ce qui implique $|f(x)-3| < \epsilon$.
*   Si $x>0$, alors $|x| < \delta_2$, ce qui implique $|f(x)-3| < \epsilon$.

Par conséquent, pour tout $\epsilon > 0$, il existe un $\delta > 0$ (défini comme ci-dessus) tel que si $|x - 0| < \delta$, alors $|f(x) - f(0)| < \epsilon$.
Ceci démontre que la fonction $f$ est continue au point $x_0 = 0$ pour $A=3$.

---

\subsubsection{Question 3 : Continuité sur les intervalles ouverts}

Étudier la continuité de $f$ sur les intervalles $]-\infty, 0[$ et $]0, +\infty[$.

\paragraph{Correction Question 3}

\textbf{Sur l'intervalle $]-\infty, 0[$ :}
Pour $x \in ]-\infty, 0[$, la fonction $f(x)$ est définie par $f(x) = \frac{\sin(3x)}{x}$.
Nous pouvons écrire $f(x) = \frac{1}{x} \cdot \sin(3x)$.
La fonction $x \mapsto x$ est une fonction polynomiale, donc elle est continue sur $\mathbb{R}$. En particulier, elle est continue et non nulle sur $]-\infty, 0[$.
La fonction $x \mapsto 3x$ est une fonction polynomiale, donc elle est continue sur $\mathbb{R}$.
La fonction $u \mapsto \sin(u)$ est une fonction trigonométrique, donc elle est continue sur $\mathbb{R}$.
Par composition, la fonction $x \mapsto \sin(3x)$ est continue sur $\mathbb{R}$.
Puisque $x \mapsto \sin(3x)$ est continue sur $]-\infty, 0[$ et $x \mapsto x$ est continue et non nulle sur $]-\infty, 0[$, le quotient $f(x) = \frac{\sin(3x)}{x}$ est continu sur $]-\infty, 0[$ en tant que quotient de fonctions continues dont le dénominateur ne s'annule pas.

\textbf{Sur l'intervalle $]0, +\infty[$ :}
Pour $x \in ]0, +\infty[$, la fonction $f(x)$ est définie par $f(x) = \frac{e^{3x} - 1}{x}$.
Nous pouvons écrire $f(x) = \frac{1}{x} \cdot (e^{3x} - 1)$.
La fonction $x \mapsto x$ est une fonction polynomiale, donc elle est continue sur $\mathbb{R}$. En particulier, elle est continue et non nulle sur $]0, +\infty[$.
La fonction $x \mapsto 3x$ est une fonction polynomiale, donc elle est continue sur $\mathbb{R}$.
La fonction $u \mapsto e^u$ est une fonction exponentielle, donc elle est continue sur $\mathbb{R}$.
Par composition, la fonction $x \mapsto e^{3x}$ est continue sur $\mathbb{R}$.
La fonction $x \mapsto e^{3x} - 1$ est continue sur $\mathbb{R}$ en tant que différence de fonctions continues.
Puisque $x \mapsto e^{3x} - 1$ est continue sur $]0, +\infty[$ et $x \mapsto x$ est continue et non nulle sur $]0, +\infty[$, le quotient $f(x) = \frac{e^{3x} - 1}{x}$ est continu sur $]0, +\infty[$ en tant que quotient de fonctions continues dont le dénominateur ne s'annule pas.

En résumé, la fonction $f$ est continue sur $]-\infty, 0[$ et sur $]0, +\infty[$.

---

\subsubsection{Question 4 : Prolongement par continuité et application du TVI}

Soit $g(x)$ la fonction définie par $g(x) = f(x)$ pour $x \neq 0$ et $g(0) = A$, où $A$ est la valeur trouvée à la Question 1.
a) Montrer que $g$ est continue sur $\mathbb{R}$.
b) Démontrer que l'équation $g(x) = 2.5$ admet au moins une solution dans l'intervalle $]-1, 1[$. Justifier l'utilisation du Théorème des Valeurs Intermédiaires (TVI).

\paragraph{Correction Question 4}

a) \textbf{Montrer que $g$ est continue sur $\mathbb{R}$.}
La fonction $g(x)$ est définie comme :
$$ g(x) = \begin{cases} \frac{\sin(3x)}{x} & \text{si } x < 0 \\ 3 & \text{si } x = 0 \\ \frac{e^{3x} - 1}{x} & \text{si } x > 0 \end{cases} $$
D'après la Question 3, la fonction $f$ (et donc $g$) est continue sur $]-\infty, 0[$ et sur $]0, +\infty[$.
D'après la Question 2, la fonction $f$ (et donc $g$) est continue au point $x_0 = 0$ pour $A=3$.
Puisque $g$ est continue sur $]-\infty, 0[$, sur $]0, +\infty[$ et au point $0$, elle est continue sur l'union de ces ensembles, c'est-à-dire sur $\mathbb{R}$.
Donc, $g$ est continue sur $\mathbb{R}$.

b) \textbf{Démontrer que l'équation $g(x) = 2.5$ admet au moins une solution dans l'intervalle $]-1, 1[$.}
Nous allons utiliser le Théorème des Valeurs Intermédiaires (TVI).
Le TVI stipule que si une fonction $h$ est continue sur un intervalle fermé $[a, b]$, alors pour toute valeur $k$ comprise entre $h(a)$ et $h(b)$ (inclusivement), il existe au moins un $c \in [a, b]$ tel que $h(c) = k$.

Appliquons le TVI à la fonction $g$ sur l'intervalle $[-1, 0]$.
1.  \textbf{Continuité de $g$ sur $[-1, 0]$ :}
    Nous avons montré en 4.a) que $g$ est continue sur $\mathbb{R}$. Par conséquent, $g$ est continue sur l'intervalle fermé $[-1, 0]$.

2.  \textbf{Calcul des valeurs aux bornes de l'intervalle :}
    *   Calculons $g(-1)$ :
        Puisque $-1 < 0$, $g(-1) = \frac{\sin(3 \cdot (-1))}{-1} = \frac{\sin(-3)}{-1} = \frac{-\sin(3)}{-1} = \sin(3)$.
        La valeur $3$ est en radians. Sachant que $\pi \approx 3.14159$, $3$ radians est légèrement inférieur à $\pi$.
        $\sin(3) \approx 0.14112$.
    *   Calculons $g(0)$ :
        Par définition, $g(0) = 3$.

3.  \textbf{Vérification de la condition du TVI :}
    Nous cherchons une solution à l'équation $g(x) = 2.5$.
    Nous avons $g(-1) = \sin(3) \approx 0.14112$ et $g(0) = 3$.
    Nous constatons que $g(-1) < 2.5 < g(0)$.
    Plus précisément, $0.14112 < 2.5 < 3$.

4.  \textbf{Conclusion par le TVI :}
    Puisque $g$ est continue sur l'intervalle fermé $[-1, 0]$ et que la valeur $2.5$ est strictement comprise entre $g(-1)$ et $g(0)$, le Théorème des Valeurs Intermédiaires garantit qu'il existe au moins un nombre réel $c \in ]-1, 0[$ tel que $g(c) = 2.5$.
    L'intervalle $]-1, 0[$ est un sous-ensemble de $]-1, 1[$.
    Par conséquent, l'équation $g(x) = 2.5$ admet au moins une solution dans l'intervalle $]-1, 1[$.

\section{Exo-07}
\section{Exercice 7 - Exploration Approfondie de la Continuité des Fonctions d'une Variable Réelle}

Cet exercice est conçu pour évaluer votre maîtrise des concepts fondamentaux et avancés de la continuité des fonctions d'une variable réelle. Il requiert une rigueur absolue dans la démonstration et l'application des théorèmes.

---

\subsection{Énoncé de l'Exercice}

\textbf{Partie A : Démonstration par la Définition Epsilon-Delta}

Soit la fonction $f: \mathbb{R} \to \mathbb{R}$ définie par $f(x) = x^3$.
Démontrez, en utilisant la définition formelle $(\epsilon, \delta)$ de la continuité, que la fonction $f$ est continue en tout point $x_0 \in \mathbb{R}$.

\textbf{Partie B : Continuité d'une Fonction Définie par Morceaux avec Paramètres}

Considérons la fonction $g: \mathbb{R} \to \mathbb{R}$ définie par :
$$
g(x) = \begin{cases}
ax^2 + 3x - 1 & \text{si } x < 1 \\
b & \text{si } x = 1 \\
\frac{\sqrt{x+3}-2}{x-1} & \text{si } x > 1
\end{cases}
$$
Déterminez les valeurs des constantes réelles $a$ et $b$ pour lesquelles la fonction $g$ est continue en $x=1$.

\textbf{Partie C : Application du Théorème des Valeurs Intermédiaires}

Soit $h: [0,1] \to [0,1]$ une fonction continue sur l'intervalle fermé $[0,1]$.
Démontrez qu'il existe au moins un point $c \in [0,1]$ tel que $h(c) = c$. Un tel point $c$ est appelé point fixe de $h$.

---

\subsection{Correction Détaillée}

\subsubsection{Partie A : Démonstration par la Définition Epsilon-Delta}

Pour démontrer que $f(x) = x^3$ est continue en tout point $x_0 \in \mathbb{R}$, nous devons montrer que pour tout $\epsilon > 0$, il existe un $\delta > 0$ tel que si $|x - x_0| < \delta$, alors $|f(x) - f(x_0)| < \epsilon$.

Soit $x_0 \in \mathbb{R}$ un point arbitraire et soit $\epsilon > 0$ un nombre réel positif arbitraire.
Nous devons analyser l'expression $|f(x) - f(x_0)|$:
$$
|f(x) - f(x_0)| = |x^3 - x_0^3|
$$
Nous utilisons l'identité algébrique de la différence de cubes, $A^3 - B^3 = (A - B)(A^2 + AB + B^2)$.
$$
|x^3 - x_0^3| = |(x - x_0)(x^2 + x x_0 + x_0^2)|
$$
Par la propriété de la valeur absolue $|AB| = |A||B|$, nous obtenons :
$$
|x^3 - x_0^3| = |x - x_0| |x^2 + x x_0 + x_0^2|
$$
Nous voulons que cette expression soit inférieure à $\epsilon$. Nous avons un facteur $|x - x_0|$ que nous pouvons rendre arbitrairement petit en choisissant $\delta$. Le défi est de majorer le second facteur, $|x^2 + x x_0 + x_0^2|$.

Pour majorer $|x^2 + x x_0 + x_0^2|$, nous allons d'abord imposer une condition sur $\delta$. Choisissons $\delta \le 1$.
Si $|x - x_0| < \delta$ et $\delta \le 1$, alors $|x - x_0| < 1$.
Cette inégalité implique que $x_0 - 1 < x < x_0 + 1$.
En utilisant la propriété de la valeur absolue $|x| \le |x_0| + |x - x_0|$, nous avons :
$$
|x| < |x_0| + 1
$$
Maintenant, nous pouvons majorer le terme $|x^2 + x x_0 + x_0^2|$ en utilisant l'inégalité triangulaire $|A+B+C| \le |A| + |B| + |C|$ :
$$
|x^2 + x x_0 + x_0^2| \le |x^2| + |x x_0| + |x_0^2|
$$
Puisque $|A^2| = |A|^2$, nous avons :
$$
|x^2 + x x_0 + x_0^2| \le |x|^2 + |x||x_0| + |x_0|^2
$$
En substituant la majoration de $|x|$ que nous avons trouvée ($|x| < |x_0| + 1$) :
$$
|x^2 + x x_0 + x_0^2| < (|x_0| + 1)^2 + (|x_0| + 1)|x_0| + |x_0|^2
$$
Développons cette expression :
$$
(|x_0| + 1)^2 = x_0^2 + 2|x_0| + 1
$$
$$
(|x_0| + 1)|x_0| = |x_0|^2 + |x_0|
$$
En sommant ces termes :
$$
|x^2 + x x_0 + x_0^2| < (x_0^2 + 2|x_0| + 1) + (|x_0|^2 + |x_0|) + |x_0|^2
$$
$$
|x^2 + x x_0 + x_0^2| < 3x_0^2 + 3|x_0| + 1
$$
Soit $M = 3x_0^2 + 3|x_0| + 1$. Notez que $M$ est une constante positive qui dépend de $x_0$.
Alors, nous avons :
$$
|f(x) - f(x_0)| = |x - x_0| |x^2 + x x_0 + x_0^2| < |x - x_0| M
$$
Puisque nous avons $|x - x_0| < \delta$, nous obtenons :
$$
|f(x) - f(x_0)| < \delta M
$$
Nous voulons que cette expression soit inférieure à $\epsilon$. Donc, nous posons $\delta M < \epsilon$, ce qui implique $\delta < \frac{\epsilon}{M}$.

Pour satisfaire toutes les conditions, nous choisissons $\delta$ comme le minimum de notre condition initiale ($\delta \le 1$) et de la nouvelle condition ($\delta < \frac{\epsilon}{M}$).
$$
\delta = \min\left(1, \frac{\epsilon}{3x_0^2 + 3|x_0| + 1}\right)
$$
Avec ce choix de $\delta$, si $|x - x_0| < \delta$, alors :
1.  $|x - x_0| < 1$, ce qui nous a permis de majorer $|x^2 + x x_0 + x_0^2|$ par $M = 3x_0^2 + 3|x_0| + 1$.
2.  $|x - x_0| < \frac{\epsilon}{M}$.
Par conséquent,
$$
|f(x) - f(x_0)| = |x - x_0| |x^2 + x x_0 + x_0^2| < \left(\frac{\epsilon}{M}\right) M = \epsilon
$$
Ainsi, pour tout $x_0 \in \mathbb{R}$ et tout $\epsilon > 0$, nous avons trouvé un $\delta > 0$ tel que si $|x - x_0| < \delta$, alors $|f(x) - f(x_0)| < \epsilon$.
La fonction $f(x) = x^3$ est donc continue en tout point $x_0 \in \mathbb{R}$.

\subsubsection{Partie B : Continuité d'une Fonction Définie par Morceaux avec Paramètres}

Pour que la fonction $g(x)$ soit continue en $x=1$, trois conditions doivent être satisfaites :
1.  La fonction $g(x)$ doit être définie en $x=1$. C'est le cas, $g(1) = b$.
2.  La limite de $g(x)$ lorsque $x$ tend vers $1$ doit exister. Cela signifie que les limites à gauche et à droite doivent exister et être égales.
3.  La limite de $g(x)$ lorsque $x$ tend vers $1$ doit être égale à $g(1)$.

\textbf{Étape 1 : Calcul de $g(1)$}
Par définition de la fonction $g(x)$, nous avons :
$$
g(1) = b
$$

\textbf{Étape 2 : Calcul de la limite à gauche en $x=1$}
Pour $x < 1$, la fonction est définie par $g(x) = ax^2 + 3x - 1$.
Puisque $P(x) = ax^2 + 3x - 1$ est une fonction polynomiale, elle est continue sur $\mathbb{R}$. Par conséquent, la limite de $P(x)$ lorsque $x$ approche $1$ par la gauche est simplement l'évaluation de $P(x)$ en $x=1$.
$$
\lim_{x \to 1^-} g(x) = \lim_{x \to 1^-} (ax^2 + 3x - 1)
$$
Par substitution directe, en raison de la continuité des polynômes :
$$
\lim_{x \to 1^-} (ax^2 + 3x - 1) = a(1)^2 + 3(1) - 1 = a + 3 - 1 = a + 2
$$

\textbf{Étape 3 : Calcul de la limite à droite en $x=1$}
Pour $x > 1$, la fonction est définie par $g(x) = \frac{\sqrt{x+3}-2}{x-1}$.
Nous calculons la limite :
$$
\lim_{x \to 1^+} g(x) = \lim_{x \to 1^+} \frac{\sqrt{x+3}-2}{x-1}
$$
Si nous substituons $x=1$ directement, nous obtenons $\frac{\sqrt{1+3}-2}{1-1} = \frac{\sqrt{4}-2}{0} = \frac{2-2}{0} = \frac{0}{0}$, qui est une forme indéterminée.
Pour lever cette indétermination, nous allons multiplier le numérateur et le dénominateur par le conjugué du numérateur, qui est $\sqrt{x+3}+2$.
$$
\lim_{x \to 1^+} \frac{\sqrt{x+3}-2}{x-1} = \lim_{x \to 1^+} \frac{(\sqrt{x+3}-2)(\sqrt{x+3}+2)}{(x-1)(\sqrt{x+3}+2)}
$$
En utilisant l'identité $(A-B)(A+B) = A^2 - B^2$ pour le numérateur :
$$
= \lim_{x \to 1^+} \frac{(x+3) - 2^2}{(x-1)(\sqrt{x+3}+2)}
$$
$$
= \lim_{x \to 1^+} \frac{x+3 - 4}{(x-1)(\sqrt{x+3}+2)}
$$
$$
= \lim_{x \to 1^+} \frac{x-1}{(x-1)(\sqrt{x+3}+2)}
$$
Puisque nous considérons la limite lorsque $x \to 1^+$, $x$ est proche de $1$ mais $x \ne 1$. Par conséquent, $x-1 \ne 0$, et nous pouvons simplifier le terme $(x-1)$ au numérateur et au dénominateur :
$$
= \lim_{x \to 1^+} \frac{1}{\sqrt{x+3}+2}
$$
Maintenant, la substitution directe de $x=1$ ne conduit plus à une forme indéterminée. La fonction $h(x) = \sqrt{x+3}+2$ est continue pour $x \ge -3$. Puisque $x=1$ est dans ce domaine, nous pouvons substituer :
$$
= \frac{1}{\sqrt{1+3}+2} = \frac{1}{\sqrt{4}+2} = \frac{1}{2+2} = \frac{1}{4}
$$

\textbf{Étape 4 : Égalisation des limites et de la valeur de la fonction}
Pour que $g(x)$ soit continue en $x=1$, les trois valeurs doivent être égales :
$$
\lim_{x \to 1^-} g(x) = \lim_{x \to 1^+} g(x) = g(1)
$$
En substituant les valeurs calculées :
$$
a + 2 = \frac{1}{4} = b
$$
De l'égalité $b = \frac{1}{4}$, nous obtenons la valeur de $b$.
De l'égalité $a + 2 = \frac{1}{4}$, nous résolvons pour $a$ :
$$
a = \frac{1}{4} - 2
$$
Pour soustraire, nous mettons $2$ au même dénominateur que $\frac{1}{4}$ : $2 = \frac{8}{4}$.
$$
a = \frac{1}{4} - \frac{8}{4} = \frac{1 - 8}{4} = -\frac{7}{4}
$$
Ainsi, pour que la fonction $g(x)$ soit continue en $x=1$, les constantes $a$ et $b$ doivent prendre les valeurs suivantes :
$$
a = -\frac{7}{4} \quad \text{et} \quad b = \frac{1}{4}
$$

\subsubsection{Partie C : Application du Théorème des Valeurs Intermédiaires}

Nous voulons démontrer qu'il existe au moins un point $c \in [0,1]$ tel que $h(c) = c$.
Pour cela, nous allons définir une nouvelle fonction auxiliaire $k(x)$ et appliquer le Théorème des Valeurs Intermédiaires (TVI).

\textbf{Étape 1 : Définition de la fonction auxiliaire}
Soit la fonction $k: [0,1] \to \mathbb{R}$ définie par $k(x) = h(x) - x$.
Un point fixe $c$ de $h$ est une valeur pour laquelle $h(c) = c$, ce qui est équivalent à $h(c) - c = 0$, c'est-à-dire $k(c) = 0$. Notre objectif est donc de montrer l'existence d'une racine pour $k(x)$ dans l'intervalle $[0,1]$.

\textbf{Étape 2 : Vérification de la continuité de $k(x)$}
Pour appliquer le TVI, la fonction $k(x)$ doit être continue sur l'intervalle fermé $[0,1]$.
1.  La fonction $h(x)$ est continue sur $[0,1]$ par hypothèse de l'énoncé.
2.  La fonction $j(x) = x$ est une fonction polynomiale (spécifiquement, la fonction identité), et les fonctions polynomiales sont continues sur tout $\mathbb{R}$, et donc sur l'intervalle $[0,1]$.
3.  Le Théorème sur la somme/différence de fonctions continues stipule que si $f_1$ et $f_2$ sont continues sur un intervalle, alors $f_1 \pm f_2$ est aussi continue sur cet intervalle.
Puisque $h(x)$ et $j(x)=x$ sont continues sur $[0,1]$, leur différence $k(x) = h(x) - x$ est également continue sur $[0,1]$.

\textbf{Étape 3 : Évaluation de $k(x)$ aux bornes de l'intervalle}
Nous évaluons la fonction $k(x)$ aux extrémités de l'intervalle $[0,1]$ :
1.  Pour $x=0$:
    $k(0) = h(0) - 0 = h(0)$.
    Par hypothèse, la fonction $h$ a pour codomaine $[0,1]$, ce qui signifie que pour tout $x \in [0,1]$, $0 \le h(x) \le 1$.
    En particulier, pour $x=0$, nous avons $0 \le h(0) \le 1$.
    Donc, $k(0) = h(0) \ge 0$.

2.  Pour $x=1$:
    $k(1) = h(1) - 1$.
    De même, pour $x=1$, nous avons $0 \le h(1) \le 1$.
    En soustrayant $1$ de cette inégalité, nous obtenons :
    $0 - 1 \le h(1) - 1 \le 1 - 1$
    $-1 \le h(1) - 1 \le 0$.
    Donc, $k(1) = h(1) - 1 \le 0$.

\textbf{Étape 4 : Application du Théorème des Valeurs Intermédiaires}
Nous avons une fonction $k(x)$ qui est continue sur l'intervalle fermé $[0,1]$.
Nous avons également déterminé que $k(0) \ge 0$ et $k(1) \le 0$.
Le Théorème des Valeurs Intermédiaires (TVI) stipule que si une fonction $f$ est continue sur un intervalle fermé $[a,b]$, et si $y_0$ est un nombre quelconque entre $f(a)$ et $f(b)$ (inclusivement), alors il existe au moins un $c \in [a,b]$ tel que $f(c) = y_0$.

Considérons les cas possibles pour $k(0)$ et $k(1)$ :
*   \textbf{Cas 1 :} Si $k(0) = 0$.
    Alors $h(0) - 0 = 0$, ce qui signifie $h(0) = 0$. Dans ce cas, $c=0$ est un point fixe.
*   \textbf{Cas 2 :} Si $k(1) = 0$.
    Alors $h(1) - 1 = 0$, ce qui signifie $h(1) = 1$. Dans ce cas, $c=1$ est un point fixe.
*   \textbf{Cas 3 :} Si $k(0) > 0$ et $k(1) < 0$.
    Dans ce cas, $k(1) < 0 < k(0)$. Le nombre $0$ est strictement compris entre $k(1)$ et $k(0)$.
    Puisque $k(x)$ est continue sur $[0,1]$, le TVI garantit qu'il existe au moins un $c \in (0,1)$ tel que $k(c) = 0$.
    Par définition de $k(x)$, $k(c) = h(c) - c$. Donc, $h(c) - c = 0$, ce qui implique $h(c) = c$.

Dans tous les cas (que $k(0)$ ou $k(1)$ soit nul, ou que $0$ soit strictement entre eux), il existe au moins un $c \in [0,1]$ tel que $k(c) = 0$.
Par conséquent, il existe au moins un $c \in [0,1]$ tel que $h(c) = c$.
La fonction $h$ admet donc au moins un point fixe dans l'intervalle $[0,1]$.

\section{Exo-08}
\section{Exercice 8 - Continuité des fonctions d'une variable réelle : Analyse approfondie et applications}

Cet exercice explore la notion de continuité sous différents angles, allant de la démonstration formelle à l'application de théorèmes fondamentaux.

\subsection{Partie A : Démonstration formelle de continuité}

Soit la fonction $f: \mathbb{R} \to \mathbb{R}$ définie par $f(x) = x^3$.
Démontrer, en utilisant la définition formelle $(\epsilon, \delta)$ de la continuité, que $f$ est continue en tout point $x_0 \in \mathbb{R}$.

\subsection{Partie B : Continuité d'une fonction définie par morceaux}

Soit la fonction $g: \mathbb{R} \to \mathbb{R}$ définie par :
$$ g(x) = \begin{cases} \frac{\sin(ax)}{x} & \text{si } x < 0 \\ b & \text{si } x = 0 \\ \frac{e^{2x}-1}{x} & \text{si } x > 0 \end{cases} $$
Déterminer les valeurs des constantes réelles $a$ et $b$ pour que la fonction $g$ soit continue sur $\mathbb{R}$.

\subsection{Partie C : Application du Théorème des Valeurs Intermédiaires}

Soit la fonction $h: \mathbb{R} \to \mathbb{R}$ définie par $h(x) = x^3 - 3x + 1$.
1.  Montrer que l'équation $h(x) = 0$ possède au moins une solution dans l'intervalle $]-2, -1[$.
2.  Montrer que l'équation $h(x) = 0$ possède au moins une solution dans l'intervalle $]0, 1[$.
3.  Montrer que l'équation $h(x) = 0$ possède au moins une solution dans l'intervalle $]1, 2[$.
4.  En déduire le nombre exact de racines réelles de l'équation $h(x) = 0$.

\subsection{Partie D : Point fixe d'une fonction continue}

Soit $I = [a, b]$ un intervalle fermé et borné de $\mathbb{R}$, avec $a < b$. Soit $f: I \to I$ une fonction continue.
Démontrer qu'il existe au moins un point $c \in I$ tel que $f(c) = c$. (Un tel point $c$ est appelé un point fixe de $f$).

---

\section{Correction de l'Exercice 8}

\subsection{Partie A : Démonstration formelle de continuité}

Nous devons démontrer que la fonction $f(x) = x^3$ est continue en tout point $x_0 \in \mathbb{R}$ en utilisant la définition formelle $(\epsilon, \delta)$.

\textbf{Définition de la continuité en un point :}
Une fonction $f: \mathbb{R} \to \mathbb{R}$ est continue en un point $x_0 \in \mathbb{R}$ si et seulement si pour tout $\epsilon > 0$, il existe un $\delta > 0$ tel que pour tout $x \in \mathbb{R}$, si $|x - x_0| < \delta$, alors $|f(x) - f(x_0)| < \epsilon$.

\textbf{Démonstration :}
Soit $x_0 \in \mathbb{R}$ un point arbitraire.
Soit $\epsilon > 0$ un nombre réel strictement positif arbitraire. Nous cherchons à trouver un $\delta > 0$ tel que si $|x - x_0| < \delta$, alors $|f(x) - f(x_0)| < \epsilon$.

Nous commençons par analyser l'expression $|f(x) - f(x_0)|$:
$$ |f(x) - f(x_0)| = |x^3 - x_0^3| $$
Nous utilisons l'identité de factorisation de la différence de cubes, $A^3 - B^3 = (A - B)(A^2 + AB + B^2)$. Ici, $A=x$ et $B=x_0$.
$$ |x^3 - x_0^3| = |(x - x_0)(x^2 + x x_0 + x_0^2)| $$
Par la propriété du produit des valeurs absolues, $|AB| = |A||B|$ :
$$ |x^3 - x_0^3| = |x - x_0| |x^2 + x x_0 + x_0^2| $$
Notre objectif est de rendre cette expression inférieure à $\epsilon$. Nous avons déjà le facteur $|x - x_0|$, que nous pouvons rendre arbitrairement petit en choisissant $\delta$ petit. Le défi est de borner le second facteur, $|x^2 + x x_0 + x_0^2|$.

Pour borner ce facteur, nous allons d'abord imposer une condition sur $\delta$. Choisissons $\delta_1 = 1$.
Si $|x - x_0| < \delta_1 = 1$, alors nous avons :
$$ -1 < x - x_0 < 1 $$
En ajoutant $x_0$ à toutes les parties de l'inégalité, nous obtenons :
$$ x_0 - 1 < x < x_0 + 1 $$
Cela implique que $|x| < |x_0| + 1$. (Si $x_0 \ge 0$, $x < x_0+1 \implies |x| < x_0+1$. Si $x_0 < 0$, $x_0-1 < x \implies |x| < |x_0-1| = |x_0|+1$).

Maintenant, nous pouvons borner le facteur $|x^2 + x x_0 + x_0^2|$ en utilisant l'inégalité triangulaire, $|A+B+C| \le |A|+|B|+|C|$ :
$$ |x^2 + x x_0 + x_0^2| \le |x^2| + |x x_0| + |x_0^2| $$
$$ |x^2 + x x_0 + x_0^2| \le |x|^2 + |x||x_0| + |x_0|^2 $$
En utilisant l'estimation $|x| < |x_0| + 1$ :
$$ |x|^2 < (|x_0| + 1)^2 = |x_0|^2 + 2|x_0| + 1 $$
$$ |x||x_0| < (|x_0| + 1)|x_0| = |x_0|^2 + |x_0| $$
En substituant ces bornes dans l'inégalité :
$$ |x^2 + x x_0 + x_0^2| < (|x_0|^2 + 2|x_0| + 1) + (|x_0|^2 + |x_0|) + |x_0|^2 $$
$$ |x^2 + x x_0 + x_0^2| < 3|x_0|^2 + 3|x_0| + 1 $$
Soit $M = 3|x_0|^2 + 3|x_0| + 1$. $M$ est une constante positive qui dépend de $x_0$.

Maintenant, nous avons :
$$ |f(x) - f(x_0)| = |x - x_0| |x^2 + x x_0 + x_0^2| < |x - x_0| M $$
Nous voulons que $|f(x) - f(x_0)| < \epsilon$. Donc, nous voulons que $|x - x_0| M < \epsilon$.
Cela implique que nous devons choisir $|x - x_0| < \frac{\epsilon}{M}$.

Nous avons deux conditions sur $\delta$:
1.  $\delta \le 1$ (pour que la borne $M$ soit valide).
2.  $\delta \le \frac{\epsilon}{M}$ (pour que $|f(x) - f(x_0)| < \epsilon$).

Nous choisissons donc $\delta = \min\left(1, \frac{\epsilon}{M}\right)$. Puisque $M = 3|x_0|^2 + 3|x_0| + 1$ est toujours positif (et même strictement positif), $\frac{\epsilon}{M}$ est bien défini et positif.

\textbf{Récapitulatif de la preuve :}
Soit $x_0 \in \mathbb{R}$ et $\epsilon > 0$.
Posons $M = 3|x_0|^2 + 3|x_0| + 1$.
Choisissons $\delta = \min\left(1, \frac{\epsilon}{M}\right)$.
Alors, si $|x - x_0| < \delta$, nous avons deux conséquences :
1.  $|x - x_0| < 1$, ce qui implique $x_0 - 1 < x < x_0 + 1$. Par conséquent, $|x| < |x_0| + 1$.
2.  $|x - x_0| < \frac{\epsilon}{M}$.

Maintenant, nous évaluons $|f(x) - f(x_0)|$:
$$ |f(x) - f(x_0)| = |x^3 - x_0^3| $$
$$ |f(x) - f(x_0)| = |(x - x_0)(x^2 + x x_0 + x_0^2)| \quad \text{(par factorisation de la différence de cubes)} $$
$$ |f(x) - f(x_0)| = |x - x_0| |x^2 + x x_0 + x_0^2| \quad \text{(par propriété du produit des valeurs absolues)} $$
En utilisant l'inégalité triangulaire pour le second facteur :
$$ |x^2 + x x_0 + x_0^2| \le |x^2| + |x x_0| + |x_0^2| \quad \text{(par inégalité triangulaire)} $$
$$ |x^2 + x x_0 + x_0^2| \le |x|^2 + |x||x_0| + |x_0|^2 $$
Puisque $|x| < |x_0| + 1$ (dû au choix $\delta \le 1$) :
$$ |x|^2 < (|x_0| + 1)^2 = |x_0|^2 + 2|x_0| + 1 $$
$$ |x||x_0| < (|x_0| + 1)|x_0| = |x_0|^2 + |x_0| $$
En substituant ces bornes :
$$ |x^2 + x x_0 + x_0^2| < (|x_0|^2 + 2|x_0| + 1) + (|x_0|^2 + |x_0|) + |x_0|^2 $$
$$ |x^2 + x x_0 + x_0^2| < 3|x_0|^2 + 3|x_0| + 1 $$
Par définition, $M = 3|x_0|^2 + 3|x_0| + 1$. Donc :
$$ |x^2 + x x_0 + x_0^2| < M $$
En combinant avec $|x - x_0| < \frac{\epsilon}{M}$ :
$$ |f(x) - f(x_0)| < \left(\frac{\epsilon}{M}\right) M $$
$$ |f(x) - f(x_0)| < \epsilon $$
Nous avons ainsi montré que pour tout $\epsilon > 0$, il existe un $\delta > 0$ (spécifiquement $\delta = \min(1, \frac{\epsilon}{3|x_0|^2 + 3|x_0| + 1})$) tel que si $|x - x_0| < \delta$, alors $|f(x) - f(x_0)| < \epsilon$.
Par conséquent, la fonction $f(x) = x^3$ est continue en tout point $x_0 \in \mathbb{R}$.

\subsection{Partie B : Continuité d'une fonction définie par morceaux}

La fonction $g: \mathbb{R} \to \mathbb{R}$ est définie par :
$$ g(x) = \begin{cases} \frac{\sin(ax)}{x} & \text{si } x < 0 \\ b & \text{si } x = 0 \\ \frac{e^{2x}-1}{x} & \text{si } x > 0 \end{cases} $$
Pour que la fonction $g$ soit continue sur $\mathbb{R}$, elle doit être continue sur les intervalles ouverts $]-\infty, 0[$ et $]0, \infty[$, et elle doit être continue au point de raccordement $x=0$.

\textbf{1. Continuité sur les intervalles ouverts :}
*   \textbf{Pour $x < 0$ :} La fonction $x \mapsto \sin(ax)$ est une composition de fonctions continues ($x \mapsto ax$ et $u \mapsto \sin(u)$), donc elle est continue. La fonction $x \mapsto x$ est continue et non nulle sur $]-\infty, 0[$. Par conséquent, le quotient $\frac{\sin(ax)}{x}$ est continu sur $]-\infty, 0[$ en tant que quotient de fonctions continues dont le dénominateur ne s'annule pas.
*   \textbf{Pour $x > 0$ :} La fonction $x \mapsto e^{2x}-1$ est une composition de fonctions continues ($x \mapsto 2x$, $u \mapsto e^u$, et $v \mapsto v-1$), donc elle est continue. La fonction $x \mapsto x$ est continue et non nulle sur $]0, \infty[$. Par conséquent, le quotient $\frac{e^{2x}-1}{x}$ est continu sur $]0, \infty[$ en tant que quotient de fonctions continues dont le dénominateur ne s'annule pas.

\textbf{2. Continuité au point $x=0$ :}
Pour que $g$ soit continue en $x=0$, il faut que la valeur de la fonction en $x=0$ soit égale à la limite de la fonction lorsque $x$ tend vers $0$. C'est-à-dire :
$$ \lim_{x \to 0^-} g(x) = g(0) = \lim_{x \to 0^+} g(x) $$

*   \textbf{Calcul de $g(0)$ :}
    Par définition de la fonction $g$, $g(0) = b$.

*   \textbf{Calcul de la limite à gauche, $\lim_{x \to 0^-} g(x)$ :}
    Pour $x < 0$, $g(x) = \frac{\sin(ax)}{x}$.
    Nous devons calculer $\lim_{x \to 0^-} \frac{\sin(ax)}{x}$.
    Nous utilisons le résultat de limite remarquable $\lim_{u \to 0} \frac{\sin(u)}{u} = 1$.
    Si $a=0$, alors $\lim_{x \to 0^-} \frac{\sin(0)}{x} = \lim_{x \to 0^-} \frac{0}{x} = \lim_{x \to 0^-} 0 = 0$.
    Si $a \ne 0$, nous pouvons réécrire l'expression :
    $$ \lim_{x \to 0^-} \frac{\sin(ax)}{x} = \lim_{x \to 0^-} a \cdot \frac{\sin(ax)}{ax} $$
    Posons $u = ax$. Lorsque $x \to 0^-$, $u \to 0$ (car $a$ est une constante finie).
    $$ \lim_{x \to 0^-} a \cdot \frac{\sin(ax)}{ax} = a \cdot \lim_{u \to 0} \frac{\sin(u)}{u} = a \cdot 1 = a $$
    Donc, $\lim_{x \to 0^-} g(x) = a$.

*   \textbf{Calcul de la limite à droite, $\lim_{x \to 0^+} g(x)$ :}
    Pour $x > 0$, $g(x) = \frac{e^{2x}-1}{x}$.
    Nous devons calculer $\lim_{x \to 0^+} \frac{e^{2x}-1}{x}$.
    Nous utilisons le résultat de limite remarquable $\lim_{u \to 0} \frac{e^u-1}{u} = 1$.
    Nous pouvons réécrire l'expression :
    $$ \lim_{x \to 0^+} \frac{e^{2x}-1}{x} = \lim_{x \to 0^+} 2 \cdot \frac{e^{2x}-1}{2x} $$
    Posons $u = 2x$. Lorsque $x \to 0^+$, $u \to 0$.
    $$ \lim_{x \to 0^+} 2 \cdot \frac{e^{2x}-1}{2x} = 2 \cdot \lim_{u \to 0} \frac{e^u-1}{u} = 2 \cdot 1 = 2 $$
    Donc, $\lim_{x \to 0^+} g(x) = 2$.

\textbf{3. Détermination des constantes $a$ et $b$ :}
Pour que $g$ soit continue en $x=0$, nous devons avoir :
$$ \lim_{x \to 0^-} g(x) = g(0) = \lim_{x \to 0^+} g(x) $$
En substituant les valeurs calculées :
$$ a = b = 2 $$
Par conséquent, pour que la fonction $g$ soit continue sur $\mathbb{R}$, les constantes $a$ et $b$ doivent prendre les valeurs $a=2$ et $b=2$.

\subsection{Partie C : Application du Théorème des Valeurs Intermédiaires}

Soit la fonction $h: \mathbb{R} \to \mathbb{R}$ définie par $h(x) = x^3 - 3x + 1$.

\textbf{Propriété de $h(x)$ :}
La fonction $h(x)$ est un polynôme. Les fonctions polynomiales sont continues sur tout $\mathbb{R}$. Cette propriété est essentielle pour appliquer le Théorème des Valeurs Intermédiaires (TVI).

\textbf{Théorème des Valeurs Intermédiaires (TVI) :}
Si une fonction $f$ est continue sur un intervalle fermé $[a, b]$, alors pour toute valeur $k$ comprise entre $f(a)$ et $f(b)$ (c'est-à-dire $f(a) \le k \le f(b)$ ou $f(b) \le k \le f(a)$), il existe au moins un $c \in [a, b]$ tel que $f(c) = k$.
Dans notre cas, nous cherchons des solutions à $h(x)=0$, donc $k=0$. Le TVI garantit l'existence d'une racine si $h(a)$ et $h(b)$ sont de signes opposés.

\textbf{1. Solution dans l'intervalle $]-2, -1[$ :}
*   La fonction $h(x) = x^3 - 3x + 1$ est continue sur l'intervalle fermé $[-2, -1]$ car c'est un polynôme.
*   Évaluons $h(x)$ aux bornes de l'intervalle :
    *   $h(-2) = (-2)^3 - 3(-2) + 1 = -8 + 6 + 1 = -1$.
    *   $h(-1) = (-1)^3 - 3(-1) + 1 = -1 + 3 + 1 = 3$.
*   Nous observons que $h(-2) = -1 < 0$ et $h(-1) = 3 > 0$.
*   Puisque $h(-2)$ et $h(-1)$ sont de signes opposés, et que $h$ est continue sur $[-2, -1]$, le Théorème des Valeurs Intermédiaires garantit qu'il existe au moins une valeur $c_1 \in ]-2, -1[$ telle que $h(c_1) = 0$.
    Ainsi, l'équation $h(x)=0$ possède au moins une solution dans $]-2, -1[$.

\textbf{2. Solution dans l'intervalle $]0, 1[$ :}
*   La fonction $h(x) = x^3 - 3x + 1$ est continue sur l'intervalle fermé $[0, 1]$ car c'est un polynôme.
*   Évaluons $h(x)$ aux bornes de l'intervalle :
    *   $h(0) = (0)^3 - 3(0) + 1 = 0 - 0 + 1 = 1$.
    *   $h(1) = (1)^3 - 3(1) + 1 = 1 - 3 + 1 = -1$.
*   Nous observons que $h(0) = 1 > 0$ et $h(1) = -1 < 0$.
*   Puisque $h(0)$ et $h(1)$ sont de signes opposés, et que $h$ est continue sur $[0, 1]$, le Théorème des Valeurs Intermédiaires garantit qu'il existe au moins une valeur $c_2 \in ]0, 1[$ telle que $h(c_2) = 0$.
    Ainsi, l'équation $h(x)=0$ possède au moins une solution dans $]0, 1[$.

\textbf{3. Solution dans l'intervalle $]1, 2[$ :}
*   La fonction $h(x) = x^3 - 3x + 1$ est continue sur l'intervalle fermé $[1, 2]$ car c'est un polynôme.
*   Évaluons $h(x)$ aux bornes de l'intervalle :
    *   $h(1) = (1)^3 - 3(1) + 1 = 1 - 3 + 1 = -1$.
    *   $h(2) = (2)^3 - 3(2) + 1 = 8 - 6 + 1 = 3$.
*   Nous observons que $h(1) = -1 < 0$ et $h(2) = 3 > 0$.
*   Puisque $h(1)$ et $h(2)$ sont de signes opposés, et que $h$ est continue sur $[1, 2]$, le Théorème des Valeurs Intermédiaires garantit qu'il existe au moins une valeur $c_3 \in ]1, 2[$ telle que $h(c_3) = 0$.
    Ainsi, l'équation $h(x)=0$ possède au moins une solution dans $]1, 2[$.

\textbf{4. Nombre exact de racines réelles :}
Nous avons montré l'existence de trois racines $c_1, c_2, c_3$ dans les intervalles $]-2, -1[$, $]0, 1[$, et $]1, 2[$ respectivement. Ces trois intervalles sont disjoints :
*   $]-2, -1[ \cap ]0, 1[ = \emptyset$
*   $]-2, -1[ \cap ]1, 2[ = \emptyset$
*   $]0, 1[ \cap ]1, 2[ = \emptyset$
Cela signifie que les trois racines $c_1, c_2, c_3$ sont distinctes.

La fonction $h(x) = x^3 - 3x + 1$ est un polynôme de degré 3. Un théorème fondamental de l'algèbre stipule qu'un polynôme de degré $n$ a au plus $n$ racines réelles (ou complexes, mais ici nous nous intéressons aux réelles).
Puisque nous avons trouvé trois racines réelles distinctes pour un polynôme de degré 3, il ne peut pas y en avoir davantage.
Par conséquent, l'équation $h(x) = 0$ possède exactement trois racines réelles.

Pour une justification plus rigoureuse de l'unicité dans chaque intervalle, nous pouvons étudier la dérivée de $h(x)$:
$h'(x) = \frac{d}{dx}(x^3 - 3x + 1) = 3x^2 - 3$.
Factorisons $h'(x)$: $h'(x) = 3(x^2 - 1) = 3(x-1)(x+1)$.
Les racines de $h'(x)$ sont $x=-1$ et $x=1$. Ce sont les points critiques de $h(x)$.
*   Pour $x < -1$, $x-1 < 0$ et $x+1 < 0$, donc $h'(x) = 3(\text{négatif})(\text{négatif}) > 0$. $h$ est strictement croissante sur $]-\infty, -1]$.
*   Pour $-1 < x < 1$, $x-1 < 0$ et $x+1 > 0$, donc $h'(x) = 3(\text{négatif})(\text{positif}) < 0$. $h$ est strictement décroissante sur $[-1, 1]$.
*   Pour $x > 1$, $x-1 > 0$ et $x+1 > 0$, donc $h'(x) = 3(\text{positif})(\text{positif}) > 0$. $h$ est strictement croissante sur $[1, \infty[$.

Calculons les valeurs de $h$ aux points critiques :
$h(-1) = (-1)^3 - 3(-1) + 1 = -1 + 3 + 1 = 3$ (maximum local).
$h(1) = (1)^3 - 3(1) + 1 = 1 - 3 + 1 = -1$ (minimum local).

*   Sur $]-\infty, -1]$ : $h$ est strictement croissante. Puisque $h(-2) = -1$ et $h(-1) = 3$, et que $0 \in [-1, 3]$, il existe une unique racine dans $]-2, -1[$ par le corollaire du TVI (Théorème de la bijection).
*   Sur $[-1, 1]$ : $h$ est strictement décroissante. Puisque $h(-1) = 3$ et $h(1) = -1$, et que $0 \in [-1, 3]$, il existe une unique racine dans $]-1, 1[$. Cette racine est $c_2 \in ]0, 1[$ car $h(0)=1$ et $h(1)=-1$.
*   Sur $[1, \infty[$ : $h$ est strictement croissante. Puisque $h(1) = -1$ et $\lim_{x \to \infty} h(x) = \lim_{x \to \infty} x^3(1 - 3/x^2 + 1/x^3) = \infty$, et que $0 \in [-1, \infty[$, il existe une unique racine dans $]1, \infty[$. Cette racine est $c_3 \in ]1, 2[$ car $h(1)=-1$ et $h(2)=3$.

Cette analyse confirme l'existence et l'unicité de trois racines réelles distinctes.

\subsection{Partie D : Point fixe d'une fonction continue}

Soit $I = [a, b]$ un intervalle fermé et borné de $\mathbb{R}$, avec $a < b$. Soit $f: I \to I$ une fonction continue.
Nous voulons démontrer qu'il existe au moins un point $c \in I$ tel que $f(c) = c$.

\textbf{Stratégie :}
Nous allons définir une nouvelle fonction auxiliaire $k(x)$ et appliquer le Théorème des Valeurs Intermédiaires à cette fonction.

\textbf{1. Définition de la fonction auxiliaire :}
Considérons la fonction $k: I \to \mathbb{R}$ définie par $k(x) = f(x) - x$.
Un point $c$ est un point fixe de $f$ si $f(c) = c$, ce qui est équivalent à $f(c) - c = 0$, c'est-à-dire $k(c) = 0$.
Notre objectif est donc de montrer qu'il existe au moins un $c \in I$ tel que $k(c) = 0$.

\textbf{2. Continuité de $k(x)$ :}
*   La fonction $f$ est continue sur l'intervalle $I$ par hypothèse.
*   La fonction $x \mapsto x$ (la fonction identité) est continue sur tout $\mathbb{R}$, et donc sur l'intervalle $I$.
*   La fonction $k(x)$ est la différence de deux fonctions continues sur $I$. Par les propriétés des fonctions continues, la différence de deux fonctions continues est continue.
*   Par conséquent, $k(x)$ est continue sur l'intervalle fermé et borné $I = [a, b]$.

\textbf{3. Évaluation de $k(x)$ aux bornes de l'intervalle :}
*   \textbf{Pour $x=a$ :}
    $k(a) = f(a) - a$.
    Puisque la fonction $f$ a pour codomaine $I = [a, b]$ (c'est-à-dire $f: I \to I$), cela signifie que pour tout $x \in I$, $f(x) \in [a, b]$.
    En particulier, $f(a) \in [a, b]$. Cela implique $f(a) \ge a$.
    Donc, $f(a) - a \ge 0$.
    Ainsi, $k(a) \ge 0$.

*   \textbf{Pour $x=b$ :}
    $k(b) = f(b) - b$.
    De même, puisque $f(b) \in [a, b]$, cela implique $f(b) \le b$.
    Donc, $f(b) - b \le 0$.
    Ainsi, $k(b) \le 0$.

\textbf{4. Application du Théorème des Valeurs Intermédiaires :}
Nous avons les conditions suivantes pour la fonction $k(x)$ sur l'intervalle $[a, b]$ :
*   $k(x)$ est continue sur $[a, b]$.
*   $k(a) \ge 0$.
*   $k(b) \le 0$.

Nous considérons trois cas possibles :
*   \textbf{Cas 1 : $k(a) = 0$.}
    Si $k(a) = 0$, alors $f(a) - a = 0$, ce qui signifie $f(a) = a$. Dans ce cas, $a$ est un point fixe de $f$.
*   \textbf{Cas 2 : $k(b) = 0$.}
    Si $k(b) = 0$, alors $f(b) - b = 0$, ce qui signifie $f(b) = b$. Dans ce cas, $b$ est un point fixe de $f$.
*   \textbf{Cas 3 : $k(a) > 0$ et $k(b) < 0$.}
    Dans ce cas, nous avons $k(a) > 0$ et $k(b) < 0$. Puisque $k(x)$ est continue sur l'intervalle fermé $[a, b]$, et que $0$ est une valeur comprise entre $k(b)$ et $k(a)$ (car $k(b) < 0 < k(a)$), le Théorème des Valeurs Intermédiaires garantit qu'il existe au moins un point $c \in ]a, b[$ tel que $k(c) = 0$.
    Si $k(c) = 0$, alors $f(c) - c = 0$, ce qui signifie $f(c) = c$. Dans ce cas, $c$ est un point fixe de $f$.

Dans tous les cas (que $k(a)=0$, $k(b)=0$, ou $k(a)>0$ et $k(b)<0$), nous avons démontré l'existence d'au moins un point $c \in I = [a, b]$ tel que $f(c) = c$.
Ceci conclut la démonstration.

\section{Exo-09}
\section{Exercice 9 - Exploration approfondie de la continuité des fonctions d'une variable réelle}

Cet exercice est conçu pour évaluer votre maîtrise des concepts fondamentaux et avancés de la continuité des fonctions d'une variable réelle. Une rigueur absolue est attendue dans toutes les démonstrations.

---

\subsection{Énoncé de l'Exercice}

Soit $f: D \to \mathbb{R}$ une fonction d'une variable réelle.

\textbf{Partie A : Définition formelle de la continuité en un point.}
Soit la fonction $f: \mathbb{R} \to \mathbb{R}$ définie par $f(x) = x\textasciicircum 2 + 5x - 3$.
Démontrez, en utilisant la définition formelle $(\epsilon, \delta)$ de la continuité, que $f$ est continue au point $x_0 = 2$.

\textbf{Partie B : Continuité des fonctions définies par morceaux.}
Considérons la fonction $g: \mathbb{R} \to \mathbb{R}$ définie par :
$$
g(x) = \begin{cases}
    \frac{\sin(ax)}{x} & \text{si } x < 0 \\
    b & \text{si } x = 0 \\
    \frac{e^{2x}-1}{x} & \text{si } x > 0
\end{cases}
$$
Déterminez les valeurs des constantes réelles $a$ et $b$ pour lesquelles la fonction $g$ est continue en $x=0$.

\textbf{Partie C : Application du Théorème des Valeurs Intermédiaires.}
Soit $h: [0,1] \to [0,1]$ une fonction continue.
Démontrez qu'il existe au moins un point $c \in [0,1]$ tel que $h(c) = c$. (Un tel point est appelé point fixe de $h$).

\textbf{Partie D : Équation fonctionnelle de Cauchy et continuité.}
Soit $k: \mathbb{R} \to \mathbb{R}$ une fonction satisfaisant l'équation fonctionnelle de Cauchy :
$$
k(x+y) = k(x) + k(y) \quad \text{pour tout } x, y \in \mathbb{R}
$$
1.  Démontrez que $k(0) = 0$.
2.  Démontrez que $k(nx) = nk(x)$ pour tout $n \in \mathbb{Z}$ et tout $x \in \mathbb{R}$.
3.  Démontrez que $k(qx) = qk(x)$ pour tout $q \in \mathbb{Q}$ et tout $x \in \mathbb{R}$.
4.  Démontrez que si $k$ est continue en un point $x_0 \in \mathbb{R}$, alors $k$ est continue sur tout $\mathbb{R}$.
5.  Démontrez que si $k$ est continue sur $\mathbb{R}$, alors il existe une constante réelle $C$ telle que $k(x) = Cx$ pour tout $x \in \mathbb{R}$.

---

\subsection{Correction Détaillée}

\subsubsection{Partie A : Définition formelle de la continuité en un point.}

Pour démontrer que $f(x) = x\textasciicircum 2 + 5x - 3$ est continue en $x_0 = 2$ en utilisant la définition $(\epsilon, \delta)$, nous devons montrer que pour tout $\epsilon > 0$, il existe un $\delta > 0$ tel que si $|x - 2| < \delta$, alors $|f(x) - f(2)| < \epsilon$.

\textbf{Étape 1 : Calcul de $f(2)$.}
Nous calculons la valeur de la fonction au point $x_0 = 2$:
$f(2) = (2)^2 + 5(2) - 3 = 4 + 10 - 3 = 11$.

\textbf{Étape 2 : Analyse de l'expression $|f(x) - f(2)|$.}
Nous considérons l'expression $|f(x) - f(2)|$:
$$
|f(x) - f(2)| = |(x^2 + 5x - 3) - 11|
$$
$$
|f(x) - f(2)| = |x^2 + 5x - 14|
$$
Nous factorisons le polynôme $x^2 + 5x - 14$. Puisque $x=2$ est une racine (car $f(2)-f(2)=0$), $(x-2)$ doit être un facteur.
Par division polynomiale ou par inspection, nous trouvons :
$x^2 + 5x - 14 = (x-2)(x+7)$.
Donc,
$$
|f(x) - f(2)| = |(x-2)(x+7)| = |x-2||x+7|
$$

\textbf{Étape 3 : Majoration du terme $|x+7|$.}
Nous voulons majorer le terme $|x+7|$ en supposant que $x$ est "proche" de $2$.
Choisissons une première contrainte sur $\delta$, par exemple $\delta \le 1$.
Si $|x - 2| < \delta$ et $\delta \le 1$, alors $|x - 2| < 1$.
Ceci implique :
$-1 < x - 2 < 1$
En ajoutant $2$ à toutes les parties de l'inégalité :
$1 < x < 3$
Maintenant, nous voulons majorer $|x+7|$. En ajoutant $7$ à toutes les parties de l'inégalité $1 < x < 3$:
$1 + 7 < x + 7 < 3 + 7$
$8 < x + 7 < 10$
Puisque $x+7$ est compris entre $8$ et $10$, il est positif, et sa valeur absolue est $x+7$.
Ainsi, $|x+7| < 10$.

\textbf{Étape 4 : Détermination de $\delta$.}
En utilisant la majoration de l'étape 3, nous avons :
$|f(x) - f(2)| = |x-2||x+7| < |x-2| \cdot 10$.
Nous voulons que cette expression soit inférieure à $\epsilon$.
Nous posons donc $|x-2| \cdot 10 < \epsilon$.
Ceci implique $|x-2| < \frac{\epsilon}{10}$.
Nous avons deux conditions sur $\delta$: $\delta \le 1$ et $\delta \le \frac{\epsilon}{10}$.
Pour satisfaire les deux conditions simultanément, nous choisissons $\delta = \min\left(1, \frac{\epsilon}{10}\right)$.

\textbf{Étape 5 : Conclusion formelle.}
Soit $\epsilon > 0$.
Choisissons $\delta = \min\left(1, \frac{\epsilon}{10}\right)$.
Supposons que $|x - 2| < \delta$.
Puisque $\delta \le 1$, nous avons $|x - 2| < 1$.
Ceci implique $1 < x < 3$, et par conséquent $8 < x+7 < 10$.
Donc, $|x+7| < 10$.
Ensuite, nous utilisons la deuxième partie de la définition de $\delta$, $\delta \le \frac{\epsilon}{10}$.
Nous avons :
$$
|f(x) - f(2)| = |x-2||x+7|
$$
Puisque $|x-2| < \delta$ et $|x+7| < 10$ (sous la condition $\delta \le 1$), nous obtenons :
$$
|f(x) - f(2)| < \delta \cdot 10
$$
Puisque $\delta \le \frac{\epsilon}{10}$, nous avons :
$$
|f(x) - f(2)| < \left(\frac{\epsilon}{10}\right) \cdot 10
$$
$$
|f(x) - f(2)| < \epsilon
$$
Par conséquent, pour tout $\epsilon > 0$, il existe un $\delta = \min\left(1, \frac{\epsilon}{10}\right)$ tel que si $|x - 2| < \delta$, alors $|f(x) - f(2)| < \epsilon$.
La fonction $f(x) = x\textasciicircum 2 + 5x - 3$ est donc continue en $x_0 = 2$.

\subsubsection{Partie B : Continuité des fonctions définies par morceaux.}

Pour que la fonction $g$ soit continue en $x=0$, trois conditions doivent être satisfaites :
1.  La fonction $g$ doit être définie en $x=0$. C'est le cas, $g(0) = b$.
2.  La limite de $g(x)$ lorsque $x$ tend vers $0$ par la gauche doit exister.
3.  La limite de $g(x)$ lorsque $x$ tend vers $0$ par la droite doit exister.
4.  Ces trois valeurs doivent être égales : $\lim_{x \to 0^-} g(x) = \lim_{x \to 0^+} g(x) = g(0)$.

\textbf{Étape 1 : Calcul de $g(0)$.}
Par définition de la fonction $g$, nous avons $g(0) = b$.

\textbf{Étape 2 : Calcul de la limite à gauche $\lim_{x \to 0^-} g(x)$.}
Pour $x < 0$, $g(x) = \frac{\sin(ax)}{x}$.
Nous calculons la limite :
$$
\lim_{x \to 0^-} g(x) = \lim_{x \to 0^-} \frac{\sin(ax)}{x}
$$
Nous utilisons la limite remarquable $\lim_{u \to 0} \frac{\sin u}{u} = 1$.
Pour appliquer cette limite, nous multiplions et divisons par $a$ (en supposant $a \ne 0$ pour l'instant) :
$$
\lim_{x \to 0^-} \frac{\sin(ax)}{x} = \lim_{x \to 0^-} a \cdot \frac{\sin(ax)}{ax}
$$
En posant $u = ax$, lorsque $x \to 0^-$, $u \to 0^-$.
$$
\lim_{x \to 0^-} a \cdot \frac{\sin(ax)}{ax} = a \cdot \lim_{u \to 0^-} \frac{\sin u}{u}
$$
Puisque $\lim_{u \to 0^-} \frac{\sin u}{u} = 1$, nous obtenons :
$$
\lim_{x \to 0^-} g(x) = a \cdot 1 = a
$$
Si $a=0$, alors $\lim_{x \to 0^-} \frac{\sin(0x)}{x} = \lim_{x \to 0^-} \frac{0}{x} = 0$. La formule $a \cdot 1 = a$ reste valide.

\textbf{Étape 3 : Calcul de la limite à droite $\lim_{x \to 0^+} g(x)$.}
Pour $x > 0$, $g(x) = \frac{e^{2x}-1}{x}$.
Nous calculons la limite :
$$
\lim_{x \to 0^+} g(x) = \lim_{x \to 0^+} \frac{e^{2x}-1}{x}
$$
Nous utilisons la limite remarquable $\lim_{u \to 0} \frac{e^u-1}{u} = 1$.
Pour appliquer cette limite, nous multiplions et divisons par $2$ :
$$
\lim_{x \to 0^+} \frac{e^{2x}-1}{x} = \lim_{x \to 0^+} 2 \cdot \frac{e^{2x}-1}{2x}
$$
En posant $v = 2x$, lorsque $x \to 0^+$, $v \to 0^+$.
$$
\lim_{x \to 0^+} 2 \cdot \frac{e^{2x}-1}{2x} = 2 \cdot \lim_{v \to 0^+} \frac{e^v-1}{v}
$$
Puisque $\lim_{v \to 0^+} \frac{e^v-1}{v} = 1$, nous obtenons :
$$
\lim_{x \to 0^+} g(x) = 2 \cdot 1 = 2
$$

\textbf{Étape 4 : Égalisation des valeurs pour la continuité.}
Pour que $g$ soit continue en $x=0$, nous devons avoir :
$\lim_{x \to 0^-} g(x) = \lim_{x \to 0^+} g(x) = g(0)$.
En substituant les valeurs calculées :
$a = 2 = b$.
Par conséquent, pour que la fonction $g$ soit continue en $x=0$, les constantes $a$ et $b$ doivent prendre les valeurs $a=2$ et $b=2$.

\subsubsection{Partie C : Application du Théorème des Valeurs Intermédiaires.}

Nous voulons démontrer qu'il existe au moins un point $c \in [0,1]$ tel que $h(c) = c$, où $h: [0,1] \to [0,1]$ est une fonction continue.
Ceci est équivalent à montrer que l'équation $h(x) - x = 0$ a au moins une solution dans l'intervalle $[0,1]$.

\textbf{Étape 1 : Définition d'une fonction auxiliaire.}
Considérons la fonction auxiliaire $k: [0,1] \to \mathbb{R}$ définie par $k(x) = h(x) - x$.

\textbf{Étape 2 : Vérification de la continuité de $k$.}
La fonction $h$ est donnée comme continue sur l'intervalle $[0,1]$.
La fonction $x \mapsto x$ est une fonction polynomiale, et donc elle est continue sur tout $\mathbb{R}$, et en particulier sur $[0,1]$.
Puisque $k(x)$ est la différence de deux fonctions continues sur $[0,1]$ ($h(x)$ et $x$), $k(x)$ est également continue sur $[0,1]$.

\textbf{Étape 3 : Évaluation de $k$ aux bornes de l'intervalle.}
Nous évaluons la fonction $k$ aux extrémités de l'intervalle $[0,1]$ :
Pour $x=0$:
$k(0) = h(0) - 0 = h(0)$.
Puisque la fonction $h$ a pour codomaine $[0,1]$, nous savons que $h(0) \in [0,1]$.
Donc, $h(0) \ge 0$. Par conséquent, $k(0) \ge 0$.

Pour $x=1$:
$k(1) = h(1) - 1$.
Puisque la fonction $h$ a pour codomaine $[0,1]$, nous savons que $h(1) \in [0,1]$.
Donc, $h(1) \le 1$. Par conséquent, $h(1) - 1 \le 0$.
Donc, $k(1) \le 0$.

\textbf{Étape 4 : Application du Théorème des Valeurs Intermédiaires (TVI).}
Nous avons établi que :
*   $k$ est continue sur l'intervalle fermé et borné $[0,1]$.
*   $k(0) \ge 0$.
*   $k(1) \le 0$.
Nous avons donc $k(0) \cdot k(1) \le 0$.
Le Théorème des Valeurs Intermédiaires stipule que si une fonction est continue sur un intervalle fermé $[a,b]$ et si $k(a)$ et $k(b)$ ont des signes opposés (ou l'un est nul), alors il existe au moins un $c \in [a,b]$ tel que $k(c) = 0$.
Dans notre cas, $a=0$ et $b=1$.
Puisque $k$ est continue sur $[0,1]$ et que $k(0) \ge 0$ et $k(1) \le 0$, le TVI garantit l'existence d'au moins un $c \in [0,1]$ tel que $k(c) = 0$.

\textbf{Étape 5 : Conclusion.}
Puisque $k(c) = 0$, par définition de $k(x)$, nous avons $h(c) - c = 0$.
Ceci implique $h(c) = c$.
Nous avons donc démontré qu'il existe au moins un point $c \in [0,1]$ tel que $h(c) = c$.

\subsubsection{Partie D : Équation fonctionnelle de Cauchy et continuité.}

Soit $k: \mathbb{R} \to \mathbb{R}$ une fonction telle que $k(x+y) = k(x) + k(y)$ pour tout $x, y \in \mathbb{R}$.

\textbf{1. Démontrez que $k(0) = 0$.}
Nous utilisons la propriété de l'équation fonctionnelle en choisissant $x=0$ et $y=0$.
$k(0+0) = k(0) + k(0)$
$k(0) = 2k(0)$
En soustrayant $k(0)$ des deux côtés de l'équation :
$0 = k(0)$
Donc, $k(0) = 0$.

\textbf{2. Démontrez que $k(nx) = nk(x)$ pour tout $n \in \mathbb{Z}$ et tout $x \in \mathbb{R}$.}

*   \textbf{Cas $n \in \mathbb{N}^*$ (entiers positifs) :}
    Nous allons procéder par récurrence sur $n$.
    *   \textbf{Initialisation :} Pour $n=1$, $k(1x) = k(x)$ et $1k(x) = k(x)$. Donc $k(1x) = 1k(x)$ est vrai.
    *   \textbf{Hérédité :} Supposons que pour un certain entier $n \ge 1$, $k(nx) = nk(x)$ est vrai.
        Nous voulons montrer que $k((n+1)x) = (n+1)k(x)$.
        En utilisant l'équation fonctionnelle de Cauchy avec $y=nx$:
        $k((n+1)x) = k(x + nx) = k(x) + k(nx)$
        Par l'hypothèse de récurrence, $k(nx) = nk(x)$.
        Donc, $k((n+1)x) = k(x) + nk(x) = (1+n)k(x) = (n+1)k(x)$.
    *   \textbf{Conclusion :} Par le principe d'induction mathématique, $k(nx) = nk(x)$ pour tout $n \in \mathbb{N}^*$.

*   \textbf{Cas $n=0$ :}
    Nous avons $k(0x) = k(0)$. D'après la question précédente, $k(0)=0$.
    Et $0k(x) = 0$.
    Donc, $k(0x) = 0k(x)$ est vrai.

*   \textbf{Cas $n \in \mathbb{Z}^-$ (entiers négatifs) :}
    Soit $n$ un entier négatif. Alors $n = -m$ pour un certain entier positif $m \in \mathbb{N}^*$.
    Nous avons $k(x+(-x)) = k(x) + k(-x)$.
    Puisque $k(x+(-x)) = k(0)$, et nous savons que $k(0)=0$, nous avons :
    $0 = k(x) + k(-x)$
    Ceci implique $k(-x) = -k(x)$.
    Maintenant, considérons $k(nx) = k(-mx)$.
    $k(-mx) = -k(mx)$ (en utilisant $k(-z)=-k(z)$ avec $z=mx$).
    Puisque $m \in \mathbb{N}^*$, nous savons que $k(mx) = mk(x)$.
    Donc, $k(-mx) = -mk(x)$.
    Puisque $n = -m$, nous avons $-m = n$.
    Par conséquent, $k(nx) = nk(x)$ pour tout $n \in \mathbb{Z}^-$.

En combinant les trois cas, nous concluons que $k(nx) = nk(x)$ pour tout $n \in \mathbb{Z}$ et tout $x \in \mathbb{R}$.

\textbf{3. Démontrez que $k(qx) = qk(x)$ pour tout $q \in \mathbb{Q}$ et tout $x \in \mathbb{R}$.}
Soit $q \in \mathbb{Q}$. Par définition, $q$ peut s'écrire sous la forme $\frac{m}{n}$ où $m \in \mathbb{Z}$ et $n \in \mathbb{Z}^*$.
Nous voulons montrer que $k\left(\frac{m}{n}x\right) = \frac{m}{n}k(x)$.
Considérons $k(nx)$. D'après la question précédente, $k(nx) = nk(x)$.
Nous pouvons aussi écrire $nx = n \left(\frac{m}{n}x\right)$.
Donc, $k(nx) = k\left(n \cdot \left(\frac{m}{n}x\right)\right)$.
En utilisant la propriété $k(Ny) = Nk(y)$ pour $N \in \mathbb{Z}$ (ici $N=n$ et $y=\frac{m}{n}x$), nous avons :
$k\left(n \cdot \left(\frac{m}{n}x\right)\right) = n \cdot k\left(\frac{m}{n}x\right)$.
Ainsi, nous avons l'égalité :
$nk(x) = n \cdot k\left(\frac{m}{n}x\right)$.
Puisque $n \in \mathbb{Z}^*$, nous pouvons diviser par $n$ :
$k(x) = k\left(\frac{m}{n}x\right)$.
Attendez, il y a une erreur dans mon raisonnement. Reprenons.

Nous savons que $k(nx) = nk(x)$ pour $n \in \mathbb{Z}$.
Soit $q = \frac{m}{n}$ avec $m \in \mathbb{Z}$ et $n \in \mathbb{Z}^*$.
Nous voulons montrer $k(qx) = qk(x)$.
Considérons $k(x)$. Nous pouvons écrire $x = n \cdot \left(\frac{1}{n}x\right)$.
Donc $k(x) = k\left(n \cdot \left(\frac{1}{n}x\right)\right)$.
En utilisant la propriété $k(Ny) = Nk(y)$ pour $N=n$ et $y=\frac{1}{n}x$:
$k(x) = n \cdot k\left(\frac{1}{n}x\right)$.
En divisant par $n$ (puisque $n \ne 0$):
$\frac{1}{n}k(x) = k\left(\frac{1}{n}x\right)$.
Ceci montre que la propriété est vraie pour $q = \frac{1}{n}$.

Maintenant, nous voulons $k\left(\frac{m}{n}x\right)$.
Nous pouvons écrire $\frac{m}{n}x = m \cdot \left(\frac{1}{n}x\right)$.
En utilisant la propriété $k(Ny) = Nk(y)$ pour $N=m$ et $y=\frac{1}{n}x$:
$k\left(\frac{m}{n}x\right) = k\left(m \cdot \left(\frac{1}{n}x\right)\right) = m \cdot k\left(\frac{1}{n}x\right)$.
En substituant $k\left(\frac{1}{n}x\right) = \frac{1}{n}k(x)$ :
$k\left(\frac{m}{n}x\right) = m \cdot \left(\frac{1}{n}k(x)\right) = \frac{m}{n}k(x)$.
Donc, $k(qx) = qk(x)$ pour tout $q \in \mathbb{Q}$ et tout $x \in \mathbb{R}$.

\textbf{4. Démontrez que si $k$ est continue en un point $x_0 \in \mathbb{R}$, alors $k$ est continue sur tout $\mathbb{R}$.}
Supposons que $k$ est continue en $x_0$. Cela signifie que $\lim_{x \to x_0} k(x) = k(x_0)$.
Nous voulons montrer que $k$ est continue en un point arbitraire $a \in \mathbb{R}$.
C'est-à-dire, nous voulons montrer que $\lim_{x \to a} k(x) = k(a)$.

Considérons la limite $\lim_{h \to 0} k(h)$.
Puisque $k$ est continue en $x_0$, pour tout $\epsilon > 0$, il existe $\delta_0 > 0$ tel que si $|x - x_0| < \delta_0$, alors $|k(x) - k(x_0)| < \epsilon$.
Soit $h = x - x_0$. Alors $x = x_0 + h$. Lorsque $x \to x_0$, $h \to 0$.
L'inégalité $|x - x_0| < \delta_0$ devient $|h| < \delta_0$.
L'inégalité $|k(x) - k(x_0)| < \epsilon$ devient $|k(x_0+h) - k(x_0)| < \epsilon$.
En utilisant l'équation fonctionnelle de Cauchy, $k(x_0+h) = k(x_0) + k(h)$.
Donc, $|(k(x_0) + k(h)) - k(x_0)| < \epsilon$.
Ceci simplifie à $|k(h)| < \epsilon$.
Ainsi, pour tout $\epsilon > 0$, il existe $\delta_0 > 0$ tel que si $|h| < \delta_0$, alors $|k(h)| < \epsilon$.
Ceci est précisément la définition de $\lim_{h \to 0} k(h) = 0$.
Puisque $k(0)=0$ (démontré en D.1), cela signifie que $k$ est continue en $0$.

Maintenant, nous allons montrer que $k$ est continue en tout point $a \in \mathbb{R}$.
Nous devons montrer que $\lim_{x \to a} k(x) = k(a)$.
Considérons la limite $\lim_{h \to 0} k(a+h)$.
En utilisant l'équation fonctionnelle de Cauchy :
$k(a+h) = k(a) + k(h)$.
Donc,
$$
\lim_{x \to a} k(x) = \lim_{h \to 0} k(a+h)
$$
$$
\lim_{h \to 0} k(a+h) = \lim_{h \to 0} (k(a) + k(h))
$$
Par la propriété de la limite d'une somme (si les limites existent) :
$$
\lim_{h \to 0} (k(a) + k(h)) = \lim_{h \to 0} k(a) + \lim_{h \to 0} k(h)
$$
Puisque $k(a)$ est une constante par rapport à $h$, $\lim_{h \to 0} k(a) = k(a)$.
Nous avons démontré précédemment que $\lim_{h \to 0} k(h) = 0$.
Donc,
$$
\lim_{x \to a} k(x) = k(a) + 0 = k(a)
$$
Par conséquent, $k$ est continue en tout point $a \in \mathbb{R}$.

\textbf{5. Démontrez que si $k$ est continue sur $\mathbb{R}$, alors il existe une constante réelle $C$ telle que $k(x) = Cx$ pour tout $x \in \mathbb{R}$.}
Nous avons déjà démontré en D.3 que pour tout $q \in \mathbb{Q}$ et tout $x \in \mathbb{R}$, $k(qx) = qk(x)$.
Soit $x=1$. Alors $k(q \cdot 1) = qk(1)$ pour tout $q \in \mathbb{Q}$.
Posons $C = k(1)$. $C$ est une constante réelle.
Alors, pour tout nombre rationnel $q$, nous avons $k(q) = Cq$.

Maintenant, considérons un nombre réel arbitraire $x \in \mathbb{R}$.
Puisque les nombres rationnels sont denses dans $\mathbb{R}$, il existe une suite de nombres rationnels $(q_n)_{n \in \mathbb{N}}$ telle que $\lim_{n \to \infty} q_n = x$.
Puisque la fonction $k$ est continue sur $\mathbb{R}$ (par hypothèse de cette question), nous pouvons utiliser la propriété de continuité séquentielle :
$$
\lim_{n \to \infty} k(q_n) = k\left(\lim_{n \to \infty} q_n\right) = k(x)
$$
D'autre part, nous savons que pour chaque $q_n$ (qui est un nombre rationnel), $k(q_n) = Cq_n$.
Donc,
$$
\lim_{n \to \infty} k(q_n) = \lim_{n \to \infty} (Cq_n)
$$
Par la propriété de la limite d'un produit par une constante :
$$
\lim_{n \to \infty} (Cq_n) = C \cdot \lim_{n \to \infty} q_n
$$
Puisque $\lim_{n \to \infty} q_n = x$, nous avons :
$$
C \cdot \lim_{n \to \infty} q_n = Cx
$$
En combinant ces résultats, nous obtenons :
$$
k(x) = Cx
$$
Ceci est vrai pour tout $x \in \mathbb{R}$.
Par conséquent, si $k$ est continue sur $\mathbb{R}$, alors il existe une constante réelle $C$ (qui est $k(1)$) telle que $k(x) = Cx$ pour tout $x \in \mathbb{R}$.

\section{Exo-10}
\section{Exercice 10 - Exploration Approfondie de la Continuité des Fonctions Réelles}

Cet exercice est conçu pour tester votre compréhension approfondie et votre capacité à appliquer rigoureusement les concepts de continuité des fonctions d'une variable réelle. Chaque partie exige une démonstration formelle et détaillée, respectant la règle de Zéro Ellipse Mathématique.

---

\textbf{Partie A : Continuité en un point par la définition $\epsilon-\delta$.}

Soit la fonction $f: \mathbb{R} \to \mathbb{R}$ définie par :
$$ f(x) = \begin{cases} \frac{\sin(x^2)}{x} & \text{si } x \neq 0 \\ 0 & \text{si } x = 0 \end{cases} $$
Démontrez, en utilisant la définition formelle $(\epsilon-\delta)$ de la continuité, que la fonction $f$ est continue au point $x_0 = 0$.

---

\textbf{Partie B : Continuité uniforme.}

Soit la fonction $g: [0, +\infty) \to \mathbb{R}$ définie par $g(x) = \sqrt{x}$.
Démontrez, en utilisant la définition formelle de la continuité uniforme, que la fonction $g$ est uniformément continue sur l'intervalle $[0, +\infty)$.

---

\textbf{Partie C : Théorèmes des valeurs intermédiaires et des bornes atteintes.}

1.  Soit $h: [0,1] \to [0,1]$ une fonction continue. Démontrez qu'il existe au moins un point $c \in [0,1]$ tel que $h(c) = c$. (Ce point $c$ est appelé un point fixe de $h$).
2.  Soit $f: [a,b] \to \mathbb{R}$ une fonction continue, où $a, b \in \mathbb{R}$ avec $a < b$. Soient $y_1, y_2, \dots, y_n$ des points quelconques de l'intervalle $[a,b]$, où $n \in \mathbb{N}^*$. Démontrez qu'il existe un point $c \in [a,b]$ tel que $f(c) = \frac{f(y_1) + f(y_2) + \dots + f(y_n)}{n}$.

---

\textbf{Partie D : Équation fonctionnelle de Cauchy.}

Trouvez toutes les fonctions $f: \mathbb{R} \to \mathbb{R}$ qui sont continues et qui satisfont l'équation fonctionnelle suivante pour tous $x, y \in \mathbb{R}$ :
$$ f(x+y) = f(x) + f(y) $$
(Cette équation est connue sous le nom d'équation fonctionnelle de Cauchy).

---

\section{Correction de l'Exercice 10}

---

\textbf{Partie A : Continuité en un point par la définition $\epsilon-\delta$.}

Pour démontrer que la fonction $f$ est continue au point $x_0 = 0$, nous devons montrer que pour tout $\epsilon > 0$, il existe un $\delta > 0$ tel que pour tout $x \in \mathbb{R}$, si $|x - 0| < \delta$, alors $|f(x) - f(0)| < \epsilon$.
Puisque $f(0) = 0$, cette condition se simplifie à : pour tout $\epsilon > 0$, il existe un $\delta > 0$ tel que pour tout $x \in \mathbb{R}$, si $|x| < \delta$, alors $|f(x)| < \epsilon$.

Soit $\epsilon > 0$ un nombre réel arbitraire. Nous cherchons à déterminer un $\delta > 0$ approprié.

Considérons deux cas pour $x$:

1.  \textbf{Cas $x=0$ :}
    Dans ce cas, $|f(0) - f(0)| = |0 - 0| = 0$. Puisque $0 < \epsilon$ pour tout $\epsilon > 0$, la condition est satisfaite pour $x=0$.

2.  \textbf{Cas $x \neq 0$ :}
    Dans ce cas, $f(x) = \frac{\sin(x^2)}{x}$. Nous pouvons réécrire $f(x)$ en multipliant et divisant par $x$ (ce qui est permis car $x \neq 0$) :
    $$ f(x) = \frac{\sin(x^2)}{x^2} \cdot x $$
    Nous voulons que $|f(x)| < \epsilon$. En utilisant la propriété du module, nous avons :
    $$ |f(x)| = \left|\frac{\sin(x^2)}{x^2} \cdot x\right| = \left|\frac{\sin(x^2)}{x^2}\right| \cdot |x| $$
    Nous savons que $\lim_{u \to 0} \frac{\sin u}{u} = 1$.
    Par la définition de la limite, pour $\epsilon_0 = \frac{1}{2} > 0$, il existe un $\delta_0 > 0$ tel que si $0 < |u| < \delta_0$, alors $\left|\frac{\sin u}{u} - 1\right| < \frac{1}{2}$.
    Cette inégalité implique :
    $$ -\frac{1}{2} < \frac{\sin u}{u} - 1 < \frac{1}{2} $$
    En ajoutant $1$ à toutes les parties de l'inégalité, nous obtenons :
    $$ \frac{1}{2} < \frac{\sin u}{u} < \frac{3}{2} $$
    Par conséquent, pour $0 < |u| < \delta_0$, nous avons $\left|\frac{\sin u}{u}\right| < \frac{3}{2}$.

    Appliquons ceci avec $u = x^2$. Si $0 < |x^2| < \delta_0$, alors $\left|\frac{\sin(x^2)}{x^2}\right| < \frac{3}{2}$.
    La condition $0 < |x^2| < \delta_0$ est équivalente à $0 < |x| < \sqrt{\delta_0}$.
    Posons $\delta_1 = \sqrt{\delta_0}$.
    Ainsi, pour tout $x$ tel que $0 < |x| < \delta_1$, nous avons $\left|\frac{\sin(x^2)}{x^2}\right| < \frac{3}{2}$.

    Maintenant, nous pouvons majorer $|f(x)|$ pour $0 < |x| < \delta_1$:
    $$ |f(x)| = \left|\frac{\sin(x^2)}{x^2}\right| \cdot |x| < \frac{3}{2} |x| $$
    Nous voulons que cette expression soit inférieure à $\epsilon$. C'est-à-dire, nous voulons $\frac{3}{2} |x| < \epsilon$.
    Ceci est équivalent à $|x| < \frac{2\epsilon}{3}$.

    Pour satisfaire toutes les conditions, nous choisissons $\delta = \min\left(\delta_1, \frac{2\epsilon}{3}\right)$.
    Avec ce choix de $\delta$, si $|x| < \delta$:
    *   Si $x=0$, nous avons déjà montré que $|f(0)-f(0)| = 0 < \epsilon$.
    *   Si $x \neq 0$, alors $0 < |x| < \delta$.
        Puisque $\delta \le \delta_1$, nous avons $0 < |x| < \delta_1$, ce qui implique $\left|\frac{\sin(x^2)}{x^2}\right| < \frac{3}{2}$.
        Puisque $\delta \le \frac{2\epsilon}{3}$, nous avons $|x| < \frac{2\epsilon}{3}$.
        En combinant ces deux inégalités, nous obtenons :
        $$ |f(x) - f(0)| = |f(x)| = \left|\frac{\sin(x^2)}{x^2}\right| \cdot |x| < \frac{3}{2} \cdot \frac{2\epsilon}{3} = \epsilon $$
    Dans tous les cas, si $|x| < \delta$, alors $|f(x) - f(0)| < \epsilon$.
    Par conséquent, la fonction $f$ est continue au point $x_0 = 0$.

---

\textbf{Partie B : Continuité uniforme.}

Une fonction $g: I \to \mathbb{R}$ est uniformément continue sur un intervalle $I$ si pour tout $\epsilon > 0$, il existe un $\delta > 0$ tel que pour tous $x, y \in I$, si $|x-y| < \delta$, alors $|g(x) - g(y)| < \epsilon$.

Soit $\epsilon > 0$ un nombre réel arbitraire. Nous cherchons à déterminer un $\delta > 0$ tel que pour tous $x, y \in [0, +\infty)$, si $|x-y| < \delta$, alors $|\sqrt{x} - \sqrt{y}| < \epsilon$.

Considérons la quantité $|\sqrt{x} - \sqrt{y}|$.
Nous allons utiliser l'inégalité suivante : pour tous $x, y \ge 0$, $|\sqrt{x} - \sqrt{y}| \le \sqrt{|x-y|}$.
Démontrons cette inégalité :
Sans perte de généralité, supposons $x \ge y \ge 0$. Alors $\sqrt{x} \ge \sqrt{y} \ge 0$.
L'inégalité à prouver devient $\sqrt{x} - \sqrt{y} \le \sqrt{x-y}$.
Puisque les deux membres de l'inégalité sont non-négatifs, nous pouvons élever au carré sans changer le sens de l'inégalité :
$$ (\sqrt{x} - \sqrt{y})^2 \le (\sqrt{x-y})^2 $$
$$ x + y - 2\sqrt{xy} \le x-y $$
En soustrayant $x$ des deux côtés :
$$ y - 2\sqrt{xy} \le -y $$
En ajoutant $y$ aux deux côtés :
$$ 2y - 2\sqrt{xy} \le 0 $$
En divisant par $2$ :
$$ y - \sqrt{xy} \le 0 $$
$$ y \le \sqrt{xy} $$
Puisque $y \ge 0$, nous pouvons élever au carré les deux côtés :
$$ y^2 \le xy $$
En soustrayant $xy$ des deux côtés :
$$ y^2 - xy \le 0 $$
En factorisant $y$ :
$$ y(y-x) \le 0 $$
Puisque nous avons supposé $x \ge y$, il s'ensuit que $y-x \le 0$.
De plus, $y \ge 0$.
Le produit d'un nombre non-négatif ($y$) et d'un nombre non-positif ($y-x$) est toujours non-positif.
Donc, $y(y-x) \le 0$ est vrai.
L'inégalité $|\sqrt{x} - \sqrt{y}| \le \sqrt{|x-y|}$ est donc démontrée pour tous $x, y \ge 0$.

Maintenant, revenons à la preuve de la continuité uniforme.
Nous voulons que $|\sqrt{x} - \sqrt{y}| < \epsilon$.
En utilisant l'inégalité que nous venons de prouver, il suffit de s'assurer que $\sqrt{|x-y|} < \epsilon$.
Pour que $\sqrt{|x-y|} < \epsilon$, nous devons avoir $|x-y| < \epsilon^2$.

Nous choisissons donc $\delta = \epsilon^2$.
Soient $x, y \in [0, +\infty)$ tels que $|x-y| < \delta$.
Alors, par notre choix de $\delta$, nous avons $|x-y| < \epsilon^2$.
En prenant la racine carrée des deux côtés (les deux côtés sont non-négatifs) :
$$ \sqrt{|x-y|} < \sqrt{\epsilon^2} $$
$$ \sqrt{|x-y|} < \epsilon $$
En utilisant l'inégalité démontrée précédemment, $|\sqrt{x} - \sqrt{y}| \le \sqrt{|x-y|}$, nous obtenons :
$$ |\sqrt{x} - \sqrt{y}| < \epsilon $$
Ainsi, pour tout $\epsilon > 0$, nous avons trouvé un $\delta = \epsilon^2 > 0$ tel que pour tous $x, y \in [0, +\infty)$, si $|x-y| < \delta$, alors $|g(x) - g(y)| < \epsilon$.
Par conséquent, la fonction $g(x) = \sqrt{x}$ est uniformément continue sur l'intervalle $[0, +\infty)$.

---

\textbf{Partie C : Théorèmes des valeurs intermédiaires et des bornes atteintes.}

\textbf{1. Point fixe pour $h: [0,1] \to [0,1]$ continue.}

Soit $h: [0,1] \to [0,1]$ une fonction continue. Nous voulons démontrer qu'il existe au moins un point $c \in [0,1]$ tel que $h(c) = c$.
Pour cela, nous allons introduire une fonction auxiliaire $k(x)$ et appliquer le Théorème des Valeurs Intermédiaires (TVI).

Définissons la fonction $k: [0,1] \to \mathbb{R}$ par $k(x) = h(x) - x$.
*   \textbf{Continuité de $k(x)$ :} La fonction $h(x)$ est continue sur $[0,1]$ par hypothèse. La fonction $j(x) = x$ (la fonction identité) est continue sur $[0,1]$ car c'est un polynôme. La différence de deux fonctions continues est continue. Par conséquent, $k(x) = h(x) - x$ est continue sur l'intervalle fermé et borné $[0,1]$.

*   \textbf{Évaluation aux bornes de l'intervalle :}
    *   Calculons $k(0)$:
        $k(0) = h(0) - 0 = h(0)$.
        Puisque $h$ est une fonction de $[0,1]$ vers $[0,1]$, la valeur $h(0)$ doit être dans l'intervalle $[0,1]$.
        Donc, $h(0) \ge 0$. Par conséquent, $k(0) \ge 0$.
    *   Calculons $k(1)$:
        $k(1) = h(1) - 1$.
        Puisque $h$ est une fonction de $[0,1]$ vers $[0,1]$, la valeur $h(1)$ doit être dans l'intervalle $[0,1]$.
        Donc, $h(1) \le 1$. Par conséquent, $h(1) - 1 \le 0$. Donc, $k(1) \le 0$.

*   \textbf{Application du Théorème des Valeurs Intermédiaires (TVI) :}
    Nous avons $k(0) \ge 0$ et $k(1) \le 0$.
    Nous distinguons trois cas :
    1.  Si $k(0) = 0$, alors $h(0) - 0 = 0$, ce qui signifie $h(0) = 0$. Dans ce cas, $c=0$ est un point fixe.
    2.  Si $k(1) = 0$, alors $h(1) - 1 = 0$, ce qui signifie $h(1) = 1$. Dans ce cas, $c=1$ est un point fixe.
    3.  Si $k(0) > 0$ et $k(1) < 0$, alors $k(1) < 0 < k(0)$.
        Puisque $k$ est continue sur l'intervalle $[0,1]$ et que $0$ est une valeur comprise entre $k(1)$ et $k(0)$, le Théorème des Valeurs Intermédiaires garantit l'existence d'au moins un point $c \in (0,1)$ tel que $k(c) = 0$.
        La condition $k(c) = 0$ signifie $h(c) - c = 0$, ce qui est équivalent à $h(c) = c$.

Dans tous les cas, il existe au moins un point $c \in [0,1]$ tel que $h(c) = c$.
Ceci démontre l'existence d'un point fixe pour la fonction $h$.

\textbf{2. Généralisation du Théorème des Valeurs Intermédiaires.}

Soit $f: [a,b] \to \mathbb{R}$ une fonction continue. Soient $y_1, y_2, \dots, y_n$ des points quelconques de l'intervalle $[a,b]$, où $n \in \mathbb{N}^*$. Nous voulons démontrer qu'il existe un point $c \in [a,b]$ tel que $f(c) = \frac{f(y_1) + f(y_2) + \dots + f(y_n)}{n}$.

Soit $A$ la moyenne arithmétique des valeurs de la fonction aux points $y_i$:
$$ A = \frac{f(y_1) + f(y_2) + \dots + f(y_n)}{n} $$

*   \textbf{Application du Théorème des Bornes Atteintes (TBA) :}
    Puisque $f$ est continue sur l'intervalle fermé et borné $[a,b]$, le Théorème des Bornes Atteintes (aussi appelé Théorème de Weierstrass) garantit que $f$ atteint son minimum et son maximum sur cet intervalle.
    Soit $m = \min_{x \in [a,b]} f(x)$ et $M = \max_{x \in [a,b]} f(x)$.
    Il existe donc $x_{min} \in [a,b]$ tel que $f(x_{min}) = m$, et il existe $x_{max} \in [a,b]$ tel que $f(x_{max}) = M$.
    Pour tout $x \in [a,b]$, nous avons $m \le f(x) \le M$.

*   \textbf{Encadrement de la moyenne $A$ :}
    Puisque $y_1, y_2, \dots, y_n$ sont des points de l'intervalle $[a,b]$, nous avons pour chaque $i \in \{1, \dots, n\}$ :
    $$ m \le f(y_i) \le M $$
    En sommant ces $n$ inégalités, nous obtenons :
    $$ \sum_{i=1}^n m \le \sum_{i=1}^n f(y_i) \le \sum_{i=1}^n M $$
    $$ n \cdot m \le f(y_1) + f(y_2) + \dots + f(y_n) \le n \cdot M $$
    Puisque $n \in \mathbb{N}^*$, $n$ est un entier strictement positif. Nous pouvons diviser l'inégalité par $n$ sans changer son sens :
    $$ m \le \frac{f(y_1) + f(y_2) + \dots + f(y_n)}{n} \le M $$
    Donc, nous avons $m \le A \le M$.

*   \textbf{Application du Théorème des Valeurs Intermédiaires (TVI) :}
    Nous avons montré que $f(x_{min}) = m$ et $f(x_{max}) = M$.
    L'inégalité $m \le A \le M$ peut donc s'écrire $f(x_{min}) \le A \le f(x_{max})$.
    Puisque $f$ est continue sur l'intervalle $[a,b]$ (qui contient $x_{min}$ et $x_{max}$), et que $A$ est une valeur comprise entre $f(x_{min})$ et $f(x_{max})$, le Théorème des Valeurs Intermédiaires garantit l'existence d'au moins un point $c$ dans l'intervalle $[a,b]$ (plus précisément, entre $x_{min}$ et $x_{max}$, donc dans $[a,b]$) tel que $f(c) = A$.
    Par conséquent, il existe un point $c \in [a,b]$ tel que $f(c) = \frac{f(y_1) + f(y_2) + \dots + f(y_n)}{n}$.
    Ceci complète la démonstration.

---

\textbf{Partie D : Équation fonctionnelle de Cauchy.}

Nous cherchons toutes les fonctions $f: \mathbb{R} \to \mathbb{R}$ qui sont continues et qui satisfont l'équation fonctionnelle $f(x+y) = f(x) + f(y)$ pour tous $x, y \in \mathbb{R}$.

\textbf{Étape 1 : Déterminer $f(0)$.}
Substituons $x=0$ et $y=0$ dans l'équation fonctionnelle :
$f(0+0) = f(0) + f(0)$
$f(0) = 2f(0)$
En soustrayant $f(0)$ des deux côtés, nous obtenons :
$f(0) = 0$.

\textbf{Étape 2 : Déterminer $f(nx)$ pour $n \in \mathbb{N}$.}
*   Pour $n=1$: $f(1x) = f(x)$.
*   Pour $n=2$: $f(2x) = f(x+x) = f(x) + f(x) = 2f(x)$.
*   Pour $n=3$: $f(3x) = f(2x+x) = f(2x) + f(x) = 2f(x) + f(x) = 3f(x)$.
Nous allons prouver par induction que $f(nx) = nf(x)$ pour tout $n \in \mathbb{N}$ et tout $x \in \mathbb{R}$.
*   \textbf{Base de l'induction :} Pour $n=0$, $f(0x) = f(0) = 0$. Et $0f(x) = 0$. Donc $f(0x) = 0f(x)$ est vrai. Pour $n=1$, $f(1x) = f(x)$ et $1f(x) = f(x)$, donc $f(1x) = 1f(x)$ est vrai.
*   \textbf{Hypothèse d'induction :} Supposons que $f(kx) = kf(x)$ pour un certain entier naturel $k \ge 0$.
*   \textbf{Étape d'induction :} Nous voulons montrer que $f((k+1)x) = (k+1)f(x)$.
    $f((k+1)x) = f(kx + x)$
    En utilisant l'équation fonctionnelle avec $X=kx$ et $Y=x$:
    $f(kx + x) = f(kx) + f(x)$
    En utilisant l'hypothèse d'induction $f(kx) = kf(x)$:
    $f(kx) + f(x) = kf(x) + f(x)$
    En factorisant $f(x)$:
    $kf(x) + f(x) = (k+1)f(x)$.
    Donc, $f((k+1)x) = (k+1)f(x)$.
Par le principe d'induction mathématique, $f(nx) = nf(x)$ pour tout $n \in \mathbb{N}$ et tout $x \in \mathbb{R}$.

\textbf{Étape 3 : Déterminer $f(nx)$ pour $n \in \mathbb{Z}$.}
Nous avons $f(0) = 0$.
Pour $x \in \mathbb{R}$, $f(x + (-x)) = f(x) + f(-x)$.
Puisque $f(x + (-x)) = f(0) = 0$, nous avons $0 = f(x) + f(-x)$.
Ceci implique $f(-x) = -f(x)$ pour tout $x \in \mathbb{R}$.
Maintenant, si $n$ est un entier négatif, soit $n = -m$ où $m \in \mathbb{N}^*$.
Alors $f(nx) = f(-mx) = -f(mx)$ (en utilisant $f(-y) = -f(y)$ avec $y=mx$).
Puisque $m \in \mathbb{N}^*$, nous savons que $f(mx) = mf(x)$.
Donc, $-f(mx) = -mf(x)$.
Ainsi, $f(nx) = -mf(x) = nf(x)$ pour tout $n \in \mathbb{Z}$.

\textbf{Étape 4 : Déterminer $f(qx)$ pour $q \in \mathbb{Q}$.}
Soit $q \in \mathbb{Q}$. Alors $q$ peut s'écrire sous la forme $\frac{m}{n}$ où $m \in \mathbb{Z}$ et $n \in \mathbb{N}^*$.
Nous voulons montrer que $f(qx) = qf(x)$.
Considérons $f(nx)$. Nous savons que $f(nx) = nf(x)$.
Soit $y = qx = \frac{m}{n}x$. Alors $ny = nx \frac{m}{n} = mx$.
Donc, $f(ny) = f(mx)$.
En utilisant la propriété $f(kx) = kf(x)$ pour $k \in \mathbb{Z}$:
$nf(y) = mf(x)$.
En substituant $y = qx$:
$nf(qx) = mf(x)$.
En divisant par $n$ (ce qui est permis car $n \in \mathbb{N}^*$, donc $n \neq 0$):
$f(qx) = \frac{m}{n}f(x)$.
Puisque $q = \frac{m}{n}$, nous avons :
$f(qx) = qf(x)$ pour tout $q \in \mathbb{Q}$ et tout $x \in \mathbb{R}$.

\textbf{Étape 5 : Déterminer $f(x)$ pour $x \in \mathbb{R}$ en utilisant la continuité.}
Posons $x=1$ dans la relation $f(qx) = qf(x)$.
Alors $f(q) = qf(1)$ pour tout $q \in \mathbb{Q}$.
Soit $a = f(1)$. C'est une constante réelle.
Donc, $f(q) = aq$ pour tout $q \in \mathbb{Q}$.

Maintenant, nous utilisons la condition de continuité de $f$.
Soit $x \in \mathbb{R}$ un nombre réel arbitraire.
Nous savons qu'il existe une suite de nombres rationnels $(q_k)_{k \in \mathbb{N}}$ telle que $\lim_{k \to \infty} q_k = x$.
Puisque la fonction $f$ est continue sur $\mathbb{R}$, par la caractérisation séquentielle de la continuité, si $\lim_{k \to \infty} q_k = x$, alors $\lim_{k \to \infty} f(q_k) = f(x)$.
Nous savons que pour chaque $q_k \in \mathbb{Q}$, $f(q_k) = aq_k$.
Donc, $\lim_{k \to \infty} f(q_k) = \lim_{k \to \infty} (aq_k)$.
Par les propriétés des limites, la limite d'un produit est le produit des limites (si elles existent) :
$\lim_{k \to \infty} (aq_k) = a \lim_{k \to \infty} q_k$.
Puisque $\lim_{k \to \infty} q_k = x$, nous avons :
$a \lim_{k \to \infty} q_k = ax$.
En combinant ces résultats, nous obtenons :
$f(x) = ax$ pour tout $x \in \mathbb{R}$.

\textbf{Étape 6 : Vérifier la solution.}
Nous avons trouvé que les fonctions continues satisfaisant l'équation de Cauchy doivent être de la forme $f(x) = ax$ pour une constante $a \in \mathbb{R}$.
Vérifions si ces fonctions satisfont les conditions données :
1.  \textbf{Continuité :} La fonction $f(x) = ax$ est une fonction linéaire (un polynôme de degré 1 ou 0 si $a=0$). Les fonctions polynomiales sont continues sur tout $\mathbb{R}$. Donc, la condition de continuité est satisfaite.
2.  \textbf{Équation fonctionnelle :} Vérifions si $f(x+y) = f(x) + f(y)$ est satisfaite.
    Membre de gauche : $f(x+y) = a(x+y) = ax + ay$.
    Membre de droite : $f(x) + f(y) = ax + ay$.
    Puisque $ax+ay = ax+ay$, l'équation fonctionnelle est satisfaite.

\textbf{Conclusion :}
Les seules fonctions $f: \mathbb{R} \to \mathbb{R}$ qui sont continues et qui satisfont l'équation fonctionnelle de Cauchy $f(x+y) = f(x) + f(y)$ pour tous $x, y \in \mathbb{R}$ sont les fonctions de la forme $f(x) = ax$, où $a$ est une constante réelle arbitraire.


\chapter{Travaux pratiques et simulations algorithmiques}
\section{TP-01}
\section{TP 1 : Recherche de Zéros par la Méthode de Dichotomie (Bissection)}

\subsection{Introduction}

En tant que Professeur d'Élite et Développeur Senior, je vous propose ce premier Travail Pratique (TP) qui explore une application fondamentale de la continuité des fonctions d'une variable réelle : la recherche de zéros (ou racines). La continuité est une propriété essentielle qui garantit, entre autres, que les fonctions ne "sautent" pas de valeurs et qu'elles prennent toutes les valeurs intermédiaires entre deux points donnés. C'est précisément cette propriété qui est exploitée par la méthode de dichotomie, également connue sous le nom de méthode de bissection.

Ce TP vous guidera à travers l'implémentation "from scratch" de cette méthode en Python pur, sans l'aide de bibliothèques numériques avancées comme NumPy ou SciPy pour la logique algorithmique. L'objectif est de solidifier votre compréhension des principes mathématiques sous-jacents et de développer vos compétences en programmation pour traduire ces principes en code robuste et vérifiable.

\subsection{Principes Mathématiques : Le Théorème des Valeurs Intermédiaires et la Méthode de Dichotomie}

\subsubsection{Le Théorème des Valeurs Intermédiaires (TVI)}

Le cœur de la méthode de dichotomie repose sur le Théorème des Valeurs Intermédiaires (TVI). Ce théorème stipule que :

Si une fonction $f$ est continue sur un intervalle fermé $[a, b]$, alors pour toute valeur $k$ comprise entre $f(a)$ et $f(b)$ (c'est-à-dire $f(a) \le k \le f(b)$ ou $f(b) \le k \le f(a)$), il existe au moins un $c$ dans l'intervalle $[a, b]$ tel que $f(c) = k$.

Un corollaire direct du TVI, crucial pour la recherche de zéros, est le suivant :

Si une fonction $f$ est continue sur un intervalle fermé $[a, b]$ et si $f(a)$ et $f(b)$ ont des signes opposés (c'est-à-dire $f(a) \cdot f(b) < 0$), alors il existe au moins un $c$ dans l'intervalle ouvert $(a, b)$ tel que $f(c) = 0$.

\subsubsection{La Méthode de Dichotomie (Bissection)}

La méthode de dichotomie est un algorithme itératif qui utilise ce corollaire pour trouver une approximation d'un zéro d'une fonction continue. L'idée est simple :

1.  \textbf{Initialisation} : On commence avec un intervalle $[a, b]$ tel que $f(a)$ et $f(b)$ ont des signes opposés. Cela garantit, par le TVI, qu'il existe au moins un zéro dans cet intervalle.
2.  \textbf{Itération} :
    *   Calculer le milieu de l'intervalle : $c = (a + b) / 2$.
    *   Évaluer $f(c)$.
    *   Si $f(c)$ est suffisamment proche de zéro (selon une tolérance donnée), alors $c$ est notre approximation du zéro.
    *   Sinon, on réduit l'intervalle de recherche :
        *   Si $f(c)$ et $f(a)$ ont des signes opposés ($f(c) \cdot f(a) < 0$), le zéro se trouve dans $[a, c]$. On remplace $b$ par $c$.
        *   Si $f(c)$ et $f(b)$ ont des signes opposés ($f(c) \cdot f(b) < 0$), le zéro se trouve dans $[c, b]$. On remplace $a$ par $c$.
        *   Notez que si $f(c) = 0$ exactement, nous avons trouvé le zéro.
3.  \textbf{Critère d'arrêt} : L'algorithme s'arrête lorsque :
    *   La largeur de l'intervalle $[a, b]$ devient inférieure à une tolérance prédéfinie.
    *   La valeur absolue de $f(c)$ est inférieure à une tolérance prédéfinie.
    *   Un nombre maximal d'itérations est atteint pour éviter une boucle infinie.

Chaque itération divise la taille de l'intervalle par deux, garantissant une convergence linéaire vers le zéro. La méthode est robuste et toujours convergente sous les conditions initiales, bien que sa vitesse de convergence soit relativement lente comparée à d'autres méthodes (comme Newton-Raphson, qui nécessite la dérivée).

\subsection{Implémentation Python "From Scratch"}

Nous allons maintenant implémenter la méthode de dichotomie en Python. L'objectif est de créer une fonction \texttt{bisection\_method} qui prend en entrée la fonction $f$, l'intervalle $[a, b]$, une tolérance et un nombre maximal d'itérations.

\begin{lstlisting}[language=python]
import math

def bisection_method(f, a, b, tol=1e-7, max_iter=100):
    """
    Implemente la methode de dichotomie (bissection) pour trouver un zero
    d'une fonction continue f sur l'intervalle [a, b].

    Args:
        f (callable): La fonction dont on cherche un zero.
                      Doit etre continue sur [a, b].
        a (float): Borne inferieure de l'intervalle de recherche.
        b (float): Borne superieure de l'intervalle de recherche.
        tol (float): Tolerance pour la convergence. L'algorithme s'arrete
                     lorsque l'intervalle [a, b] est plus petit que tol,
                     ou lorsque |f(c)| est plus petit que tol.
        max_iter (int): Nombre maximal d'iterations pour eviter une boucle infinie.

    Returns:
        float: Une approximation du zero de f dans l'intervalle [a, b].

    Raises:
        ValueError: Si f(a) et f(b) n'ont pas de signes opposes,
                    ce qui ne garantit pas l'existence d'un zero par le TVI.
        RuntimeError: Si le nombre maximal d'iterations est atteint sans convergence.
    """

    # Verification de la condition du Theoreme des Valeurs Intermediaires
    fa = f(a)
    fb = f(b)

    if fa * fb > 0:
        raise ValueError(
            "Les valeurs de la fonction aux bornes de l'intervalle "
            "doivent avoir des signes opposes (f(a)*f(b) < 0) pour garantir "
            "l'existence d'un zero par le TVI."
        )

    # Cas ou l'une des bornes est deja un zero
    if abs(fa) < tol:
        return a
    if abs(fb) < tol:
        return b

    # Boucle d'iteration de la methode de dichotomie
    for i in range(max_iter):
        c = (a + b) / 2
        fc = f(c)

        # Critere d'arret base sur la valeur de la fonction
        if abs(fc) < tol:
            return c

        # Reduction de l'intervalle
        # On met a jour fa ou fb pour eviter de recalculer f(a) ou f(b)
        # a chaque iteration si l'intervalle ne change pas de cote.
        if fa * fc < 0:  # Le zero est dans [a, c]
            b = c
            fb = fc
        else:  # Le zero est dans [c, b]
            a = c
            fa = fc

        # Critere d'arret base sur la largeur de l'intervalle
        if (b - a) < tol:
            return (a + b) / 2 # Retourne le milieu du dernier intervalle

    # Si le nombre maximal d'iterations est atteint sans convergence
    raise RuntimeError(
        f"La methode de dichotomie n'a pas converge apres {max_iter} iterations. "
        "Considerez d'augmenter max_iter ou tol."
    )

\end{lstlisting}

\subsection{Exemples et Validation Mathématique avec Assertions}

Nous allons maintenant tester notre implémentation avec plusieurs fonctions et utiliser des assertions (\texttt{assert}) pour valider mathématiquement la logique. Les assertions sont cruciales pour garantir que notre code respecte les propriétés attendues du résultat.

\subsubsection{Exemple 1 : Fonction Polynomiale Simple $f(x) = x\textasciicircum 2 - 2$}

Nous cherchons les zéros de $f(x) = x\textasciicircum 2 - 2$. Les zéros exacts sont $\pm\sqrt{2}$. Nous allons chercher $\sqrt{2}$ dans l'intervalle $[1, 2]$.

\begin{lstlisting}[language=python]
# --- Exemple 1: f(x) = x\textasciicircum 2 - 2 ---
print("--- Exemple 1: f(x) = x\textasciicircum 2 - 2 ---")

def f1(x):
    return x**2 - 2

# Intervalle [1, 2] ou f(1) = -1 et f(2) = 2, donc un zero existe.
a1, b1 = 1.0, 2.0
tol1 = 1e-9

print(f"Recherche du zero de f(x) = x\textasciicircum 2 - 2 dans [{a1}, {b1}] avec tol={tol1}")

root1 = bisection_method(f1, a1, b1, tol=tol1)
print(f"Zero trouve : {root1}")
print(f"f({root1}) = {f1(root1)}")

# Validation mathematique avec assertions
# 1. La valeur de la fonction au zero trouve doit etre tres proche de zero.
assert abs(f1(root1)) < tol1, f"Assertion Failed: f(root1) = {f1(root1)} n'est pas proche de zero."
print(f"Assertion OK: |f({root1})| < {tol1}")

# 2. Le zero trouve doit etre dans l'intervalle initial.
assert a1 <= root1 <= b1, f"Assertion Failed: Le zero {root1} n'est pas dans [{a1}, {b1}]."
print(f"Assertion OK: {root1} est dans [{a1}, {b1}]")

# 3. Comparaison avec la valeur exacte (si connue).
expected_root1 = math.sqrt(2)
assert abs(root1 - expected_root1) < tol1, f"Assertion Failed: {root1} n'est pas assez proche de {expected_root1}."
print(f"Assertion OK: {root1} est proche de la valeur exacte {expected_root1}")

print("-" * 30)
\end{lstlisting}

\subsubsection{Exemple 2 : Fonction Transcendantale $f(x) = \cos(x) - x$}

Nous cherchons le zéro de $f(x) = \cos(x) - x$. Ce zéro est également le point fixe de $g(x) = \cos(x)$. Nous allons chercher dans l'intervalle $[0, \pi/2]$.

\begin{lstlisting}[language=python]
# --- Exemple 2: f(x) = cos(x) - x ---
print("--- Exemple 2: f(x) = cos(x) - x ---")

def f2(x):
    return math.cos(x) - x

# Intervalle [0, pi/2] ou f(0) = 1 et f(pi/2) = -pi/2, donc un zero existe.
a2, b2 = 0.0, math.pi / 2
tol2 = 1e-8

print(f"Recherche du zero de f(x) = cos(x) - x dans [{a2}, {b2}] avec tol={tol2}")

root2 = bisection_method(f2, a2, b2, tol=tol2)
print(f"Zero trouve : {root2}")
print(f"f({root2}) = {f2(root2)}")

# Validation mathematique avec assertions
assert abs(f2(root2)) < tol2, f"Assertion Failed: f(root2) = {f2(root2)} n'est pas proche de zero."
print(f"Assertion OK: |f({root2})| < {tol2}")

assert a2 <= root2 <= b2, f"Assertion Failed: Le zero {root2} n'est pas dans [{a2}, {b2}]."
print(f"Assertion OK: {root2} est dans [{a2}, {b2}]")

# Pour cette fonction, il n'y a pas de forme analytique simple pour le zero,
# donc nous nous fions principalement a f(root) ~ 0.
# On peut aussi verifier que le point est un point fixe de cos(x)
assert abs(math.cos(root2) - root2) < tol2, f"Assertion Failed: {root2} n'est pas un point fixe de cos(x)."
print(f"Assertion OK: {root2} est un point fixe de cos(x)")

print("-" * 30)
\end{lstlisting}

\subsubsection{Exemple 3 : Fonction avec plusieurs zéros, ciblant un intervalle spécifique $f(x) = x^3 - x - 1$}

Cette fonction a une seule racine réelle (la constante de plastique). Nous allons la chercher dans l'intervalle $[1, 2]$.

\begin{lstlisting}[language=python]
# --- Exemple 3: f(x) = x^3 - x - 1 ---
print("--- Exemple 3: f(x) = x^3 - x - 1 ---")

def f3(x):
    return x**3 - x - 1

# Intervalle [1, 2] ou f(1) = -1 et f(2) = 5, donc un zero existe.
a3, b3 = 1.0, 2.0
tol3 = 1e-9

print(f"Recherche du zero de f(x) = x^3 - x - 1 dans [{a3}, {b3}] avec tol={tol3}")

root3 = bisection_method(f3, a3, b3, tol=tol3)
print(f"Zero trouve : {root3}")
print(f"f({root3}) = {f3(root3)}")

# Validation mathematique avec assertions
assert abs(f3(root3)) < tol3, f"Assertion Failed: f(root3) = {f3(root3)} n'est pas proche de zero."
print(f"Assertion OK: |f({root3})| < {tol3}")

assert a3 <= root3 <= b3, f"Assertion Failed: Le zero {root3} n'est pas dans [{a3}, {b3}]."
print(f"Assertion OK: {root3} est dans [{a3}, {b3}]")

print("-" * 30)
\end{lstlisting}

\subsubsection{Exemple 4 : Cas où une borne est déjà le zéro}

Testons le cas où $f(a)=0$ ou $f(b)=0$.

\begin{lstlisting}[language=python]
# --- Exemple 4: f(x) = x - 5, intervalle [5, 10] ---
print("--- Exemple 4: f(x) = x - 5, borne est le zero ---")

def f4(x):
    return x - 5

a4, b4 = 5.0, 10.0
tol4 = 1e-9

print(f"Recherche du zero de f(x) = x - 5 dans [{a4}, {b4}] avec tol={tol4}")

root4 = bisection_method(f4, a4, b4, tol=tol4)
print(f"Zero trouve : {root4}")
print(f"f({root4}) = {f4(root4)}")

# Validation mathematique avec assertions
assert abs(f4(root4)) < tol4, f"Assertion Failed: f(root4) = {f4(root4)} n'est pas proche de zero."
print(f"Assertion OK: |f({root4})| < {tol4}")

assert a4 <= root4 <= b4, f"Assertion Failed: Le zero {root4} n'est pas dans [{a4}, {b4}]."
print(f"Assertion OK: {root4} est dans [{a4}, {b4}]")

assert abs(root4 - 5.0) < tol4, f"Assertion Failed: {root4} n'est pas assez proche de 5.0."
print(f"Assertion OK: {root4} est proche de la valeur exacte 5.0")

print("-" * 30)
\end{lstlisting}

\subsubsection{Exemple 5 : Gestion d'erreur - Pas de changement de signe}

Testons le cas où la condition du TVI n'est pas remplie ($f(a) \cdot f(b) > 0$).

\begin{lstlisting}[language=python]
# --- Exemple 5: Gestion d'erreur - Pas de changement de signe ---
print("--- Exemple 5: Gestion d'erreur - Pas de changement de signe ---")

def f5(x):
    return x**2 + 1 # Toujours positif

a5, b5 = -2.0, 2.0
tol5 = 1e-9

print(f"Tentative de recherche du zero de f(x) = x\textasciicircum 2 + 1 dans [{a5}, {b5}]")

try:
    root5 = bisection_method(f5, a5, b5, tol=tol5)
    # Cette ligne ne devrait pas etre atteinte
    assert False, "Assertion Failed: ValueError n'a pas ete levee pour f(a)*f(b) > 0."
except ValueError as e:
    print(f"Erreur attendue levee : {e}")
    print("Assertion OK: ValueError a ete correctement levee.")
except Exception as e:
    assert False, f"Assertion Failed: Une erreur inattendue a ete levee : {e}"

print("-" * 30)
\end{lstlisting}

\subsubsection{Exemple 6 : Gestion d'erreur - Nombre maximal d'itérations atteint}

Testons le cas où la convergence n'est pas atteinte dans le nombre d'itérations spécifié.

\begin{lstlisting}[language=python]
# --- Exemple 6: Gestion d'erreur - Nombre maximal d'iterations atteint ---
print("--- Exemple 6: Gestion d'erreur - Nombre maximal d'iterations atteint ---")

def f6(x):
    return x - 0.1 # Zero a 0.1

a6, b6 = 0.0, 1.0
tol6 = 1e-15 # Tres petite tolerance
max_iter6 = 5 # Tres peu d'iterations

print(f"Tentative de recherche du zero de f(x) = x - 0.1 dans [{a6}, {b6}] avec tol={tol6} et max_iter={max_iter6}")

try:
    root6 = bisection_method(f6, a6, b6, tol=tol6, max_iter=max_iter6)
    # Cette ligne ne devrait pas etre atteinte si max_iter est trop petit
    assert False, "Assertion Failed: RuntimeError n'a pas ete levee pour max_iter atteint."
except RuntimeError as e:
    print(f"Erreur attendue levee : {e}")
    print("Assertion OK: RuntimeError a ete correctement levee.")
except Exception as e:
    assert False, f"Assertion Failed: Une erreur inattendue a ete levee : {e}"

print("-" * 30)
\end{lstlisting}

\subsection{Analyse de la Convergence et Limitations}

La méthode de dichotomie est une méthode de recherche de zéros très robuste, car elle garantit toujours la convergence vers un zéro si les conditions initiales ($f(a) \cdot f(b) < 0$) sont remplies et si la fonction est continue.

\subsubsection{Vitesse de Convergence}

La convergence de la méthode de dichotomie est linéaire. À chaque itération, la taille de l'intervalle de recherche est divisée par deux. Si l'intervalle initial est de longueur $L_0 = b - a$, après $n$ itérations, la longueur de l'intervalle $L_n$ sera $L_n = L_0 / 2^n$. Pour atteindre une tolérance $\epsilon$, il faut environ $\log_2(L_0 / \epsilon)$ itérations. C'est une convergence prévisible mais relativement lente.

\subsubsection{Limitations}

1.  \textbf{Nécessite un changement de signe} : La méthode ne peut trouver un zéro que si la fonction change de signe sur l'intervalle initial. Elle ne peut pas trouver de zéros où la fonction touche l'axe des x sans le traverser (par exemple, le zéro de $f(x) = x\textasciicircum 2$ à $x=0$).
2.  \textbf{Lenteur} : Comme mentionné, sa convergence est linéaire, ce qui peut être lent pour atteindre une très haute précision par rapport à des méthodes d'ordre supérieur (comme la méthode de Newton, qui est quadratique).
3.  \textbf{Ne trouve qu'un seul zéro} : Si plusieurs zéros existent dans l'intervalle initial, la méthode de dichotomie ne trouvera qu'un seul d'entre eux. Pour trouver tous les zéros, il faudrait diviser l'intervalle en sous-intervalles plus petits et appliquer la méthode à chacun.
4.  \textbf{Pas de dérivée requise} : C'est aussi un avantage, car d'autres méthodes qui utilisent la dérivée (comme Newton) peuvent converger plus rapidement mais nécessitent cette information supplémentaire.

\subsection{Conclusion}

Ce TP vous a permis d'implémenter la méthode de dichotomie, une technique fondamentale pour la recherche de zéros de fonctions continues. Vous avez pu constater comment le Théorème des Valeurs Intermédiaires, un pilier de la théorie de la continuité, est directement appliqué pour garantir la convergence de l'algorithme. L'implémentation "from scratch" en Python, accompagnée de validations rigoureuses par assertions, renforce non seulement votre compréhension mathématique mais aussi votre capacité à écrire du code fiable et vérifiable.

La méthode de dichotomie, bien que simple et robuste, illustre parfaitement l'importance de la continuité pour résoudre des problèmes numériques. Elle sert de base à des algorithmes plus avancés et reste un outil précieux dans l'arsenal de tout ingénieur ou scientifique.

\section{TP-02}
\section{TP 2/5 : Théorème du Point Fixe et Méthodes Itératives}

\subsection{Introduction}

Dans le cadre de notre exploration approfondie de la continuité des fonctions d'une variable réelle, ce deuxième travail pratique se concentre sur une application fondamentale : le Théorème du Point Fixe. Ce théorème, qui trouve ses racines dans les propriétés topologiques des fonctions continues, est non seulement d'une grande élégance mathématique mais aussi d'une utilité pratique considérable, notamment dans la résolution numérique d'équations.

Un point fixe d'une fonction $f$ est une valeur $x^*$ telle que $f(x^*) = x^*$. La recherche de points fixes est équivalente à la résolution de l'équation $f(x) - x = 0$. De nombreuses équations peuvent être reformulées sous cette forme, rendant les méthodes de point fixe essentielles en analyse numérique. Ce TP vous guidera à travers l'implémentation "from scratch" d'une méthode itérative pour trouver ces points fixes, en s'appuyant sur les principes de la continuité et de la convergence.

\subsection{Objectifs Pédagogiques}

À l'issue de ce TP, vous devriez être capable de :

1.  \textbf{Comprendre le Théorème du Point Fixe} : Saisir les conditions d'existence et d'unicité d'un point fixe pour une fonction continue.
2.  \textbf{Implémenter une Méthode Itérative} : Développer un algorithme en Python pur pour approximer un point fixe.
3.  \textbf{Analyser la Convergence} : Identifier les facteurs qui influencent la convergence de la méthode itérative du point fixe, notamment le rôle de la dérivée de la fonction.
4.  \textbf{Valider Numériquement} : Utiliser des assertions pour vérifier la correction mathématique de l'implémentation.
5.  \textbf{Maîtriser Python "from scratch"} : Renforcer vos compétences en programmation sans recourir à des bibliothèques numériques avancées pour la logique de base.


\subsubsection{1.1 Définition d'un Point Fixe}

Soit $f: I \to \mathbb{R}$ une fonction définie sur un intervalle $I$. Un point $x^* \in I$ est appelé un \textbf{point fixe} de $f$ si $f(x^*) = x^*$.

Géométriquement, les points fixes sont les abscisses des points d'intersection de la courbe $y = f(x)$ avec la droite $y = x$.

\subsubsection{1.2 Théorème du Point Fixe (Version Simple)}

Un théorème fondamental, directement lié à la continuité et au Théorème des Valeurs Intermédiaires (TVI), garantit l'existence d'un point fixe sous certaines conditions :

\textbf{Théorème (Existence)} : Si $f: [a, b] \to [a, b]$ est une fonction continue sur l'intervalle fermé $[a, b]$ et que son image est contenue dans cet intervalle (c'est-à-dire $a \le f(x) \le b$ pour tout $x \in [a, b]$), alors $f$ admet au moins un point fixe dans $[a, b]$.

\textit{Preuve esquissée} : Considérons la fonction auxiliaire $g(x) = f(x) - x$.
Puisque $f$ est continue et $x$ est continue, $g$ est continue sur $[a, b]$.
De plus, $f(a) \ge a$ implique $g(a) = f(a) - a \ge 0$.
Et $f(b) \le b$ implique $g(b) = f(b) - b \le 0$.
Par le TVI, il existe $x^* \in [a, b]$ tel que $g(x^*) = 0$, ce qui signifie $f(x^*) - x^* = 0$, ou $f(x^*) = x^*$.

Ce théorème garantit l'existence, mais pas l'unicité, ni une méthode constructive pour trouver le point fixe.

\subsubsection{1.3 Théorème du Point Fixe de Banach (Théorème de l'Application Contractante)}

Pour garantir l'unicité et la convergence d'une méthode itérative, nous nous tournons vers une version plus forte du théorème, qui introduit la notion d'application contractante :

\textbf{Définition (Application Contractante)} : Une fonction $f: I \to I$ est dite \textbf{contractante} sur un intervalle $I$ s'il existe une constante $k \in [0, 1)$ telle que pour tout $x, y \in I$, on ait :
$|f(x) - f(y)| \le k|x - y|$.
La constante $k$ est appelée le coefficient de contraction.

Si $f$ est dérivable sur $I$, la condition de contraction est souvent vérifiée si $|f'(x)| \le k < 1$ pour tout $x \in I$.

\textbf{Théorème (Unicité et Convergence - Banach)} : Soit $I$ un intervalle fermé non vide. Si $f: I \to I$ est une application contractante sur $I$, alors $f$ admet un unique point fixe $x^* \in I$. De plus, pour toute valeur initiale $x_0 \in I$, la suite définie par l'itération $x_{n+1} = f(x_n)$ converge vers ce point fixe $x^*$.

Ce théorème est la base de la méthode itérative que nous allons implémenter. La continuité de $f$ est une conséquence de la propriété de contraction.

\subsubsection{1.4 Lien avec la Continuité}

La continuité est la pierre angulaire de ces théorèmes.
*   Le Théorème du Point Fixe simple utilise directement le Théorème des Valeurs Intermédiaires, qui est une propriété fondamentale des fonctions continues.
*   Une application contractante est nécessairement uniformément continue, et donc continue. La condition de contraction assure que les valeurs de la fonction ne s'éloignent pas trop rapidement, ce qui est essentiel pour la convergence de la suite itérative.

\subsection{Partie 2 : Implémentation de la Méthode Itérative du Point Fixe}

La méthode itérative du point fixe consiste à construire une suite $(x_n)_{n \in \mathbb{N}}$ définie par $x_{n+1} = f(x_n)$, à partir d'une valeur initiale $x_0$. Si la fonction $f$ est une contraction sur un intervalle contenant $x_0$ et le point fixe, cette suite converge vers le point fixe unique.

L'algorithme s'arrête lorsque la différence entre deux itérations successives devient suffisamment petite (inférieure à une tolérance \texttt{tol}), ou lorsque le nombre maximal d'itérations (\texttt{max\_iter}) est atteint pour éviter les boucles infinies en cas de non-convergence.

\begin{lstlisting}[language=python]
import math

def fixed_point_iteration(f, x0, tol=1e-6, max_iter=1000):
    """
    Implemente la methode iterative du point fixe pour trouver x tel que f(x) = x.

    Args:
        f (callable): La fonction pour laquelle on cherche un point fixe.
                      Doit prendre un float et retourner un float.
        x0 (float): L'estimation initiale du point fixe.
        tol (float): La tolerance pour la convergence. L'iteration s'arrete
                     quand |x_n+1 - x_n| < tol.
        max_iter (int): Le nombre maximal d'iterations.

    Returns:
        float: Le point fixe trouve si la convergence est atteinte, sinon None.
    """
    x_prev = x0
    print(f"Debut de l'iteration du point fixe avec x0 = {x0:.6f}, tol = {tol}, max_iter = {max_iter}")
    for i in range(max_iter):
        x_next = f(x_prev)
        # Critere d'arret: la difference entre l'iteration actuelle et la precedente est inferieure a la tolerance.
        # C'est une estimation de l'erreur absolue |x_next - x_exact| si la convergence est lineaire.
        if abs(x_next - x_prev) < tol:
            print(f"Convergence atteinte en {i+1} iterations. Point fixe: {x_next:.8f}")
            return x_next
        x_prev = x_next
    print(f"Avertissement: Nombre maximal d'iterations ({max_iter}) atteint sans convergence. Derniere valeur: {x_prev:.8f}")
    return None

\end{lstlisting}

\textbf{Explication du code :}

*   \texttt{f} : C'est la fonction mathématique $f(x)$ dont nous cherchons un point fixe. Elle est passée comme un objet \texttt{callable} (une fonction Python ou une lambda).
*   \texttt{x0} : La valeur initiale de notre suite itérative. Le choix de \texttt{x0} peut être crucial pour la convergence, surtout si la fonction n'est pas contractante sur tout l'intervalle.
*   \texttt{tol} : La tolérance. C'est le seuil en dessous duquel nous considérons que la convergence est atteinte. Plus \texttt{tol} est petit, plus la précision du point fixe sera élevée, mais plus le nombre d'itérations sera grand.
*   \texttt{max\_iter} : Le nombre maximal d'itérations. C'est une sécurité pour éviter que l'algorithme ne tourne indéfiniment si la suite ne converge pas ou converge trop lentement.
*   La boucle \texttt{for} effectue les itérations $x_{n+1} = f(x_n)$.
*   Le critère d'arrêt \texttt{abs(x\_next - x\_prev) < tol} est une heuristique courante pour détecter la convergence. Il est basé sur l'idée que si la suite converge, les termes successifs doivent devenir de plus en plus proches.
*   Si la convergence est atteinte, la fonction retourne le point fixe. Sinon, après \texttt{max\_iter} itérations, elle retourne \texttt{None}, indiquant un échec de convergence dans les limites données.

\subsection{Partie 3 : Validation et Tests Mathématiques avec Assertions}

Pour garantir la robustesse et la correction de notre implémentation, nous allons créer une série de tests utilisant des assertions (\texttt{assert}). Ces tests couvriront différents scénarios : fonctions convergentes, fonctions non convergentes, et l'impact des paramètres.

\begin{lstlisting}[language=python]
# --- Partie 3: Validation et Tests Mathematiques avec Assertions ---

def run_tests():
    print("--- Execution des tests pour la methode du point fixe ---")

    # Test Case 1: Fonction lineaire convergente f(x) = 0.5x + 1
    # Point fixe analytique: x = 0.5x + 1 => 0.5x = 1 => x = 2.0
    # Derivee: f'(x) = 0.5. Puisque |f'(x)| = 0.5 < 1, c'est une contraction, la convergence est garantie.
    print("Test 1: f(x) = 0.5x + 1, point fixe attendu = 2.0")
    f1 = lambda x: 0.5 * x + 1
    x0_1 = 0.0
    fixed_point_1 = fixed_point_iteration(f1, x0_1, tol=1e-7)
    assert fixed_point_1 is not None, "Test 1 echoue: La fonction devrait converger."
    assert abs(fixed_point_1 - 2.0) < 1e-6, f"Test 1 echoue: Point fixe incorrect. Attendu 2.0, obtenu {fixed_point_1}"
    assert abs(f1(fixed_point_1) - fixed_point_1) < 1e-7, "Test 1 echoue: f(x*) != x* (verification de la propriete du point fixe)"
    print(f"Test 1 Reussi: Point fixe trouve {fixed_point_1:.6f} est proche de 2.0")

    # Test Case 2: Fonction trigonometrique f(x) = cos(x)
    # Point fixe: x = cos(x). f'(x) = -sin(x).
    # Pour x dans un voisinage du point fixe (environ 0.739), |f'(x)| = |-sin(x)| < 1.
    # La convergence est attendue.
    print("Test 2: f(x) = cos(x), point fixe attendu ~ 0.739085")
    f2 = lambda x: math.cos(x)
    x0_2 = 0.5 # Une estimation initiale raisonnable
    fixed_point_2 = fixed_point_iteration(f2, x0_2, tol=1e-8, max_iter=5000)
    assert fixed_point_2 is not None, "Test 2 echoue: La fonction devrait converger."
    expected_fp_2 = 0.7390851332151607 # Valeur connue avec haute precision
    assert abs(fixed_point_2 - expected_fp_2) < 1e-7, f"Test 2 echoue: Point fixe incorrect. Attendu {expected_fp_2:.6f}, obtenu {fixed_point_2:.6f}"
    assert abs(f2(fixed_point_2) - fixed_point_2) < 1e-8, "Test 2 echoue: f(x*) != x* (verification de la propriete du point fixe)"
    print(f"Test 2 Reussi: Point fixe trouve {fixed_point_2:.6f} est proche de {expected_fp_2:.6f}")

    # Test Case 3: Fonction non convergente f(x) = 2x + 1
    # Derivee: f'(x) = 2. Puisque |f'(x)| = 2 > 1, ce n'est PAS une contraction. La methode ne devrait pas converger.
    print("Test 3: f(x) = 2x + 1, ne devrait PAS converger.")
    f3 = lambda x: 2 * x + 1
    x0_3 = 0.0
    fixed_point_3 = fixed_point_iteration(f3, x0_3, tol=1e-6, max_iter=100) # max_iter faible pour tester rapidement l'echec
    assert fixed_point_3 is None, "Test 3 echoue: La fonction ne devrait PAS converger."
    print("Test 3 Reussi: La methode n'a pas converge comme attendu.")

    # Test Case 4: Fonction avec un point fixe mais convergence lente ou dependante de x0
    # f(x) = sqrt(x+1). Point fixe: x^2 = x+1 => x^2 - x - 1 = 0. La racine positive est le nombre d'or phi = (1 + sqrt(5))/2.
    # Derivee: f'(x) = 1 / (2*sqrt(x+1)). Pour x > 0, |f'(x)| < 1 (ex: pour x=1, f'(1) = 1/(2*sqrt(2)) ~ 0.35 < 1).
    # La convergence est attendue.
    print("Test 4: f(x) = sqrt(x+1), point fixe attendu = phi (nombre d'or) ~ 1.618034")
    f4 = lambda x: math.sqrt(x + 1)
    x0_4 = 1.0
    fixed_point_4 = fixed_point_iteration(f4, x0_4, tol=1e-7, max_iter=2000)
    assert fixed_point_4 is not None, "Test 4 echoue: La fonction devrait converger."
    golden_ratio = (1 + math.sqrt(5)) / 2
    assert abs(fixed_point_4 - golden_ratio) < 1e-6, f"Test 4 echoue: Point fixe incorrect. Attendu {golden_ratio:.6f}, obtenu {fixed_point_4:.6f}"
    assert abs(f4(fixed_point_4) - fixed_point_4) < 1e-7, "Test 4 echoue: f(x*) != x* (verification de la propriete du point fixe)"
    print(f"Test 4 Reussi: Point fixe trouve {fixed_point_4:.6f} est proche de {golden_ratio:.6f}")

    # Test Case 5: Recherche d'une racine de h(x) = x^2 - 3x + 2 = 0 via point fixe de g(x) = (x^2 + 2)/3
    # Les racines de h(x) sont 1 et 2.
    # Pour g(x) = (x^2 + 2)/3, g'(x) = 2x/3.
    # Pour x=1, g'(1) = 2/3 < 1 (convergence vers 1).
    # Pour x=2, g'(2) = 4/3 > 1 (divergence ou convergence difficile).
    print("Test 5: Recherche d'une racine de x^2 - 3x + 2 = 0 via point fixe de g(x) = (x^2 + 2)/3")
    g5 = lambda x: (x**2 + 2) / 3

    # Test 5a: Convergence vers 1
    x0_5_a = 0.5 # Devrait converger vers 1
    fixed_point_5a = fixed_point_iteration(g5, x0_5_a, tol=1e-7, max_iter=1000)
    assert fixed_point_5a is not None, "Test 5a echoue: Devrait converger vers 1."
    assert abs(fixed_point_5a - 1.0) < 1e-6, f"Test 5a echoue: Point fixe incorrect. Attendu 1.0, obtenu {fixed_point_5a:.6f}"
    assert abs(g5(fixed_point_5a) - fixed_point_5a) < 1e-7, "Test 5a echoue: g(x*) != x*"
    print(f"Test 5a Reussi: Point fixe trouve {fixed_point_5a:.6f} est proche de 1.0")

    # Test 5b: Divergence ou non-convergence rapide vers 2
    x0_5_b = 2.5 # Devrait diverger ou ne pas converger rapidement vers 2 car |g'(2.5)| > 1
    fixed_point_5b = fixed_point_iteration(g5, x0_5_b, tol=1e-7, max_iter=100) # max_iter faible pour confirmer l'echec
    assert fixed_point_5b is None, "Test 5b echoue: Devrait diverger ou ne pas converger rapidement."
    print("Test 5b Reussi: La methode n'a pas converge comme attendu pour x0_5_b.")

    print("--- Tous les tests ont ete executes avec succes ! ---")

if __name__ == "__main__":
    run_tests()

\end{lstlisting}

\textbf{Explication des tests :}

Chaque test est conçu pour valider un aspect spécifique de la méthode du point fixe :

*   \textbf{Test 1 (\texttt{f(x) = 0.5x + 1})} : Une fonction linéaire simple avec un point fixe évident et une dérivée constante inférieure à 1 en valeur absolue. C'est un cas idéal pour la convergence rapide.
    *   Assertions : Vérifient que le point fixe est trouvé, qu'il est numériquement proche de la valeur attendue, et que la propriété $f(x^*) = x^*$ est satisfaite avec la tolérance donnée.
*   \textbf{Test 2 (\texttt{f(x) = cos(x)})} : Une fonction non linéaire classique dont le point fixe est souvent utilisé comme exemple. La dérivée varie, mais reste inférieure à 1 en valeur absolue autour du point fixe, garantissant la convergence.
    *   Assertions : Similaires au Test 1, avec une tolérance plus stricte pour un point fixe connu avec plus de décimales.
*   \textbf{Test 3 (\texttt{f(x) = 2x + 1})} : Une fonction dont la dérivée est supérieure à 1 en valeur absolue. La méthode du point fixe est censée diverger dans ce cas.
    *   Assertion : Vérifie que la fonction retourne \texttt{None}, indiquant l'absence de convergence dans le nombre maximal d'itérations.
*   \textbf{Test 4 (\texttt{f(x) = sqrt(x+1)})} : Une autre fonction non linéaire dont le point fixe est le nombre d'or. La convergence est assurée mais peut être plus lente que pour le Test 1.
    *   Assertions : Vérifient la convergence et la précision par rapport à la valeur analytique du nombre d'or.
*   \textbf{Test 5 (\texttt{g(x) = (x\textasciicircum 2 + 2)/3})} : Cet exemple illustre comment le choix de la fonction $g(x)$ (qui est une réécriture de $h(x)=0$) et de l'initialisation $x_0$ affecte la convergence. Pour $h(x) = x^2 - 3x + 2 = 0$, les racines sont $x=1$ et $x=2$.
    *   \textbf{Test 5a} : Montre la convergence vers $x=1$ car $|g'(1)| = 2/3 < 1$.
    *   \textbf{Test 5b} : Montre la non-convergence (ou divergence) pour une initialisation proche de $x=2$ car $|g'(2)| = 4/3 > 1$.

Ces tests couvrent les scénarios les plus importants pour évaluer la robustesse de l'implémentation.

\subsection{Consignes Supplémentaires et Questions}

1.  \textbf{Analyse de la Convergence} :
    *   Modifiez les valeurs de \texttt{x0} pour les fonctions convergentes (Test 1, 2, 4). Observez comment le nombre d'itérations change. Expliquez pourquoi un \texttt{x0} "proche" du point fixe est généralement préférable.
    *   Modifiez la tolérance \texttt{tol}. Comment cela affecte-t-il le nombre d'itérations et la précision du résultat ?
    *   Pour le Test 2 (\texttt{f(x) = cos(x)}), essayez \texttt{x0 = 100.0}. La méthode converge-t-elle toujours ? Pourquoi ? (Indice : la fonction cosinus est contractante sur $\mathbb{R}$ avec $k=1$ si on ne considère pas la dérivée, mais la dérivée est bornée par 1 en valeur absolue, ce qui est suffisant pour la convergence).

2.  \textbf{Limites de la Méthode} :
    *   Discutez des principales limitations de la méthode itérative du point fixe. Quand cette méthode est-elle inefficace ou inapplicable ?
    *   Comment pourrait-on détecter si une fonction n'est pas contractante sans calculer sa dérivée analytiquement ?

3.  \textbf{Améliorations Possibles} :
    *   Proposez une modification de la fonction \texttt{fixed\_point\_iteration} pour qu'elle retourne également le nombre d'itérations effectuées.
    *   Recherchez et décrivez brièvement une technique d'accélération de convergence pour la méthode du point fixe (par exemple, le procédé $\Delta^2$ d'Aitken). Expliquez son principe sans l'implémenter.
    *   Comment pourrait-on adapter cette méthode pour trouver les racines d'une équation $h(x) = 0$ (c'est-à-dire, trouver $x^*$ tel que $h(x^*) = 0$) ? Donnez un exemple de réarrangement pour $h(x) = x^3 - x - 1 = 0$.

Ce TP vous a permis d'implémenter un algorithme numérique fondamental basé sur les propriétés de continuité. La compréhension des conditions de convergence et des limites de ces méthodes est cruciale pour leur application judicieuse en ingénierie et en sciences.

\section{TP-03}
\section{TP 3 : Exploration Empirique de l'Uniforme Continuité et du Module de Continuité}

\subsection{Introduction}

Chers étudiants,

Dans les jalons précédents, nous avons approfondi la notion fondamentale de continuité d'une fonction en un point et sur un intervalle. Nous avons exploré la définition epsilon-delta qui en est le cœur, et ses implications topologiques et analytiques. Aujourd'hui, nous allons aborder une propriété plus forte et souvent cruciale en analyse numérique et en théorie de l'approximation : l'uniforme continuité.

La continuité d'une fonction $f$ sur un intervalle $I$ signifie que pour chaque point $x_0 \in I$ et pour tout $\epsilon > 0$, il existe un $\delta > 0$ tel que si $x$ est suffisamment proche de $x_0$ (i.e., $|x - x_0| < \delta$), alors $f(x)$ est suffisamment proche de $f(x_0)$ (i.e., $|f(x) - f(x_0)| < \epsilon$). La subtilité réside dans le fait que ce $\delta$ peut dépendre non seulement de $\epsilon$, mais aussi du point $x_0$.

L'uniforme continuité est une propriété plus exigeante. Elle stipule que le $\delta$ peut être choisi indépendamment du point $x_0$ (ou des points $x, y$ considérés). Cette distinction est capitale, car elle garantit une "régularité" de la fonction sur l'ensemble de l'intervalle, ce qui est essentiel pour des concepts comme l'intégrabilité de Riemann ou la convergence uniforme de suites de fonctions.

Dans ce Travail Pratique, nous allons implémenter en Python une méthode empirique pour évaluer le \textbf{module de continuité} d'une fonction. Le module de continuité est une fonction qui quantifie précisément la "dépendance" de $\delta$ par rapport à $\epsilon$ (ou plutôt, la dépendance de la variation de $f$ par rapport à la variation de sa variable). Son comportement lorsque $\delta$ tend vers zéro nous permettra de distinguer empiriquement les fonctions uniformément continues des autres.

\subsection{Contexte Mathématique}


Une fonction $f: I \to \mathbb{R}$ est continue sur un intervalle $I$ si pour tout $x_0 \in I$, pour tout $\epsilon > 0$, il existe un $\delta > 0$ (qui peut dépendre de $\epsilon$ et de $x_0$) tel que pour tout $x \in I$, si $|x - x_0| < \delta$, alors $|f(x) - f(x_0)| < \epsilon$.

\subsubsection{Définition : Uniforme Continuité}

Une fonction $f: I \to \mathbb{R}$ est uniformément continue sur un intervalle $I$ si pour tout $\epsilon > 0$, il existe un $\delta > 0$ (qui dépend \textit{uniquement} de $\epsilon$, et non des points $x, y \in I$) tel que pour tout $x, y \in I$, si $|x - y| < \delta$, alors $|f(x) - f(y)| < \epsilon$.

La différence cruciale réside dans l'ordre des quantificateurs :
*   Continuité : $\forall x_0 \in I, \forall \epsilon > 0, \exists \delta > 0, \dots$
*   Uniforme Continuité : $\forall \epsilon > 0, \exists \delta > 0, \forall x, y \in I, \dots$

\subsubsection{Définition : Module de Continuité}

Le module de continuité d'une fonction $f: I \to \mathbb{R}$ est une fonction $\omega_f: \mathbb{R}^+ \to \mathbb{R}^+$ définie par :
$$ \omega_f(\delta) = \sup \{ |f(x) - f(y)| : x, y \in I, |x - y| < \delta \} $$
où $\sup$ désigne le supremum (la borne supérieure).

Propriétés clés du module de continuité :
1.  $\omega_f(\delta)$ est une fonction croissante de $\delta$.
2.  $f$ est uniformément continue sur $I$ si et seulement si $\lim_{\delta \to 0^+} \omega_f(\delta) = 0$.

Notre objectif est d'estimer $\omega_f(\delta)$ empiriquement pour différentes valeurs de $\delta$ et d'observer son comportement.

\subsection{Stratégie d'Implémentation}

Le calcul exact du supremum sur un ensemble continu de points est généralement impossible numériquement sans outils d'analyse symbolique ou d'optimisation avancée. Nous allons donc adopter une approche empirique et discrète :

1.  \textbf{Discrétisation de l'intervalle} : Pour un intervalle $[a, b]$, nous allons générer un grand nombre de points équidistants dans cet intervalle. Plus le nombre de points est élevé, meilleure sera l'approximation.
2.  \textbf{Recherche du maximum de différence} : Pour une valeur $\delta$ donnée, nous allons parcourir toutes les paires de points $(x, y)$ générées. Si la distance $|x - y|$ est inférieure à $\delta$, nous calculons la différence $|f(x) - f(y)|$. Le maximum de toutes ces différences sera notre estimation empirique du module de continuité $\omega_f(\delta)$.
3.  \textbf{Observation du comportement} : En calculant $\omega_f(\delta)$ pour une séquence de $\delta$ décroissantes, nous pourrons observer si cette valeur tend vers zéro, ce qui indiquerait une uniforme continuité.

Il est crucial de comprendre que cette méthode est une \textit{approximation numérique}. La qualité de l'approximation dépendra fortement du nombre de points échantillonnés. Un nombre insuffisant de points pourrait masquer le vrai comportement de la fonction.

\subsection{Implémentation Python}

Nous allons implémenter les fonctions suivantes :
*   \texttt{generate\_points(interval\_start, interval\_end, num\_points)} : Génère une liste de points équidistants dans un intervalle donné.
*   \texttt{empirical\_modulus\_of\_continuity(f, interval\_start, interval\_end, delta, num\_sample\_points)} : Calcule l'estimation empirique de $\omega_f(\delta)$.

\begin{lstlisting}[language=python]
import math

def generate_points(interval_start: float, interval_end: float, num_points: int) -> list[float]:
    """
    Genere une liste de points equidistants dans un intervalle donne.

    Args:
        interval_start (float): Le debut de l'intervalle.
        interval_end (float): La fin de l'intervalle.
        num_points (int): Le nombre de points a generer. Doit etre >= 2.

    Returns:
        list[float]: Une liste de points flottants.
    """
    if num_points < 2:
        raise ValueError("Le nombre de points doit etre au moins 2 pour definir un intervalle.")
    if interval_start >= interval_end:
        raise ValueError("L'intervalle doit etre non vide et interval_start < interval_end.")

    points = []
    step = (interval_end - interval_start) / (num_points - 1)
    for i in range(num_points):
        points.append(interval_start + i * step)
    return points

def empirical_modulus_of_continuity(
    f,
    interval_start: float,
    interval_end: float,
    delta: float,
    num_sample_points: int
) -> float:
    """
    Calcule une estimation empirique du module de continuite omega_f(delta)
    pour une fonction f sur un intervalle [interval_start, interval_end].

    Args:
        f (callable): La fonction f(x) a evaluer.
        interval_start (float): Le debut de l'intervalle.
        interval_end (float): La fin de l'intervalle.
        delta (float): La valeur delta pour laquelle calculer omega_f(delta).
        num_sample_points (int): Le nombre de points a echantillonner dans l'intervalle.

    Returns:
        float: L'estimation empirique de omega_f(delta).
    """
    if delta <= 0:
        raise ValueError("delta doit etre strictement positif.")
    if num_sample_points < 2:
        raise ValueError("num_sample_points doit etre au moins 2.")

    points = generate_points(interval_start, interval_end, num_sample_points)
    max_diff = 0.0

    # Iteration sur toutes les paires de points
    # On utilise deux boucles pour s'assurer que toutes les paires (x, y) sont considerees,
    # y compris celles ou x et y sont differents mais tres proches.
    # On peut optimiser en ne considerant que x < y, car |f(x)-f(y)| = |f(y)-f(x)|.
    for i in range(len(points)):
        x = points[i]
        for j in range(i + 1, len(points)): # Optimisation: j commence a i+1
            y = points[j]
            if abs(x - y) < delta:
                diff = abs(f(x) - f(y))
                if diff > max_diff:
                    max_diff = diff
            else:
                # Puisque les points sont tries, si abs(x-y) >= delta pour y > x,
                # alors pour tout y' > y, abs(x-y') > delta.
                # On peut potentiellement optimiser en cassant la boucle interne,
                # mais cela n'est vrai que si f est monotone ou si l'on cherche
                # le sup sur des intervalles de longueur delta.
                # Pour le sup sur TOUTES les paires, il faut continuer.
                # Cependant, si y est deja trop loin de x, les y suivants le seront aussi.
                # Mais on doit considerer d'autres x.
                # La version actuelle est la plus simple et la plus robuste pour l'approximation.
                pass # Pas d'optimisation prematuree ici pour la robustesse.

    return max_diff

# --- Fonctions de test ---
def f_x_squared(x: float) -> float:
    """Fonction f(x) = x\textasciicircum 2."""
    return x * x

def f_one_over_x(x: float) -> float:
    """Fonction f(x) = 1/x."""
    if x == 0:
        raise ValueError("Division par zero pour f(x) = 1/x.")
    return 1.0 / x

def f_sin_x(x: float) -> float:
    """Fonction f(x) = sin(x)."""
    return math.sin(x)

# --- Bloc d'execution principal et validations ---
if __name__ == "__main__":
    print("--- TP 3 : Exploration Empirique de l'Uniforme Continuite et du Module de Continuite ---")
    print("Validation des fonctions utilitaires...")
    # Test generate_points
    points_test = generate_points(0, 1, 5)
    assert len(points_test) == 5
    assert points_test[0] == 0.0
    assert points_test[-1] == 1.0
    assert points_test == [0.0, 0.25, 0.5, 0.75, 1.0]
    print("generate_points: OK")

    try:
        generate_points(1, 0, 5)
        assert False, "generate_points devrait lever une erreur pour interval_start >= interval_end"
    except ValueError:
        assert True
    print("generate_points (erreur intervalle): OK")

    try:
        empirical_modulus_of_continuity(f_x_squared, 0, 1, 0.1, 1)
        assert False, "empirical_modulus_of_continuity devrait lever une erreur pour num_sample_points < 2"
    except ValueError:
        assert True
    print("empirical_modulus_of_continuity (erreur num_sample_points): OK")

    print("--- Exemples et Validation ---")

    # --- Exemple 1 : Fonction f(x) = x\textasciicircum 2 sur [0, 1] (Uniformement Continue) ---
    # Theoriquement, pour f(x) = x\textasciicircum 2 sur [0, 1], |f(x) - f(y)| = |x^2 - y^2| = |x - y||x + y|.
    # Si |x - y| < delta et x, y in [0, 1], alors |x + y| <= 2.
    # Donc, |f(x) - f(y)| < delta * 2.
    # Ainsi, omega_f(delta) <= 2 * delta.
    print("### Exemple 1 : Fonction f(x) = x\textasciicircum 2 sur [0, 1] (Uniformement Continue)")
    interval_a1, interval_b1 = 0.0, 1.0
    num_samples1 = 10000 # Un grand nombre de points pour une bonne approximation

    print(f"Evaluation de f(x) = x\textasciicircum 2 sur [{interval_a1}, {interval_b1}]")
    delta_values1 = [0.2, 0.1, 0.05, 0.01, 0.005, 0.001]
    results1 = []
    for d in delta_values1:
        omega_est = empirical_modulus_of_continuity(f_x_squared, interval_a1, interval_b1, d, num_samples1)
        results1.append((d, omega_est))
        print(f"  delta = {d:.3f}, omega_f(delta)  {omega_est:.6f} (Attendu <= {2*d:.6f})")

    # Assertions pour f(x) = x\textasciicircum 2
    # On s'attend a ce que omega_f(delta) soit proche de 2*delta et tende vers 0.
    assert results1[0][1] < 2 * delta_values1[0] + 0.01 # Tolerance pour l'approximation
    assert results1[1][1] < 2 * delta_values1[1] + 0.005
    assert results1[3][1] < 2 * delta_values1[3] + 0.001
    assert results1[-1][1] < 0.003 # Doit etre tres petit pour delta=0.001
    assert results1[-1][1] < results1[0][1] # Doit decroitre
    print("Assertions pour f(x) = x\textasciicircum 2: OK. Le module de continuite tend vers 0.")

    # --- Exemple 2 : Fonction f(x) = 1/x sur (0, 1] (Non Uniformement Continue) ---
    # Theoriquement, pour f(x) = 1/x sur (0, 1], le module de continuite ne tend pas vers 0.
    # En effet, si on prend x = delta et y = 2*delta (avec delta petit),
    # |x - y| = delta. Alors |f(x) - f(y)| = |1/delta - 1/(2*delta)| = 1/(2*delta).
    # Cette valeur tend vers l'infini quand delta tend vers 0.
    # Nous devons choisir un intervalle [epsilon_min, 1] pour eviter la division par zero.
    print("### Exemple 2 : Fonction f(x) = 1/x sur (0, 1] (Non Uniformement Continue)")
    interval_a2, interval_b2 = 0.0001, 1.0 # On commence a une petite valeur > 0
    num_samples2 = 20000 # Encore plus de points car la fonction varie fortement

    print(f"Evaluation de f(x) = 1/x sur [{interval_a2}, {interval_b2}]")
    delta_values2 = [0.01, 0.005, 0.001, 0.0005, 0.0001]
    results2 = []
    for d in delta_values2:
        omega_est = empirical_modulus_of_continuity(f_one_over_x, interval_a2, interval_b2, d, num_samples2)
        results2.append((d, omega_est))
        print(f"  delta = {d:.5f}, omega_f(delta)  {omega_est:.6f}")

    # Assertions pour f(x) = 1/x
    # On s'attend a ce que omega_f(delta) reste grand ou meme augmente (si l'intervalle est bien choisi)
    # et ne tende PAS vers 0.
    # Pour delta = 0.0001, on peut avoir des points x, y tres proches de 0.0001.
    # Par exemple, x=0.0001, y=0.0002. |x-y|=0.0001. |f(x)-f(y)| = |10000 - 5000| = 5000.
    # L'estimation devrait etre une valeur significative.
    assert results2[0][1] > 10.0 # Doit etre significatif meme pour delta "grand"
    assert results2[-1][1] > 1000.0 # Doit etre tres grand pour un petit delta
    # On peut meme s'attendre a ce que la valeur augmente ou reste elevee pour des deltas plus petits
    assert results2[-1][1] > results2[0][1] # Le module de continuite augmente ou reste eleve
    print("Assertions pour f(x) = 1/x: OK. Le module de continuite ne tend PAS vers 0.")

    # --- Exemple 3 : Fonction f(x) = sin(x) sur [0, 2*pi] (Uniformement Continue) ---
    # La fonction sin(x) est uniformement continue sur tout intervalle borne.
    # Theoriquement, |sin(x) - sin(y)| <= |x - y| (car |cos(c)| <= 1).
    # Donc, omega_f(delta) <= delta.
    print("### Exemple 3 : Fonction f(x) = sin(x) sur [0, 2*pi] (Uniformement Continue)")
    interval_a3, interval_b3 = 0.0, 2 * math.pi
    num_samples3 = 10000

    print(f"Evaluation de f(x) = sin(x) sur [{interval_a3:.2f}, {interval_b3:.2f}]")
    delta_values3 = [0.2, 0.1, 0.05, 0.01, 0.005, 0.001]
    results3 = []
    for d in delta_values3:
        omega_est = empirical_modulus_of_continuity(f_sin_x, interval_a3, interval_b3, d, num_samples3)
        results3.append((d, omega_est))
        print(f"  delta = {d:.3f}, omega_f(delta)  {omega_est:.6f} (Attendu <= {d:.6f})")

    # Assertions pour f(x) = sin(x)
    assert results3[0][1] < delta_values3[0] + 0.01
    assert results3[1][1] < delta_values3[1] + 0.005
    assert results3[3][1] < delta_values3[3] + 0.001
    assert results3[-1][1] < 0.002 # Tres petit pour delta=0.001
    assert results3[-1][1] < results3[0][1]
    print("Assertions pour f(x) = sin(x): OK. Le module de continuite tend vers 0.")

    print("--- Tous les tests et validations sont passes avec succes. ---")

\end{lstlisting}

\subsection{Exemples et Validation}

Nous avons testé notre implémentation avec trois fonctions représentatives pour illustrer les concepts d'uniforme continuité et de non-uniforme continuité.

\subsubsection{Exemple 1 : Fonction $f(x) = x\textasciicircum 2$ sur $[0, 1]$ (Uniformément Continue)}

Sur un intervalle fermé et borné, toute fonction continue est uniformément continue (Théorème de Heine). $f(x) = x\textasciicircum 2$ est continue sur $[0, 1]$, elle est donc uniformément continue.
Théoriquement, pour $x, y \in [0, 1]$, on a $|f(x) - f(y)| = |x^2 - y^2| = |x - y||x + y|$. Puisque $x, y \in [0, 1]$, on a $x+y \le 2$. Donc, si $|x - y| < \delta$, alors $|f(x) - f(y)| < 2\delta$. Le module de continuité $\omega_f(\delta)$ est donc majoré par $2\delta$. Nous nous attendons à ce que $\omega_f(\delta)$ tende vers 0 linéairement avec $\delta$.

Les résultats empiriques montrent que pour des valeurs de $\delta$ décroissantes, l'estimation de $\omega_f(\delta)$ décroît également et tend vers zéro, confirmant l'uniforme continuité. Les assertions vérifient que les valeurs obtenues sont cohérentes avec la borne théorique $2\delta$ et que la tendance est bien à la décroissance vers zéro.

\subsubsection{Exemple 2 : Fonction $f(x) = 1/x$ sur $(0, 1]$ (Non Uniformément Continue)}

La fonction $f(x) = 1/x$ est continue sur $(0, 1]$, mais elle n'est pas uniformément continue. La raison est que la pente de la fonction devient arbitrairement grande à mesure que $x$ s'approche de 0. Pour un $\epsilon$ donné, il est impossible de trouver un $\delta$ unique qui fonctionne pour tous les points. Si nous prenons $x$ et $y$ très proches de 0, même si $|x - y|$ est très petit, la différence $|f(x) - f(y)|$ peut être arbitrairement grande.
Par exemple, si $x = \delta$ et $y = 2\delta$ (pour un $\delta$ très petit), alors $|x - y| = \delta$. Mais $|f(x) - f(y)| = |1/\delta - 1/(2\delta)| = 1/(2\delta)$, qui tend vers l'infini lorsque $\delta \to 0$.

Les résultats empiriques, obtenus en échantillonnant l'intervalle $[0.0001, 1.0]$ pour éviter la division par zéro, montrent que l'estimation de $\omega_f(\delta)$ reste très élevée, voire augmente, lorsque $\delta$ diminue. Cela indique clairement que la fonction n'est pas uniformément continue sur cet intervalle. Les assertions confirment que le module de continuité ne tend pas vers zéro.

\subsubsection{Exemple 3 : Fonction $f(x) = \sin(x)$ sur $[0, 2\pi]$ (Uniformément Continue)}

La fonction $f(x) = \sin(x)$ est continue sur tout $\mathbb{R}$, et donc sur l'intervalle fermé et borné $[0, 2\pi]$. Elle est donc uniformément continue.
Théoriquement, en utilisant le théorème des accroissements finis, $|\sin(x) - \sin(y)| = |\cos(c)(x - y)|$ pour un certain $c$ entre $x$ et $y$. Puisque $|\cos(c)| \le 1$, on a $|\sin(x) - \sin(y)| \le |x - y|$. Ainsi, le module de continuité $\omega_f(\delta)$ est majoré par $\delta$.

Les résultats empiriques pour $\sin(x)$ confirment cette attente : l'estimation de $\omega_f(\delta)$ décroît et tend vers zéro à mesure que $\delta$ diminue, et les valeurs sont proches de $\delta$. Les assertions valident cette tendance et la conformité avec la borne théorique.

\subsection{Conclusion}

Ce Travail Pratique nous a permis d'explorer la notion d'uniforme continuité et du module de continuité à travers une implémentation Python "from scratch". Nous avons vu comment une approche empirique, basée sur la discrétisation d'un intervalle et la recherche du maximum de différences de fonction, peut nous donner une intuition précieuse sur le comportement de ces propriétés mathématiques.

Nous avons clairement distingué les fonctions uniformément continues (où $\omega_f(\delta) \to 0$ quand $\delta \to 0$) des fonctions non uniformément continues (où $\omega_f(\delta)$ ne tend pas vers 0). Cette distinction est fondamentale en analyse et a des répercussions directes sur la validité de nombreux théorèmes et méthodes numériques.

Il est important de noter les limitations de cette approche empirique :
*   \textbf{Précision} : L'estimation dépend fortement du nombre de points échantillonnés. Un nombre insuffisant de points peut conduire à des conclusions erronées.
*   \textbf{Intervalle} : Le comportement de l'uniforme continuité est lié à l'intervalle. Une fonction peut être uniformément continue sur un intervalle et non sur un autre (par exemple, $x^2$ sur $[0, 1]$ vs. sur $[0, \infty)$).
\textit{   \textbf{Preuve formelle} : Cette méthode fournit une }indication* empirique, mais ne constitue en aucun cas une preuve mathématique formelle de l'uniforme continuité.

Malgré ces limitations, l'exercice de construire un tel outil nous aide à solidifier notre compréhension des définitions epsilon-delta et à apprécier la complexité et la richesse des concepts de continuité en analyse réelle.

\section{TP-04}
\section{TP 4: Exploration Empirique du Module de Continuité et de l'Uniforme Continuité}

\subsection{Introduction}

Chers étudiants,

Bienvenue à ce quatrième Travail Pratique, qui s'inscrit dans le cadre de notre Jalon 18 sur la "Continuité des fonctions d'une variable réelle". En tant que Professeur d'Élite et Développeur Senior, je vous propose d'aller au-delà des définitions théoriques pour explorer numériquement une notion fondamentale : le \textbf{module de continuité} et, par extension, la \textbf{continuité uniforme}.

La continuité est une propriété locale, décrivant le comportement d'une fonction autour d'un point. La continuité uniforme, en revanche, est une propriété globale, garantissant que la variation de la fonction peut être rendue arbitrairement petite sur l'ensemble du domaine, pour des points suffisamment proches, \textit{indépendamment de leur position}. Cette distinction est cruciale en analyse réelle et a des implications profondes en calcul numérique, notamment pour l'intégration et l'approximation de fonctions.

L'objectif de ce TP est d'implémenter en Python pur, "from scratch", un outil permettant d'approximer numériquement le module de continuité d'une fonction sur un intervalle donné, et d'utiliser cette approximation pour évaluer empiriquement si une fonction est uniformément continue.



\subsubsection{Continuité}

Une fonction $f: I \to \mathbb{R}$ est dite continue en un point $x_0 \in I$ si pour tout $\varepsilon > 0$, il existe $\delta > 0$ tel que pour tout $x \in I$, si $|x - x_0| < \delta$, alors $|f(x) - f(x_0)| < \varepsilon$.

\subsubsection{Continuité Uniforme}

Une fonction $f: I \to \mathbb{R}$ est dite uniformément continue sur l'intervalle $I$ si pour tout $\varepsilon > 0$, il existe $\delta > 0$ tel que pour tout $x, y \in I$, si $|x - y| < \delta$, alors $|f(x) - f(y)| < \varepsilon$.
La différence clé avec la continuité simple est que le $\delta$ ne dépend que de $\varepsilon$ (et de la fonction $f$), et non des points $x$ et $y$ spécifiques.

\subsubsection{Module de Continuité}

Le module de continuité $\omega_f(\delta)$ (souvent noté simplement $\omega(\delta)$) d'une fonction $f$ sur un intervalle $I$ est défini comme :
$$ \omega(\delta) = \sup \{ |f(x) - f(y)| : x, y \in I, |x - y| < \delta \} $$
où $\sup$ désigne le supremum (la borne supérieure).

Propriétés importantes du module de continuité :
1.  $\omega(\delta)$ est une fonction croissante de $\delta$.
2.  $\omega(\delta) \ge 0$.
3.  Une fonction $f$ est uniformément continue sur $I$ si et seulement si $\lim_{\delta \to 0^+} \omega(\delta) = 0$.

\subsection{Objectif du TP}

Votre tâche est de développer un script Python qui :
1.  \textbf{Implémente une fonction \texttt{evaluate\_modulus\_of\_continuity}} : Cette fonction prendra en entrée une fonction $f$, un intervalle $I$, une valeur $\delta$, et un nombre d'échantillons. Elle retournera une approximation numérique de $\omega(\delta)$ en échantillonnant aléatoirement des paires de points $(x, y)$ dans $I$ telles que $|x - y| \le \delta$, et en trouvant la différence maximale $|f(x) - f(y)|$.
2.  \textbf{Implémente une fonction \texttt{check\_uniform\_continuity\_empirically}} : Cette fonction utilisera la précédente pour évaluer $\omega(\delta)$ pour une série de valeurs de $\delta$ décroissantes vers zéro. Elle analysera la tendance de $\omega(\delta)$ pour déterminer empiriquement si la fonction est uniformément continue sur l'intervalle donné.
3.  \textbf{Valide la logique avec des assertions mathématiques} : Des tests rigoureux seront effectués sur des fonctions connues pour être uniformément continues ou non, afin de vérifier la robustesse de l'implémentation.

\textbf{Contrainte majeure :} Vous devez utiliser Python pur. L'utilisation de bibliothèques comme \texttt{numpy}, \texttt{scipy} ou \texttt{matplotlib} est interdite pour la logique de calcul du module de continuité. \texttt{math} et \texttt{random} sont autorisés.

\subsection{Implémentation Python}

Nous allons commencer par définir les fonctions nécessaires.

\begin{lstlisting}[language=python]
import random
import math

def evaluate_modulus_of_continuity(f, interval, delta, num_samples=100000):
    """
    Approximates the modulus of continuity omega(delta) for a function f
    on a given interval.

    omega(delta) = sup { |f(x) - f(y)| : x, y in interval, |x - y| <= delta }

    The supremum is approximated by sampling a large number of random pairs (x, y)
    that satisfy the conditions and finding the maximum difference |f(x) - f(y)|.
    We use |x - y| <= delta for practical numerical sampling, which is equivalent
    to |x - y| < delta for continuous functions when taking the supremum.

    Args:
        f (callable): The function f(x).
        interval (tuple): A tuple (a, b) representing the closed interval [a, b].
        delta (float): The delta value. Must be positive.
        num_samples (int): Number of random pairs (x, y) to sample to approximate the supremum.
                           A higher number of samples leads to a better approximation but
                           increases computation time.

    Returns:
        float: An approximation of omega(delta).
    """
    a, b = interval
    if not (a <= b):
        raise ValueError("Interval start 'a' must be less than or equal to interval end 'b'.")
    if delta <= 0:
        raise ValueError("Delta must be positive.")
    if num_samples <= 0:
        raise ValueError("Number of samples must be positive.")

    max_diff = 0.0

    for _ in range(num_samples):
        # 1. Pick a random x within the interval [a, b]
        x = random.uniform(a, b)

        # 2. Define the valid range for y such that |x - y| <= delta and y is also in [a, b]
        # This range is [max(a, x - delta), min(b, x + delta)]
        y_lower_bound = max(a, x - delta)
        y_upper_bound = min(b, x + delta)

        # 3. Pick a random y within this constrained range.
        # Handle cases where the range is degenerate (y_lower_bound >= y_upper_bound).
        # This typically means x is very close to a boundary and delta is small,
        # or delta is so small that only x itself (or points extremely close to x) are valid.
        if y_lower_bound >= y_upper_bound:
            # If the range is degenerate, the only valid y is effectively x itself.
            # In this case, |f(x) - f(y)| = |f(x) - f(x)| = 0.
            y = x
        else:
            y = random.uniform(y_lower_bound, y_upper_bound)

        # Calculate the absolute difference and update max_diff if it's larger
        diff = abs(f(x) - f(y))
        if diff > max_diff:
            max_diff = diff

    return max_diff

def check_uniform_continuity_empirically(f, interval, delta_values, num_samples_per_delta=100000, tolerance=0.04):
    """
    Empirically checks for uniform continuity of a function f on an interval
    by observing the trend of omega(delta) as delta approaches zero.

    A function f is uniformly continuous on I if and only if lim (delta->0+) omega(delta) = 0.
    This function approximates omega(delta) for a series of decreasing delta values
    and checks if the last (smallest delta) omega value is below a given tolerance.

    Args:
        f (callable): The function f(x).
        interval (tuple): A tuple (a, b) representing the closed interval [a, b].
        delta_values (list): A list of positive delta values, typically decreasing towards zero.
                             Example: [0.1, 0.05, 0.01, 0.005, 0.001].
        num_samples_per_delta (int): Number of samples for each omega(delta) evaluation.
        tolerance (float): The threshold for the smallest omega(delta) value to be considered
                           "close to zero" and thus indicative of uniform continuity.
                           This value needs to be chosen carefully, as the rate at which
                           omega(delta) approaches zero varies between functions.

    Returns:
        tuple: (bool, list of floats) - True if the function is likely uniformly continuous
               on the interval, False otherwise. Also returns the list of computed omega(delta) values.
    """
    if not delta_values:
        raise ValueError("delta_values list cannot be empty.")
    if not all(d > 0 for d in delta_values):
        raise ValueError("All delta values must be positive.")

    # Optional warning if delta_values are not strictly decreasing, as this is the typical use case.
    if len(delta_values) > 1 and not all(delta_values[i] > delta_values[i+1] for i in range(len(delta_values)-1)):
        print("Warning: delta_values are not strictly decreasing. This might affect trend interpretation.")

    omega_values = []
    print(f"\n--- Evaluation de la continuite uniforme pour la fonction {f.__name__} sur {interval} ---")
    for delta in delta_values:
        omega = evaluate_modulus_of_continuity(f, interval, delta, num_samples_per_delta)
        omega_values.append(omega)
        print(f"  delta = {delta:.6f}, omega(delta) = {omega:.6f}")

    # Le critere heuristique principal : verifier si le module de continuite pour le plus petit delta
    # est inferieur a la tolerance specifiee.
    is_uniformly_continuous = omega_values[-1] < tolerance

    print(f"Conclusion: La fonction {f.__name__} est {'PROBABLEMENT' if is_uniformly_continuous else 'IMPROBABLEMENT'} uniformement continue sur {interval} (omega({delta_values[-1]:.6f}) = {omega_values[-1]:.6f} vs tolerance {tolerance:.6f}).")

    return is_uniformly_continuous, omega_values

# --- Fonctions de test ---
def f_linear(x):
    """Une fonction lineaire simple: f(x) = x."""
    return x

def f_x_squared(x):
    """Une fonction quadratique: f(x) = x\textasciicircum 2."""
    return x * x

def f_one_over_x(x):
    """Une fonction qui n'est pas uniformement continue sur (0, 1]."""
    # Pour eviter la division par zero et simuler (0, 1], nous utiliserons un intervalle [epsilon, 1].
    if x == 0:
        raise ValueError("La fonction 1/x est indefinie en x=0. Ajustez l'intervalle pour eviter 0.")
    return 1.0 / x

def f_sin_x(x):
    """La fonction sinus: f(x) = sin(x)."""
    return math.sin(x)

def f_sqrt_x(x):
    """La fonction racine carree: f(x) = sqrt(x)."""
    if x < 0:
        raise ValueError("sqrt(x) n'est pas definie pour x < 0 dans les nombres reels.")
    return math.sqrt(x)

\end{lstlisting}

\subsection{Validation Mathématique et Tests}

Nous allons maintenant valider notre implémentation avec une série de tests basés sur des propriétés connues de fonctions. Les assertions \texttt{assert} nous permettront de vérifier que notre approche numérique concorde avec la théorie mathématique.

Pour chaque test, nous utiliserons un ensemble de \texttt{delta\_values} décroissantes et une \texttt{tolerance} pour juger de la continuité uniforme. Le \texttt{num\_samples\_per\_delta} est fixé à \texttt{100000} pour une meilleure précision. La \texttt{DEFAULT\_TOLERANCE} est choisie pour permettre aux fonctions uniformément continues de passer le test pour le plus petit \texttt{delta} utilisé.

\begin{lstlisting}[language=python]
# Configuration des tests
DEFAULT_DELTA_VALUES = [0.1, 0.05, 0.01, 0.005, 0.001]
DEFAULT_NUM_SAMPLES = 100000
DEFAULT_TOLERANCE = 0.04 # Tolerance ajustee pour les verifications empiriques

print("--- Demarrage des tests de validation ---")

# Test 1: Fonction lineaire f(x) = x sur [0, 1]
# Theorie: Uniformement continue (Lipschitz, omega(delta) = delta)
print("--- Test 1: f(x) = x sur [0, 1] ---")
is_uc, omega_vals = check_uniform_continuity_empirically(
    f_linear, (0, 1), DEFAULT_DELTA_VALUES, DEFAULT_NUM_SAMPLES, DEFAULT_TOLERANCE
)
assert is_uc, "Test 1 FAILED: f(x)=x devrait etre uniformement continue."
# Verifions que omega(delta) est proche de delta
assert abs(omega_vals[0] - DEFAULT_DELTA_VALUES[0]) < 0.01, "Test 1 FAILED: omega(0.1) pour f(x)=x non proche de 0.1"
assert abs(omega_vals[-1] - DEFAULT_DELTA_VALUES[-1]) < 0.01, "Test 1 FAILED: omega(0.001) pour f(x)=x non proche de 0.001"
print("Test 1 PASSED: f(x)=x est correctement identifiee comme uniformement continue.")

# Test 2: Fonction quadratique f(x) = x\textasciicircum 2 sur [0, 1]
# Theorie: Uniformement continue (Lipschitz, omega(delta) approx 2*delta pour petit delta)
print("--- Test 2: f(x) = x\textasciicircum 2 sur [0, 1] ---")
is_uc, omega_vals = check_uniform_continuity_empirically(
    f_x_squared, (0, 1), DEFAULT_DELTA_VALUES, DEFAULT_NUM_SAMPLES, DEFAULT_TOLERANCE
)
assert is_uc, "Test 2 FAILED: f(x)=x\textasciicircum 2 devrait etre uniformement continue."
# Pour f(x)=x\textasciicircum 2 sur [0,1], omega(delta) = 2*delta - delta^2.
# Pour delta=0.1, omega(0.1) = 0.2 - 0.01 = 0.19
# Pour delta=0.001, omega(0.001) = 0.002 - 0.000001 = 0.001999
assert abs(omega_vals[0] - (2*DEFAULT_DELTA_VALUES[0] - DEFAULT_DELTA_VALUES[0]**2)) < 0.01, "Test 2 FAILED: omega(0.1) pour f(x)=x\textasciicircum 2 non proche de la valeur theorique."
assert abs(omega_vals[-1] - (2*DEFAULT_DELTA_VALUES[-1] - DEFAULT_DELTA_VALUES[-1]**2)) < 0.01, "Test 2 FAILED: omega(0.001) pour f(x)=x\textasciicircum 2 non proche de la valeur theorique."
print("Test 2 PASSED: f(x)=x\textasciicircum 2 est correctement identifiee comme uniformement continue.")

# Test 3: Fonction f(x) = 1/x sur (0, 1]
# Theorie: NON uniformement continue. omega(delta) ne tend pas vers 0.
# Pour eviter la division par zero et simuler (0, 1], nous utilisons [epsilon, 1].
print("--- Test 3: f(x) = 1/x sur [1e-6, 1] ---")
# On s'attend a ce que omega(delta) reste eleve meme pour de petits delta.
# La tolerance par defaut (0.04) est suffisamment petite pour que cette fonction echoue au test.
is_uc, omega_vals = check_uniform_continuity_empirically(
    f_one_over_x, (1e-6, 1), DEFAULT_DELTA_VALUES, DEFAULT_NUM_SAMPLES, DEFAULT_TOLERANCE
)
assert not is_uc, "Test 3 FAILED: f(x)=1/x ne devrait PAS etre uniformement continue."
assert omega_vals[-1] > 1.0, "Test 3 FAILED: omega(delta) pour f(x)=1/x devrait etre grand." # Verifie qu'il est effectivement grand
print("Test 3 PASSED: f(x)=1/x est correctement identifiee comme NON uniformement continue.")

# Test 4: Fonction f(x) = sin(x) sur [0, 2*pi]
# Theorie: Uniformement continue (Lipschitz, |f'(x)| <= 1, donc omega(delta) <= delta)
print("--- Test 4: f(x) = sin(x) sur [0, 2*pi] ---")
is_uc, omega_vals = check_uniform_continuity_empirically(
    f_sin_x, (0, 2 * math.pi), DEFAULT_DELTA_VALUES, DEFAULT_NUM_SAMPLES, DEFAULT_TOLERANCE
)
assert is_uc, "Test 4 FAILED: f(x)=sin(x) devrait etre uniformement continue."
# Pour sin(x), omega(delta) <= delta.
assert omega_vals[-1] < DEFAULT_DELTA_VALUES[-1] * 1.1, "Test 4 FAILED: omega(0.001) pour f(x)=sin(x) non proche de 0.001."
print("Test 4 PASSED: f(x)=sin(x) est correctement identifiee comme uniformement continue.")

# Test 5: Fonction f(x) = sqrt(x) sur [0, 1]
# Theorie: Uniformement continue (mais pas Lipschitz sur [0,1] a cause de la derivee en 0).
# omega(delta) = sqrt(delta) pour cette fonction.
print("--- Test 5: f(x) = sqrt(x) sur [0, 1] ---")
is_uc, omega_vals = check_uniform_continuity_empirically(
    f_sqrt_x, (0, 1), DEFAULT_DELTA_VALUES, DEFAULT_NUM_SAMPLES, DEFAULT_TOLERANCE
)
assert is_uc, "Test 5 FAILED: f(x)=sqrt(x) devrait etre uniformement continue."
# Pour f(x)=sqrt(x) sur [0,1], omega(delta) = sqrt(delta).
# Pour delta=0.001, omega(0.001) = sqrt(0.001) approx 0.0316
assert abs(omega_vals[-1] - math.sqrt(DEFAULT_DELTA_VALUES[-1])) < 0.01, "Test 5 FAILED: omega(0.001) pour f(x)=sqrt(x) non proche de sqrt(0.001)."
print("Test 5 PASSED: f(x)=sqrt(x) est correctement identifiee comme uniformement continue.")

print("--- Tous les tests de validation ont ete executes avec succes ! ---")

\end{lstlisting}

\subsection{Conclusion}

Félicitations ! Vous avez implémenté avec succès un outil d'exploration numérique pour le module de continuité et la continuité uniforme. Ce TP vous a permis de :

*   \textbf{Comprendre en profondeur} les définitions de la continuité et de la continuité uniforme, ainsi que le rôle central du module de continuité $\omega(\delta)$.
*   \textbf{Développer une intuition numérique} pour ces concepts en observant le comportement de $\omega(\delta)$ lorsque $\delta$ tend vers zéro.
*   \textbf{Mettre en pratique vos compétences en programmation Python} pour traduire des définitions mathématiques abstraites en algorithmes concrets, en respectant des contraintes strictes (Python pur, "from scratch").
*   \textbf{Valider votre implémentation} à l'aide d'assertions basées sur des résultats théoriques connus, renforçant ainsi la fiabilité de votre code.

Il est important de noter que l'approche par échantillonnage aléatoire est une approximation. La précision de $\omega(\delta)$ dépend fortement du nombre d'échantillons (\texttt{num\_samples}). Pour des fonctions plus complexes ou des exigences de précision plus élevées, un nombre d'échantillons beaucoup plus grand serait nécessaire, ou des méthodes d'optimisation plus sophistiquées pour trouver le supremum. Néanmoins, pour une exploration empirique et une compréhension conceptuelle, cette méthode est très efficace.

Ce TP démontre comment les outils de développement peuvent être utilisés pour sonder et valider des concepts mathématiques fondamentaux, jetant un pont entre la théorie et la pratique. Continuez à explorer !

\section{TP-05}
\section{TP 5 : Exploration Empirique de la Continuité Uniforme et du Module de Continuité}

\subsection{Introduction}

Dans le domaine de l'analyse réelle, la notion de continuité est fondamentale. Une fonction est continue si de petites variations de l'entrée entraînent de petites variations de la sortie. Cependant, il existe une distinction cruciale entre la continuité simple et la continuité uniforme. La continuité simple permet que la "taille" de la variation de l'entrée ($\delta$) dépende non seulement de la "taille" de la variation de la sortie désirée ($\epsilon$), mais aussi du point spécifique où la continuité est évaluée. La continuité uniforme, en revanche, exige qu'un unique $\delta$ puisse être trouvé pour un $\epsilon$ donné, et ce $\delta$ fonctionne pour \textit{tous} les points de l'intervalle.

Le module de continuité, noté $\omega_f(\delta)$, est un outil puissant pour caractériser la continuité uniforme. Il mesure la plus grande variation possible de la fonction pour des points d'entrée séparés par une distance inférieure à $\delta$. Une fonction est uniformément continue si et seulement si son module de continuité tend vers zéro lorsque $\delta$ tend vers zéro.

Ce Travail Pratique vous mettra au défi d'implémenter de manière empirique l'estimation du module de continuité et d'utiliser cette estimation pour distinguer les fonctions uniformément continues des fonctions qui ne le sont pas, en utilisant uniquement Python pur.

\subsection{Objectifs Pédagogiques}

À la fin de ce TP, vous devriez être capable de :
1.  Comprendre la définition formelle et la signification intuitive de la continuité uniforme et du module de continuité.
2.  Implémenter une méthode d'estimation numérique du module de continuité pour une fonction donnée sur un intervalle.
3.  Utiliser cette estimation pour analyser la tendance du module de continuité et en déduire empiriquement la continuité uniforme.
4.  Démontrer la différence de comportement du module de continuité entre des fonctions uniformément continues et non uniformément continues.
5.  Valider la logique mathématique de votre implémentation à l'aide d'assertions.



\subsubsection{Continuité d'une fonction $f$ en un point $x_0$}

Une fonction $f: I \to \mathbb{R}$ est continue en $x_0 \in I$ si pour tout $\epsilon > 0$, il existe un $\delta > 0$ tel que pour tout $x \in I$, si $|x - x_0| < \delta$, alors $|f(x) - f(x_0)| < \epsilon$.
Notez que $\delta$ peut dépendre de $\epsilon$ et de $x_0$.

\subsubsection{Continuité Uniforme d'une fonction $f$ sur un intervalle $I$}

Une fonction $f: I \to \mathbb{R}$ est uniformément continue sur $I$ si pour tout $\epsilon > 0$, il existe un $\delta > 0$ tel que pour tous $x, y \in I$, si $|x - y| < \delta$, alors $|f(x) - f(y)| < \epsilon$.
La différence cruciale est que $\delta$ ne dépend \textit{que} de $\epsilon$, et non des points spécifiques $x, y$.

\subsubsection{Module de Continuité $\omega_f(\delta)$}

Le module de continuité d'une fonction $f$ sur un intervalle $I$ est défini pour $\delta > 0$ par :
$$ \omega_f(\delta) = \sup \{ |f(x) - f(y)| : x, y \in I, |x - y| < \delta \} $$
Une fonction $f$ est uniformément continue sur $I$ si et seulement si $\lim_{\delta \to 0^+} \omega_f(\delta) = 0$.

\subsection{Partie 1 : Implémentation du Module de Continuité Empirique}

L'objectif de cette partie est d'implémenter une fonction Python qui estime empiriquement le module de continuité $\omega_f(\delta)$ pour une fonction $f$ donnée sur un intervalle $[a, b]$ et pour une valeur $\delta$ spécifique. Puisque le calcul du \texttt{sup} (supremum) est difficile numériquement de manière exacte, nous allons l'approximer en échantillonnant un grand nombre de paires de points $(x, y)$ et en prenant le maximum des différences $|f(x) - f(y)|$.

\subsubsection{Consignes :}
1.  Créez une fonction \texttt{estimate\_modulus\_of\_continuity(f, a, b, delta, num\_pairs=10000)}.
2.  La fonction \texttt{f} sera une fonction Python prenant un argument numérique.
3.  \texttt{a} et \texttt{b} définissent l'intervalle $[a, b]$.
4.  \texttt{delta} est la valeur pour laquelle nous estimons $\omega_f(\delta)$.
5.  \texttt{num\_pairs} est le nombre de paires $(x, y)$ à échantillonner.
6.  Pour chaque échantillon :
    *   Choisissez un point \texttt{x} aléatoirement et uniformément dans $[a, b]$.
    *   Choisissez un point \texttt{y} tel que \texttt{|x - y| < delta} et \texttt{y} soit également dans $[a, b]$. Une stratégie efficace est de générer un décalage \texttt{d} aléatoirement dans l'intervalle \texttt{(-delta, delta)}, puis de définir \texttt{y = x + d}. Ensuite, assurez-vous que \texttt{y} reste dans l'intervalle \texttt{[a, b]} en le "clampant" (par exemple, \texttt{y = max(a, min(b, y))}). Cette approche garantit que \texttt{|x - y| < delta} est toujours respecté.
    *   Calculez \texttt{abs(f(x) - f(y))}.
7.  La fonction doit retourner la valeur maximale de \texttt{abs(f(x) - f(y))} trouvée parmi toutes les paires échantillonnées.
8.  Ajoutez des validations pour les arguments (par exemple, \texttt{delta > 0}, \texttt{a <= b}, \texttt{num\_pairs > 0}).

\subsection{Partie 2 : Vérification Empirique de la Continuité Uniforme}

Dans cette partie, vous utiliserez la fonction implémentée précédemment pour analyser la tendance du module de continuité lorsque $\delta$ tend vers zéro.

\subsubsection{Consignes :}
1.  Créez une fonction \texttt{check\_uniform\_continuity\_trend(f, a, b, delta\_values, num\_pairs=10000)}.
2.  \texttt{delta\_values} sera une liste de valeurs de $\delta$, idéalement décroissantes (par exemple, \texttt{[0.5, 0.2, 0.1, 0.05, 0.01, 0.005, 0.001]}).
3.  Pour chaque \texttt{delta} dans \texttt{delta\_values}, appelez \texttt{estimate\_modulus\_of\_continuity} et stockez le résultat.
4.  La fonction doit retourner une liste de tuples \texttt{(delta, omega\_f(delta)\_estimated)}.
5.  Cette liste permettra d'observer si \texttt{omega\_f(delta)} tend vers zéro à mesure que \texttt{delta} diminue.

\subsection{Partie 3 : Exemples et Assertions}

Vous allez maintenant tester votre implémentation avec des fonctions connues pour être uniformément continues et d'autres qui ne le sont pas.

\subsubsection{Consignes :}
1.  Définissez les fonctions de test suivantes :
    *   \texttt{f\_x\_squared(x)} : $f(x) = x\textasciicircum 2$
    *   \texttt{f\_sin\_x(x)} : $f(x) = \sin(x)$
    *   \texttt{f\_one\_over\_x(x)} : $f(x) = 1/x$ (attention à l'intervalle de définition)
    *   \texttt{f\_sqrt\_x(x)} : $f(x) = \sqrt{x}$ (attention à l'intervalle de définition)
2.  \textbf{Scénario 1 : Fonction uniformément continue.}
    *   Testez $f(x) = x\textasciicircum 2$ sur l'intervalle $[0, 1]$.
    *   Utilisez \texttt{check\_uniform\_continuity\_trend} avec une liste de \texttt{delta\_values} décroissantes.
    *   Affichez les résultats.
    *   Ajoutez des assertions pour vérifier que :
        *   Le module de continuité estimé est généralement décroissant à mesure que $\delta$ diminue.
        *   La valeur du module de continuité pour le plus petit $\delta$ est très proche de zéro (utilisez une petite tolérance).
3.  \textbf{Scénario 2 : Fonction non uniformément continue.}
    *   Testez $f(x) = 1/x$ sur l'intervalle $(0, 1]$ (par exemple, \texttt{[0.0001, 1.0]} pour éviter la singularité en 0).
    *   Utilisez \texttt{check\_uniform\_continuity\_trend} avec une liste de \texttt{delta\_values} décroissantes.
    *   Affichez les résultats.
    *   Ajoutez des assertions pour vérifier que :
        \textit{   La valeur du module de continuité pour le plus petit $\delta$ n'est }pas* proche de zéro, mais reste significativement élevée.
4.  \textbf{Scénario 3 : Autre fonction uniformément continue.}
    *   Testez $f(x) = \sqrt{x}$ sur l'intervalle $[0, 1]$.
    *   Répétez les étapes d'affichage et d'assertions similaires au Scénario 1.

---

\texttt{}`python
import random
import math

\section{--- Partie 1 : Implémentation du Module de Continuité Empirique ---}

def estimate\_modulus\_of\_continuity(f, a, b, delta, num\_pairs=10000):
    """
    Estime empiriquement le module de continuité omega\_f(delta) pour une fonction f
    sur l'intervalle [a, b].

    Args:
        f (callable): La fonction f(x) à évaluer.
        a (float): Borne inférieure de l'intervalle.
        b (float): Borne supérieure de l'intervalle.
        delta (float): La valeur delta pour laquelle estimer omega\_f(delta).
        num\_pairs (int): Le nombre de paires (x, y) à échantillonner pour l'estimation.

    Returns:
        float: L'estimation empirique de omega\_f(delta).
    """
    if not (a <= b):
        raise ValueError("La borne inférieure 'a' doit être inférieure ou égale à la borne supérieure 'b'.")
    if not (delta > 0):
        raise ValueError("Delta doit être strictement positif.")
    if not (num\_pairs > 0):
        raise ValueError("Le nombre de paires d'échantillons doit être strictement positif.")

    max\_diff = 0.0

    for \_ in range(num\_pairs):
        \# Échantillonner x uniformément dans [a, b]
        x = random.uniform(a, b)

        \# Générer un décalage 'd' dans (-delta, delta)
        d = random.uniform(-delta, delta)
        y = x + d

        \# S'assurer que y reste dans l'intervalle [a, b]
        \# Cette opération garantit que y est dans [a, b] et que |x - y| < delta.
        \# En effet, si y est clampé, la distance |x - y\_clamped| sera inférieure ou égale à |x - y|,
        \# et donc toujours inférieure à delta.
        y = max(a, min(b, y))

        \# Calculer la différence et mettre à jour le maximum
        diff = abs(f(x) - f(y))
        if diff > max\_diff:
            max\_diff = diff

    return max\_diff

\section{--- Partie 2 : Vérification Empirique de la Continuité Uniforme ---}

def check\_uniform\_continuity\_trend(f, a, b, delta\_values, num\_pairs=10000):
    """
    Vérifie la tendance du module de continuité pour évaluer empiriquement
    la continuité uniforme.

    Args:
        f (callable): La fonction f(x) à évaluer.
        a (float): Borne inférieure de l'intervalle.
        b (float): Borne supérieure de l'intervalle.
        delta\_values (list of float): Une liste de valeurs delta, idéalement décroissantes,
                                      pour lesquelles le module de continuité sera estimé.
        num\_pairs (int): Le nombre de paires (x, y) à échantillonner pour chaque estimation.

    Returns:
        list of tuple: Une liste de (delta, omega\_f(delta)\_estimated) pour chaque delta.
    """
    if not isinstance(delta\_values, list) or not all(isinstance(d, (int, float)) and d > 0 for d in delta\_values):
        raise ValueError("delta\_values doit être une liste de nombres flottants positifs.")

    results = []
    \# Itérer sur les valeurs delta fournies, en supposant qu'elles sont ordonnées
    \# pour l'analyse de tendance (généralement décroissantes).
    for delta in delta\_values:
        omega\_delta = estimate\_modulus\_of\_continuity(f, a, b, delta, num\_pairs)
        results.append((delta, omega\_delta))
    return results

\section{--- Partie 3 : Exemples et Assertions ---}

\section{Fonctions de test}
def f\_x\_squared(x):
    """Fonction f(x) = x\textasciicircum 2"""
    return x * x

def f\_sin\_x(x):
    """Fonction f(x) = sin(x)"""
    return math.sin(x)

def f\_one\_over\_x(x):
    """Fonction f(x) = 1/x. Non définie en 0."""
    if x == 0:
        raise ValueError("f\_one\_over\_x n'est pas définie pour x=0.")
    return 1.0 / x

def f\_sqrt\_x(x):
    """Fonction f(x) = sqrt(x). Non définie pour x<0."""
    if x < 0:
        raise ValueError("f\_sqrt\_x n'est pas définie pour x<0.")
    return math.sqrt(x)

\section{--- Scénario 1 : Fonction uniformément continue (f(x) = x\textasciicircum 2 sur [0, 1]) ---}
print("--- Scénario 1 : f(x) = x\textasciicircum 2 sur [0, 1] (Uniformément Continue) ---")
interval\_uc = (0.0, 1.0)
delta\_test\_values\_uc = [0.5, 0.2, 0.1, 0.05, 0.01, 0.005, 0.001]
results\_uc = check\_uniform\_continuity\_trend(f\_x\_squared, interval\_uc[0], interval\_uc[1], delta\_test\_values\_uc)

print("Résultats pour f(x) = x\textasciicircum 2 sur [0, 1]:")
for delta, omega in results\_uc:
    print(f"  delta = {delta:.4f}, omega\_f(delta) $\approx$ {omega:.6f}")

\section{Assertions pour la continuité uniforme:}
\section{Le module de continuité doit tendre vers 0 à mesure que delta tend vers 0.}
\section{Nous vérifions que chaque omega\_f(delta) est inférieur ou égal au précédent (tendance non croissante)}
\section{et que la dernière valeur est très proche de 0.}
assert len(results\_uc) > 1, "Pas assez de points pour vérifier la tendance."
for i in range(1, len(results\_uc)):
    \# Permettre de légères fluctuations dues à l'échantillonnage, mais la tendance générale doit être non croissante.
    assert results\_uc[i][1] <= results\_uc[i-1][1] + 1e-3, \
        f"Tendance non décroissante pour f(x)=x\textasciicircum 2: omega({results\_uc[i][0]})={results\_uc[i][1]:.6f} > omega({results\_uc[i-1][0]})={results\_uc[i-1][1]:.6f}"

\section{La dernière valeur (pour le plus petit delta) doit être très proche de 0.}
assert results\_uc[-1][1] < 0.005, \
    f"omega\_f(delta) pour le plus petit delta ({results\_uc[-1][0]:.4f}) n'est pas proche de 0: {results\_uc[-1][1]:.6f}"
print("Assertions pour f(x) = x\textasciicircum 2 passées avec succès : Tendance non croissante et limite proche de 0.")

\section{--- Scénario 2 : Fonction non uniformément continue (f(x) = 1/x sur (0, 1]) ---}
print("--- Scénario 2 : f(x) = 1/x sur (0, 1] (NON Uniformément Continue) ---")
\section{On doit éviter x=0, donc l'intervalle commence légèrement au-dessus de 0.}
\section{L'intervalle [0.0001, 1.0] est choisi pour illustrer le comportement près de la singularité.}
interval\_nuc = (0.0001, 1.0)
delta\_test\_values\_nuc = [0.5, 0.2, 0.1, 0.05, 0.01, 0.005, 0.001]
results\_nuc = check\_uniform\_continuity\_trend(f\_one\_over\_x, interval\_nuc[0], interval\_nuc[1], delta\_test\_values\_nuc)

print("Résultats pour f(x) = 1/x sur [0.0001, 1.0]:")
for delta, omega in results\_nuc:
    print(f"  delta = {delta:.4f}, omega\_f(delta) $\approx$ {omega:.6f}")

\section{Assertions pour la non-continuité uniforme:}
\section{Le module de continuité ne doit PAS tendre vers 0, ou doit rester significativement grand.}
\section{Pour f(x)=1/x sur (0,1], omega\_f(delta) ne tend pas vers 0.}
\section{L'estimation ici est limitée par la borne inférieure de l'intervalle (0.0001).}
\section{Pour delta=0.001, on peut trouver x=0.0001, y=0.0002 (si 0.0002 est dans [x-delta, x+delta] et [a,b]).}
\section{Alors |f(x)-f(y)| = |1/0.0001 - 1/0.0002| = |10000 - 5000| = 5000.}
\section{Donc, omega\_f(delta) devrait rester grand.}
assert results\_nuc[-1][1] > 100.0, \
    f"omega\_f(delta) pour le plus petit delta ({results\_nuc[-1][0]:.4f}) est trop proche de 0: {results\_nuc[-1][1]:.6f}"
print("Assertions pour f(x) = 1/x passées avec succès : Le module de continuité reste élevé.")

\section{--- Scénario 3 : Fonction uniformément continue (f(x) = sqrt(x) sur [0, 1]) ---}
print("--- Scénario 3 : f(x) = sqrt(x) sur [0, 1] (Uniformément Continue) ---")
interval\_sqrt = (0.0, 1.0)
delta\_test\_values\_sqrt = [0.5, 0.2, 0.1, 0.05, 0.01, 0.005, 0.001]
results\_sqrt = check\_uniform\_continuity\_trend(f\_sqrt\_x, interval\_sqrt[0], interval\_sqrt[1], delta\_test\_values\_sqrt)

print("Résultats pour f(x) = sqrt(x) sur [0, 1]:")
for delta, omega in results\_sqrt:
    print(f"  delta = {delta:.4f}, omega\_f(delta) $\approx$ {omega:.6f}")

\section{Assertions pour la continuité uniforme:}
assert len(results\_sqrt) > 1, "Pas assez de points pour vérifier la tendance."
for i in range(1, len(results\_sqrt)):
    assert results\_sqrt[i][1] <= results\_sqrt[i-1][1] + 1e-3, \
        f"Tendance non décroissante pour f(x)=sqrt(x): omega({results\_sqrt[i][0]})={results\_sqrt[i][1]:.6f} > omega({results\_sqrt[i-1][0]})={results\_sqrt[i-1][1]:.6f}"

assert results\_sqrt[-1][1] < 0.005, \
    f"omega\_f(delta) pour le plus petit delta ({results\_sqrt[-1][0]:.4f}) n'est pas proche de 0: {results\_sqrt[-1][1]:.6f}"
print("Assertions pour f(x) = sqrt(x) passées avec succès : Tendance non croissante et limite proche de 0.")


\end{document}
